% Supplement of the v1 paper (input by main.tex). Source keys as in main.tex.
% Placement (Lam, PR HEAD 19 Nov "after the draft review" (2)): P0, P1, P2, P5 here; P3, P6, P7, P9, G1
% already placed here (PR l.249-252, l.379-381, l.409-410). P4, P8 in the main text (Section IV-F).
% P10 (21 Nov): tables here, a summary paragraph in Section IV-F (agent default, coordinator brief; Lam decides).
% Phase-1 E1-E4: section S-D at the end, typeset only when the generated tables exist.
\clearpage
\onecolumn
\setcounter{table}{0}
\renewcommand{\thetable}{S\arabic{table}}
\renewcommand{\theHtable}{S\arabic{table}}
\setcounter{figure}{0}
\renewcommand{\thefigure}{S\arabic{figure}}
\renewcommand{\theHfigure}{S\arabic{figure}}
\begin{center}
{\Large Supplementary material}
\end{center}

\section*{S-A. Locked results that A3 places in the supplement}
These are pre-specified results from the locked run (results archive \texttt{results-20261114T093033Z.zip}); A3 places them in the supplement only. The Kermany all-250 false-rare rate, which A3 lists in the main-text secondary row, is in Table~\ref{tab:frr}.
% src: AM l.184-190; PR l.114-118

\subsection*{S1. Logistic probe (secondary classifier)}
Table~\ref{tab:logistic} gives five-way macro recall with the fixed L2-regularised logistic probe on the hold-out, first 25 support draws per $K$. In development the probe ran at $K=4$ only (125 episodes per arm); the locked run wrote no development summary for it, and its only summary is the post hoc P6 (Table~\ref{tab:p6}).
% src: ARC support_draw_summary.csv (classifier=logistic), paired_differences.csv (logistic); ARC runs/<arm>/development/logistic/ (per-episode files only); methods_results_v2.md section 8.2 S1

\begin{table}[!htbp]
\caption{Hold-out five-way macro recall with the fixed logistic probe (secondary classifier), first 25 support draws per $K$. Paired rows: difference from Inception-v3 on the same draws.}
\label{tab:logistic}
\centering\footnotesize
% src: ARC reports/legacy/support_draw_summary.csv classifier=logistic; reports/legacy/paired_differences.csv classifier=logistic.
\begin{tabular}{@{}lcccccc@{}}
\toprule
Arm & $K=1$ & $K=2$ & $K=4$ & $K=8$ & $K=16$ & \makecell{Mean $K$ 1--8\\(SD; 5th--95th)} \\
\midrule
MIRAGE & 0.564 & 0.670 & 0.801 & 0.873 & 0.913 & 0.727 (0.068; 0.600--0.815) \\
DINOv2 & 0.417 & 0.522 & 0.630 & 0.716 & 0.744 & 0.571 (0.059; 0.465--0.670) \\
Inception-v3 & 0.416 & 0.498 & 0.586 & 0.662 & 0.687 & 0.541 (0.064; 0.466--0.652) \\
RETFound & 0.258 & 0.312 & 0.372 & 0.424 & 0.501 & 0.341 (0.063; 0.235--0.426) \\
\midrule
MIRAGE $-$ Inception-v3 & +0.148 & +0.172 & +0.216 & +0.210 & +0.226 & +0.187 (0.080; +0.073 to +0.309) \\
DINOv2 $-$ Inception-v3 & +0.001 & +0.024 & +0.044 & +0.054 & +0.056 & +0.031 (0.076; $-$0.096 to +0.126) \\
RETFound $-$ Inception-v3 & $-$0.158 & $-$0.186 & $-$0.213 & $-$0.239 & $-$0.186 & $-$0.199 (0.075; $-$0.315 to $-$0.097) \\
\bottomrule
\end{tabular}

\end{table}

\subsection*{S2. Five-way output shares on external disease images}
\begin{table}[!htbp]
\caption{Five-way head on external images labelled CSC/CSR or MH by OCTID or OCTDL: share of queries given their own class, mean over draws, by $K$. Supports are benchmark images only. These measure agreement with the datasets' published labels (OCTID CSC and MH: common classes in OCTID, tier D; OCTDL MH: tier B). MacTel, Stargardt and RP receive no external claim. Mean counts per predicted class are in \texttt{external\_transfer\_outputs.csv} (released).}
\label{tab:transfer}
\centering\footnotesize
% src: ARC reports/external/external_transfer_summary.csv columns mean, sd, p05, p95 (share of queries given their
%      own class; over the available draws: prototype 100, logistic probe 25).
\begin{tabular}{@{}lll*{5}{c}@{}}
\toprule
Images & Arm & Head & $K=1$ & $K=2$ & $K=4$ & $K=8$ & $K=16$ \\
\midrule
OCTID CSC (102) & MIRAGE & prototype & \makecell{0.441\\{\scriptsize (0.231; 0.068--0.845)}} & \makecell{0.460\\{\scriptsize (0.235; 0.107--0.883)}} & \makecell{0.495\\{\scriptsize (0.193; 0.175--0.795)}} & \makecell{0.573\\{\scriptsize (0.144; 0.343--0.794)}} & \makecell{0.611\\{\scriptsize (0.079; 0.480--0.735)}} \\
 &  & logistic & \makecell{0.346\\{\scriptsize (0.208; 0.053--0.676)}} & \makecell{0.378\\{\scriptsize (0.261; 0.029--0.782)}} & \makecell{0.425\\{\scriptsize (0.175; 0.133--0.659)}} & \makecell{0.505\\{\scriptsize (0.104; 0.373--0.625)}} & \makecell{0.504\\{\scriptsize (0.075; 0.375--0.606)}} \\
 & DINOv2 & prototype & \makecell{0.165\\{\scriptsize (0.195; 0.000--0.551)}} & \makecell{0.162\\{\scriptsize (0.138; 0.010--0.471)}} & \makecell{0.199\\{\scriptsize (0.136; 0.039--0.461)}} & \makecell{0.246\\{\scriptsize (0.101; 0.098--0.403)}} & \makecell{0.290\\{\scriptsize (0.069; 0.176--0.402)}} \\
 &  & logistic & \makecell{0.238\\{\scriptsize (0.239; 0.000--0.655)}} & \makecell{0.175\\{\scriptsize (0.131; 0.041--0.369)}} & \makecell{0.258\\{\scriptsize (0.157; 0.094--0.573)}} & \makecell{0.328\\{\scriptsize (0.100; 0.180--0.494)}} & \makecell{0.384\\{\scriptsize (0.055; 0.278--0.451)}} \\
 & Inception-v3 & prototype & \makecell{0.268\\{\scriptsize (0.208; 0.048--0.677)}} & \makecell{0.293\\{\scriptsize (0.167; 0.078--0.599)}} & \makecell{0.339\\{\scriptsize (0.144; 0.157--0.608)}} & \makecell{0.407\\{\scriptsize (0.114; 0.255--0.609)}} & \makecell{0.452\\{\scriptsize (0.068; 0.333--0.550)}} \\
 &  & logistic & \makecell{0.289\\{\scriptsize (0.195; 0.055--0.675)}} & \makecell{0.300\\{\scriptsize (0.126; 0.139--0.537)}} & \makecell{0.345\\{\scriptsize (0.126; 0.182--0.557)}} & \makecell{0.400\\{\scriptsize (0.099; 0.288--0.573)}} & \makecell{0.416\\{\scriptsize (0.063; 0.337--0.518)}} \\
 & RETFound & prototype & \makecell{0.115\\{\scriptsize (0.275; 0.000--0.952)}} & \makecell{0.094\\{\scriptsize (0.246; 0.000--0.853)}} & \makecell{0.048\\{\scriptsize (0.161; 0.000--0.272)}} & \makecell{0.020\\{\scriptsize (0.072; 0.000--0.148)}} & \makecell{0.004\\{\scriptsize (0.013; 0.000--0.040)}} \\
 &  & logistic & \makecell{0.120\\{\scriptsize (0.285; 0.000--0.892)}} & \makecell{0.033\\{\scriptsize (0.099; 0.000--0.298)}} & \makecell{0.067\\{\scriptsize (0.222; 0.000--0.549)}} & \makecell{0.024\\{\scriptsize (0.087; 0.000--0.106)}} & \makecell{0.003\\{\scriptsize (0.014; 0.000--0.000)}} \\
\addlinespace[2pt]
OCTID MH (102) & MIRAGE & prototype & \makecell{0.392\\{\scriptsize (0.292; 0.000--0.883)}} & \makecell{0.451\\{\scriptsize (0.300; 0.000--0.912)}} & \makecell{0.572\\{\scriptsize (0.256; 0.097--0.912)}} & \makecell{0.702\\{\scriptsize (0.184; 0.362--0.942)}} & \makecell{0.774\\{\scriptsize (0.096; 0.607--0.892)}} \\
 &  & logistic & \makecell{0.465\\{\scriptsize (0.308; 0.000--0.963)}} & \makecell{0.524\\{\scriptsize (0.318; 0.006--0.941)}} & \makecell{0.570\\{\scriptsize (0.274; 0.120--0.910)}} & \makecell{0.689\\{\scriptsize (0.161; 0.467--0.900)}} & \makecell{0.711\\{\scriptsize (0.094; 0.524--0.812)}} \\
 & DINOv2 & prototype & \makecell{0.242\\{\scriptsize (0.197; 0.000--0.578)}} & \makecell{0.274\\{\scriptsize (0.182; 0.020--0.549)}} & \makecell{0.319\\{\scriptsize (0.177; 0.049--0.559)}} & \makecell{0.367\\{\scriptsize (0.132; 0.127--0.550)}} & \makecell{0.421\\{\scriptsize (0.079; 0.275--0.549)}} \\
 &  & logistic & \makecell{0.304\\{\scriptsize (0.227; 0.012--0.622)}} & \makecell{0.347\\{\scriptsize (0.198; 0.033--0.563)}} & \makecell{0.415\\{\scriptsize (0.157; 0.139--0.608)}} & \makecell{0.507\\{\scriptsize (0.087; 0.380--0.606)}} & \makecell{0.561\\{\scriptsize (0.048; 0.476--0.618)}} \\
 & Inception-v3 & prototype & \makecell{0.314\\{\scriptsize (0.219; 0.039--0.717)}} & \makecell{0.363\\{\scriptsize (0.176; 0.087--0.657)}} & \makecell{0.422\\{\scriptsize (0.142; 0.186--0.627)}} & \makecell{0.455\\{\scriptsize (0.110; 0.302--0.647)}} & \makecell{0.471\\{\scriptsize (0.056; 0.382--0.569)}} \\
 &  & logistic & \makecell{0.404\\{\scriptsize (0.212; 0.082--0.763)}} & \makecell{0.446\\{\scriptsize (0.153; 0.222--0.688)}} & \makecell{0.500\\{\scriptsize (0.135; 0.241--0.675)}} & \makecell{0.541\\{\scriptsize (0.089; 0.431--0.694)}} & \makecell{0.563\\{\scriptsize (0.047; 0.504--0.643)}} \\
 & RETFound & prototype & \makecell{0.137\\{\scriptsize (0.297; 0.000--0.972)}} & \makecell{0.141\\{\scriptsize (0.313; 0.000--0.991)}} & \makecell{0.150\\{\scriptsize (0.307; 0.000--1.000)}} & \makecell{0.145\\{\scriptsize (0.314; 0.000--0.980)}} & \makecell{0.008\\{\scriptsize (0.033; 0.000--0.041)}} \\
 &  & logistic & \makecell{0.159\\{\scriptsize (0.317; 0.000--0.865)}} & \makecell{0.280\\{\scriptsize (0.420; 0.000--1.000)}} & \makecell{0.255\\{\scriptsize (0.407; 0.000--1.000)}} & \makecell{0.138\\{\scriptsize (0.303; 0.000--0.927)}} & \makecell{0.034\\{\scriptsize (0.082; 0.000--0.165)}} \\
\addlinespace[2pt]
OCTDL MH (39) & MIRAGE & prototype & \makecell{0.195\\{\scriptsize (0.248; 0.000--0.771)}} & \makecell{0.228\\{\scriptsize (0.288; 0.000--0.923)}} & \makecell{0.223\\{\scriptsize (0.271; 0.000--0.897)}} & \makecell{0.252\\{\scriptsize (0.255; 0.026--0.774)}} & \makecell{0.256\\{\scriptsize (0.164; 0.077--0.590)}} \\
 &  & logistic & \makecell{0.281\\{\scriptsize (0.332; 0.000--0.933)}} & \makecell{0.281\\{\scriptsize (0.338; 0.000--0.944)}} & \makecell{0.278\\{\scriptsize (0.305; 0.000--0.928)}} & \makecell{0.237\\{\scriptsize (0.245; 0.010--0.754)}} & \makecell{0.222\\{\scriptsize (0.149; 0.056--0.508)}} \\
 & DINOv2 & prototype & \makecell{0.379\\{\scriptsize (0.251; 0.000--0.769)}} & \makecell{0.453\\{\scriptsize (0.226; 0.101--0.822)}} & \makecell{0.543\\{\scriptsize (0.189; 0.205--0.846)}} & \makecell{0.671\\{\scriptsize (0.116; 0.486--0.847)}} & \makecell{0.695\\{\scriptsize (0.063; 0.615--0.821)}} \\
 &  & logistic & \makecell{0.479\\{\scriptsize (0.293; 0.051--0.867)}} & \makecell{0.552\\{\scriptsize (0.221; 0.149--0.913)}} & \makecell{0.714\\{\scriptsize (0.153; 0.415--0.892)}} & \makecell{0.819\\{\scriptsize (0.104; 0.621--0.923)}} & \makecell{0.890\\{\scriptsize (0.071; 0.749--0.949)}} \\
 & Inception-v3 & prototype & \makecell{0.239\\{\scriptsize (0.273; 0.000--0.797)}} & \makecell{0.260\\{\scriptsize (0.241; 0.000--0.721)}} & \makecell{0.302\\{\scriptsize (0.275; 0.000--0.773)}} & \makecell{0.378\\{\scriptsize (0.253; 0.026--0.822)}} & \makecell{0.390\\{\scriptsize (0.182; 0.127--0.694)}} \\
 &  & logistic & \makecell{0.333\\{\scriptsize (0.332; 0.000--0.862)}} & \makecell{0.314\\{\scriptsize (0.292; 0.026--0.836)}} & \makecell{0.359\\{\scriptsize (0.281; 0.031--0.872)}} & \makecell{0.383\\{\scriptsize (0.233; 0.051--0.764)}} & \makecell{0.430\\{\scriptsize (0.209; 0.164--0.718)}} \\
 & RETFound & prototype & \makecell{0.178\\{\scriptsize (0.307; 0.000--0.874)}} & \makecell{0.179\\{\scriptsize (0.329; 0.000--1.000)}} & \makecell{0.212\\{\scriptsize (0.352; 0.000--1.000)}} & \makecell{0.176\\{\scriptsize (0.338; 0.000--0.976)}} & \makecell{0.027\\{\scriptsize (0.095; 0.000--0.208)}} \\
 &  & logistic & \makecell{0.255\\{\scriptsize (0.354; 0.000--0.821)}} & \makecell{0.320\\{\scriptsize (0.434; 0.000--1.000)}} & \makecell{0.337\\{\scriptsize (0.429; 0.000--1.000)}} & \makecell{0.195\\{\scriptsize (0.331; 0.000--0.918)}} & \makecell{0.059\\{\scriptsize (0.156; 0.000--0.256)}} \\
\bottomrule
\end{tabular}

\end{table}
% src: ARC reports/external/external_transfer_summary.csv; AM l.32-40 (A1 rules 1-3)

\subsection*{S3. Nine-way per-class recall with 95\% intervals}
\begin{table}[!htbp]
\caption{Nine-way per-class recall: count/denominator (rate; 95\% interval). Interval methods as in Table~\ref{tab:frr}. The Kermany rows include the 22 hold-out rare images.}
\label{tab:perclass}
\centering\scriptsize
\setlength{\tabcolsep}{3pt}
% src: ARC reports/external/external_source_audit.csv endpoints recall_<class>: numerator/denominator (value; interval_low--interval_high).
\begin{tabular}{@{}llcccc@{}}
\toprule
Source & Class & MIRAGE & DINOv2 & Inception-v3 & RETFound \\
\midrule
Kermany + hold-out & NORMAL & 250/250 (1.000; 1.000--1.000) & 218/250 (0.872; 0.828--0.914) & 223/250 (0.892; 0.851--0.929) & 215/250 (0.860; 0.816--0.901) \\
Kermany + hold-out & CNV & 229/250 (0.916; 0.868--0.957) & 209/250 (0.836; 0.771--0.890) & 170/250 (0.680; 0.588--0.765) & 147/250 (0.588; 0.515--0.664) \\
Kermany + hold-out & DME & 200/250 (0.800; 0.731--0.859) & 152/250 (0.608; 0.541--0.676) & 138/250 (0.552; 0.485--0.615) & 99/250 (0.396; 0.332--0.459) \\
Kermany + hold-out & DRUSEN & 230/250 (0.920; 0.873--0.962) & 162/250 (0.648; 0.576--0.720) & 163/250 (0.652; 0.579--0.722) & 104/250 (0.416; 0.352--0.481) \\
Kermany + hold-out & CSC & 5/5 (1.000; 1.000--1.000) & 4/5 (0.800; 0.400--1.000) & 4/5 (0.800; 0.400--1.000) & 1/5 (0.200; 0.000--0.600) \\
Kermany + hold-out & MH & 5/5 (1.000; 1.000--1.000) & 4/5 (0.800; 0.400--1.000) & 4/5 (0.800; 0.400--1.000) & 1/5 (0.200; 0.000--0.600) \\
Kermany + hold-out & MacTel & 2/4 (0.500; 0.000--1.000) & 4/4 (1.000; 1.000--1.000) & 1/4 (0.250; 0.000--0.750) & 2/4 (0.500; 0.000--1.000) \\
Kermany + hold-out & Stargardt\textsuperscript{a} & 4/4 (1.000; 1.000--1.000) & 2/4 (0.500; 0.000--1.000) & 3/4 (0.750; 0.250--1.000) & 1/4 (0.250; 0.000--0.750) \\
Kermany + hold-out & RP\textsuperscript{a} & 4/4 (1.000; 1.000--1.000) & 1/4 (0.250; 0.000--0.750) & 1/4 (0.250; 0.000--0.750) & 4/4 (1.000; 1.000--1.000) \\
OCTID & NORMAL & 202/206 (0.981; 0.951--0.995) & 188/206 (0.913; 0.865--0.947) & 182/206 (0.883; 0.832--0.924) & 189/206 (0.917; 0.871--0.951) \\
OCTID & CSC & 23/102 (0.225; 0.149--0.319) & 30/102 (0.294; 0.208--0.393) & 22/102 (0.216; 0.140--0.308) & 0/102 (0.000; 0.000--0.036) \\
OCTID & MH & 10/102 (0.098; 0.048--0.173) & 48/102 (0.471; 0.371--0.572) & 27/102 (0.265; 0.182--0.361) & 0/102 (0.000; 0.000--0.036) \\
OCTDL & NORMAL & 26/332 (0.078; 0.049--0.113) & 145/332 (0.437; 0.383--0.493) & 66/332 (0.199; 0.160--0.239) & 137/332 (0.413; 0.341--0.475) \\
OCTDL & MH & 3/39 (0.077; 0.000--0.206) & 26/39 (0.667; 0.472--0.830) & 0/39 (0.000; 0.000--0.000) & 0/39 (0.000; 0.000--0.000) \\
\bottomrule
\end{tabular}

\par\smallskip
\parbox{\textwidth}{\footnotesize \fnRS The folder labels for RP and Stargardt disagree with the counts in the original paper; macro recall is unaffected.}
\end{table}

\subsection*{S4. Nine-way confusion counts}
\begin{table}[!htbp]
\caption{Nine-way confusion counts (rows: true class; columns: assigned class, in the order NORMAL, CNV, DME, DRUSEN, CSC, MH, MacTel, Stargardt\fnRS, RP\fnRS). ``Kermany + hold-out'': the 1,000 Kermany test images and the 22 hold-out rare images. Empty true rows are dropped.}
\label{tab:confusion}
\centering\scriptsize
\setlength{\tabcolsep}{3.5pt}
% src: ARC runs/<arm>/reporting/nine_way/metrics.json and runs/<arm>/external/{OCTID,OCTDL}/nine_way/metrics.json,
%      keys classes and confusion (rows true, columns assigned); empty true rows dropped.
\begin{tabular}{@{}lll*{9}{r}@{}}
\toprule
Arm & Source & True & NOR & CNV & DME & DRU & CSC & MH & MacT & Starg & RP \\
\midrule
MIRAGE & Kermany + hold-out & NOR & 250 & 0 & 0 & 0 & 0 & 0 & 0 & 0 & 0 \\
 &  & CNV & 0 & 229 & 4 & 17 & 0 & 0 & 0 & 0 & 0 \\
 &  & DME & 26 & 12 & 200 & 6 & 0 & 0 & 0 & 0 & 6 \\
 &  & DRU & 6 & 11 & 3 & 230 & 0 & 0 & 0 & 0 & 0 \\
 &  & CSC & 0 & 0 & 0 & 0 & 5 & 0 & 0 & 0 & 0 \\
 &  & MH & 0 & 0 & 0 & 0 & 0 & 5 & 0 & 0 & 0 \\
 &  & MacT & 0 & 0 & 0 & 0 & 0 & 1 & 2 & 1 & 0 \\
 &  & Starg & 0 & 0 & 0 & 0 & 0 & 0 & 0 & 4 & 0 \\
 &  & RP & 0 & 0 & 0 & 0 & 0 & 0 & 0 & 0 & 4 \\
\addlinespace[1pt]
 & OCTID & NOR & 202 & 0 & 0 & 0 & 0 & 0 & 4 & 0 & 0 \\
 &  & CSC & 8 & 28 & 23 & 2 & 23 & 1 & 1 & 15 & 1 \\
 &  & MH & 8 & 10 & 72 & 1 & 0 & 10 & 0 & 1 & 0 \\
\addlinespace[1pt]
 & OCTDL & NOR & 26 & 148 & 118 & 10 & 0 & 6 & 24 & 0 & 0 \\
 &  & MH & 0 & 6 & 30 & 0 & 0 & 3 & 0 & 0 & 0 \\
\addlinespace[1pt]
\midrule
DINOv2 & Kermany + hold-out & NOR & 218 & 5 & 3 & 14 & 2 & 0 & 1 & 0 & 7 \\
 &  & CNV & 2 & 209 & 20 & 8 & 4 & 0 & 3 & 0 & 4 \\
 &  & DME & 24 & 22 & 152 & 14 & 4 & 0 & 14 & 0 & 20 \\
 &  & DRU & 14 & 65 & 5 & 162 & 0 & 0 & 0 & 0 & 4 \\
 &  & CSC & 0 & 0 & 0 & 0 & 4 & 0 & 1 & 0 & 0 \\
 &  & MH & 0 & 0 & 0 & 0 & 0 & 4 & 0 & 0 & 1 \\
 &  & MacT & 0 & 0 & 0 & 0 & 0 & 0 & 4 & 0 & 0 \\
 &  & Starg & 0 & 0 & 0 & 0 & 0 & 0 & 0 & 2 & 2 \\
 &  & RP & 0 & 1 & 0 & 0 & 0 & 0 & 0 & 2 & 1 \\
\addlinespace[1pt]
 & OCTID & NOR & 188 & 1 & 0 & 0 & 0 & 0 & 3 & 3 & 11 \\
 &  & CSC & 11 & 19 & 1 & 2 & 30 & 0 & 3 & 30 & 6 \\
 &  & MH & 4 & 2 & 6 & 5 & 2 & 48 & 23 & 2 & 10 \\
\addlinespace[1pt]
 & OCTDL & NOR & 145 & 29 & 1 & 16 & 6 & 0 & 10 & 58 & 67 \\
 &  & MH & 0 & 0 & 0 & 0 & 1 & 26 & 8 & 2 & 2 \\
\addlinespace[1pt]
\midrule
Inception-v3 & Kermany + hold-out & NOR & 223 & 0 & 9 & 11 & 0 & 0 & 5 & 1 & 1 \\
 &  & CNV & 7 & 170 & 30 & 22 & 3 & 1 & 0 & 13 & 4 \\
 &  & DME & 44 & 35 & 138 & 16 & 0 & 3 & 2 & 3 & 9 \\
 &  & DRU & 20 & 30 & 11 & 163 & 0 & 0 & 0 & 24 & 2 \\
 &  & CSC & 0 & 1 & 0 & 0 & 4 & 0 & 0 & 0 & 0 \\
 &  & MH & 0 & 0 & 0 & 0 & 1 & 4 & 0 & 0 & 0 \\
 &  & MacT & 0 & 0 & 0 & 0 & 0 & 0 & 1 & 1 & 2 \\
 &  & Starg & 0 & 0 & 0 & 0 & 0 & 0 & 1 & 3 & 0 \\
 &  & RP & 0 & 0 & 1 & 0 & 1 & 0 & 0 & 1 & 1 \\
\addlinespace[1pt]
 & OCTID & NOR & 182 & 0 & 9 & 0 & 12 & 0 & 3 & 0 & 0 \\
 &  & CSC & 3 & 16 & 26 & 11 & 22 & 0 & 2 & 14 & 8 \\
 &  & MH & 6 & 5 & 45 & 2 & 2 & 27 & 0 & 1 & 14 \\
\addlinespace[1pt]
 & OCTDL & NOR & 66 & 196 & 21 & 38 & 0 & 0 & 0 & 0 & 11 \\
 &  & MH & 2 & 25 & 11 & 0 & 1 & 0 & 0 & 0 & 0 \\
\addlinespace[1pt]
\midrule
RETFound & Kermany + hold-out & NOR & 215 & 1 & 2 & 31 & 0 & 0 & 1 & 0 & 0 \\
 &  & CNV & 3 & 147 & 83 & 16 & 0 & 0 & 0 & 0 & 1 \\
 &  & DME & 56 & 50 & 99 & 41 & 0 & 0 & 0 & 0 & 4 \\
 &  & DRU & 65 & 36 & 45 & 104 & 0 & 0 & 0 & 0 & 0 \\
 &  & CSC & 0 & 0 & 1 & 0 & 1 & 2 & 0 & 1 & 0 \\
 &  & MH & 0 & 0 & 1 & 0 & 1 & 1 & 1 & 0 & 1 \\
 &  & MacT & 0 & 1 & 0 & 0 & 1 & 0 & 2 & 0 & 0 \\
 &  & Starg & 0 & 0 & 0 & 0 & 0 & 0 & 0 & 1 & 3 \\
 &  & RP & 0 & 0 & 0 & 0 & 0 & 0 & 0 & 0 & 4 \\
\addlinespace[1pt]
 & OCTID & NOR & 189 & 0 & 3 & 14 & 0 & 0 & 0 & 0 & 0 \\
 &  & CSC & 13 & 12 & 72 & 5 & 0 & 0 & 0 & 0 & 0 \\
 &  & MH & 8 & 3 & 80 & 11 & 0 & 0 & 0 & 0 & 0 \\
\addlinespace[1pt]
 & OCTDL & NOR & 137 & 8 & 167 & 2 & 0 & 0 & 12 & 1 & 5 \\
 &  & MH & 0 & 0 & 36 & 0 & 0 & 0 & 0 & 0 & 3 \\
\addlinespace[1pt]
\bottomrule
\end{tabular}

\end{table}
% src: ARC runs/<arm>/{reporting,external/OCTID,external/OCTDL}/nine_way/metrics.json

\clearpage
\subsection*{S5. Pilot timing, repeat checks and external screen}
\begin{table}[!htbp]
\caption{Locked pilot on one NVIDIA L4 (code \texttt{2f48654}): median warm time per image over 100 images, loading time, peak GPU memory and the repeat check; one pilot sequence in the session.}
\label{tab:pilot}
\centering\footnotesize
\begin{tabular}{@{}lrrrcc@{}}
\toprule
Arm & ms per image & Load (s) & Peak GPU memory (GB) & Worst difference & Lowest cosine \\
\midrule
MIRAGE & 13.57 & 5.5 & 0.37 & 0 & 1.000000 \\
DINOv2 & 14.03 & 5.4 & 0.37 & 0 & 1.000000 \\
Inception-v3 & 323.66 & 9.5 & 2.33 & 0 & 1.000000 \\
RETFound & 23.91 & 22.0 & 1.24 & 0 & 1.000000 \\
\bottomrule
\end{tabular}
\par\smallskip
\parbox{\textwidth}{\footnotesize RETFound loaded with missing keys \texttt{fc\_norm.weight}, \texttt{fc\_norm.bias}, \texttt{head.weight}, \texttt{head.bias} and 106 unexpected keys, including \texttt{norm.weight} (Section~\ref{sec:models}). External screen: OCTID 410 queries from 410 candidates, OCTDL 371 from 371, none excluded; OCTID 410 and OCTDL 113 dependence groups; flagged pairs among external images 25 (difference hash) and 4 (identical pixels); none between an external image and a fitting or hold-out image.}
\end{table}
% src: DRV step11_pilots.txt lines 1-15; DRV step14_external.txt lines 1-3

\subsection*{S6. OCTID two-way test: per-class recall and confusion counts}
\begin{table}[!htbp]
\caption{OCTID two-way test per $K$: macro recall, per-class recall and mean confusion counts over draws (denominators 102 CSC and 102 MH), and draws in which every query went to one class.}
\label{tab:twoway}
\centering\footnotesize
% src: ARC reports/external/external_two_way.csv lines 2-31 (recall_CSC_mean, recall_MH_mean, mean confusion counts, collapsed_draws).
\begin{tabular}{@{}lrcccccccr@{}}
\toprule
Arm & $K$ & Macro & CSC recall & MH recall & CSC$\to$CSC & CSC$\to$MH & MH$\to$CSC & MH$\to$MH & Collapsed \\
\midrule
RETFound & 1 & 0.507 & 0.562 & 0.452 & 57.36 & 44.64 & 55.86 & 46.14 & 60 \\
 & 2 & 0.514 & 0.602 & 0.426 & 61.39 & 40.61 & 58.58 & 43.42 & 56 \\
 & 4 & 0.510 & 0.494 & 0.525 & 50.39 & 51.61 & 48.43 & 53.57 & 50 \\
 & 8 & 0.518 & 0.501 & 0.535 & 51.06 & 50.94 & 47.47 & 54.53 & 44 \\
 & 16 & 0.541 & 0.463 & 0.618 & 47.24 & 54.76 & 38.95 & 63.05 & 20 \\
\addlinespace[2pt]
DINOv2 & 1 & 0.669 & 0.685 & 0.654 & 69.87 & 32.13 & 35.34 & 66.66 & 0 \\
 & 2 & 0.755 & 0.816 & 0.695 & 83.22 & 18.78 & 31.10 & 70.90 & 0 \\
 & 4 & 0.805 & 0.876 & 0.734 & 89.32 & 12.68 & 27.17 & 74.83 & 0 \\
 & 8 & 0.844 & 0.924 & 0.765 & 94.26 & 7.74 & 23.99 & 78.01 & 0 \\
 & 16 & 0.863 & 0.957 & 0.769 & 97.60 & 4.40 & 23.61 & 78.39 & 0 \\
\addlinespace[2pt]
MIRAGE & 1 & 0.647 & 0.707 & 0.587 & 72.16 & 29.84 & 42.11 & 59.89 & 3 \\
 & 2 & 0.690 & 0.754 & 0.627 & 76.89 & 25.11 & 38.08 & 63.92 & 0 \\
 & 4 & 0.773 & 0.813 & 0.732 & 82.92 & 19.08 & 27.33 & 74.67 & 0 \\
 & 8 & 0.837 & 0.849 & 0.825 & 86.63 & 15.37 & 17.81 & 84.19 & 0 \\
 & 16 & 0.898 & 0.915 & 0.881 & 93.29 & 8.71 & 12.14 & 89.86 & 0 \\
\addlinespace[2pt]
Inception-v3 & 1 & 0.623 & 0.628 & 0.618 & 64.07 & 37.93 & 38.99 & 63.01 & 0 \\
 & 2 & 0.676 & 0.686 & 0.666 & 70.00 & 32.00 & 34.04 & 67.96 & 0 \\
 & 4 & 0.742 & 0.760 & 0.724 & 77.54 & 24.46 & 28.20 & 73.80 & 0 \\
 & 8 & 0.802 & 0.852 & 0.752 & 86.90 & 15.10 & 25.34 & 76.66 & 0 \\
 & 16 & 0.844 & 0.890 & 0.799 & 90.76 & 11.24 & 20.54 & 81.46 & 0 \\
\addlinespace[2pt]
Size-only & 1 & 0.500 & 1.000 & 0.000 & 102.00 & 0.00 & 102.00 & 0.00 & 100 \\
 & 2 & 0.500 & 1.000 & 0.000 & 102.00 & 0.00 & 102.00 & 0.00 & 100 \\
 & 4 & 0.500 & 1.000 & 0.000 & 102.00 & 0.00 & 102.00 & 0.00 & 100 \\
 & 8 & 0.500 & 1.000 & 0.000 & 102.00 & 0.00 & 102.00 & 0.00 & 100 \\
 & 16 & 0.500 & 1.000 & 0.000 & 102.00 & 0.00 & 102.00 & 0.00 & 100 \\
\bottomrule
\end{tabular}

\end{table}
% src: ARC reports/external/external_two_way.csv lines 2-31; AM l.37-40 (A1 rule 3)

\subsection*{S7. Image geometry by source, class and role (diagnostic)}
\begin{table}[!htbp]
\caption{Image geometry by source, class and role, from the original files: median (min--max). Border share: the share of image area in complete edge rows and columns whose grey range is at most 8. Near-white share: pixels at 250 or above (median). Retained fraction: the share of the image each input transform keeps (MIRAGE and Inception-v3 resize the whole image; DINOv2 and RETFound resize and centre-crop). Descriptive; it gates nothing.}
\label{tab:geometry}
\centering\scriptsize
\setlength{\tabcolsep}{3pt}
% src: ARC reports/geometry/MIRAGE/geometry_summary.csv lines 2-24 (all columns except retained fraction);
%      retained fraction: same file (MIRAGE = Inception-v3, 1.0 everywhere) and reports/geometry/DINOv2/geometry_summary.csv
%      (DINOv2 = RETFound, files identical by MD5 per the draft, methods_results_v2.md section 8, Table 6).
\begin{tabular}{@{}lllrcccccc@{}}
\toprule
Source & Class & Role & $n$ & Width, px & Height, px & Border share & Near-white share & \makecell{Retained:\\MIRAGE, Inc.} & \makecell{Retained:\\DINOv2, RETFound} \\
\midrule
benchmark & CSC & development & 27 & 414 (410--644) & 413 (345--417) & 0.502 (0.007--0.746) & 0.000 & 1.000 & 0.763 (0.496--0.766) \\
benchmark & CSC & hold-out & 5 & 415 (413--417) & 412 (410--416) & 0.411 (0.058--0.566) & 0.000 & 1.000 & 0.766 (0.760--0.766) \\
benchmark & MH & development & 26 & 327 (321--491) & 326 (286--387) & 0.445 (0.089--0.674) & 0.000 & 1.000 & 0.763 (0.605--0.766) \\
benchmark & MH & hold-out & 5 & 327 (324--330) & 326 (311--327) & 0.386 (0.219--0.630) & 0.000 & 1.000 & 0.760 (0.734--0.766) \\
benchmark & MacTel & development & 25 & 482 (366--505) & 435 (286--437) & 0.374 (0.215--0.600) & 0.000 & 1.000 & 0.693 (0.557--0.763) \\
benchmark & MacTel & hold-out & 4 & 482 (481--484) & 436 (434--436) & 0.561 (0.272--0.638) & 0.000 & 1.000 & 0.691 (0.690--0.695) \\
benchmark & Stargardt\textsuperscript{a} & development & 19 & 410 (332--580) & 387 (258--525) & 0.374 (0.000--0.643) & 0.000 & 1.000 & 0.754 (0.543--0.766) \\
benchmark & Stargardt\textsuperscript{a} & hold-out & 4 & 374.5 (374--376) & 374.5 (373--375) & 0.369 (0.261--0.411) & 0.000 & 1.000 & 0.766 (0.763--0.766) \\
benchmark & RP\textsuperscript{a} & development & 27 & 393 (292--497) & 390 (269--491) & 0.339 (0.000--0.563) & 0.000 & 1.000 & 0.757 (0.628--0.766) \\
benchmark & RP\textsuperscript{a} & hold-out & 4 & 393 (290--401) & 354.5 (277--392) & 0.218 (0.000--0.378) & 0.000 & 1.000 & 0.746 (0.611--0.766) \\
Kermany & NORMAL & fitting & 1000 & 512 (512--1536) & 496 (496--512) & 0.000 (0.000--0.355) & 0.016 & 1.000 & 0.742 (0.247--0.766) \\
Kermany & NORMAL & test & 250 & 512 (512--768) & 496 (496--496) & 0.000 (0.000--0.290) & 0.022 & 1.000 & 0.742 (0.495--0.742) \\
Kermany & CNV & fitting & 1000 & 512 (512--1536) & 496 (496--496) & 0.000 (0.000--0.423) & 0.015 & 1.000 & 0.742 (0.247--0.742) \\
Kermany & CNV & test & 250 & 512 (512--768) & 496 (496--496) & 0.000 (0.000--0.306) & 0.025 & 1.000 & 0.742 (0.495--0.742) \\
Kermany & DME & fitting & 1000 & 512 (512--1536) & 512 (496--512) & 0.002 (0.000--0.516) & 0.022 & 1.000 & 0.766 (0.247--0.766) \\
Kermany & DME & test & 250 & 512 (512--1536) & 496 (496--496) & 0.000 (0.000--0.226) & 0.018 & 1.000 & 0.742 (0.247--0.742) \\
Kermany & DRUSEN & fitting & 1000 & 512 (512--1536) & 496 (496--496) & 0.000 (0.000--0.419) & 0.021 & 1.000 & 0.742 (0.247--0.742) \\
Kermany & DRUSEN & test & 250 & 512 (512--768) & 496 (496--496) & 0.000 (0.000--0.226) & 0.027 & 1.000 & 0.742 (0.495--0.742) \\
OCTID & CSC & query & 102 & 750 (750--879) & 500 (500--586) & 0.000 (0.000--0.024) & 0.001 & 1.000 & 0.510 (0.510--0.510) \\
OCTID & MH & query & 102 & 750 (750--750) & 500 (500--500) & 0.000 (0.000--0.010) & 0.001 & 1.000 & 0.510 (0.510--0.510) \\
OCTID & NORMAL & query & 206 & 750 (750--879) & 500 (500--586) & 0.000 (0.000--0.000) & 0.004 & 1.000 & 0.510 (0.510--0.510) \\
OCTDL & MH & query & 39 & 912 (738--1146) & 321 (288--412) & 0.000 (0.000--0.000) & 0.001 & 1.000 & 0.280 (0.241--0.326) \\
OCTDL & NORMAL & query & 332 & 882 (614--1270) & 321 (228--536) & 0.000 (0.000--0.000) & 0.001 & 1.000 & 0.278 (0.241--0.402) \\
\bottomrule
\end{tabular}

\par\smallskip
\parbox{\textwidth}{\footnotesize Benchmark roles: development (fitting) and hold-out (query); Kermany: fitting and test; OCTID and OCTDL: query only. \fnRS The folder labels for RP and Stargardt disagree with the counts in the original paper.}
\end{table}

\clearpage
\section*{S-B. Post hoc results (15--21 November 2026)}
Every item below is post hoc and carries the label of its source file. Rules were committed before the numbers (16 November 2026, commit \texttt{1eb3964}; P9 on 18 November, commit \texttt{39c3d73}; P10 on 18 November, commit \texttt{dcbb323}), except P0, P6 and P7, which were first computed before the rules were timestamped. No post hoc result replaces a locked one. P4 and P8 are in the main text, with a summary of P10 (Section~\ref{sec:res-posthoc}); R1 sits beside every false-rare rate (Table~\ref{tab:frr}).
% src: PR l.226-236 (rules 1-4), l.241-256, l.329-354; PR l.275-278; PH CSV posthoc columns; end date: P10_sensitivity.csv computed 2026-11-21 (last post hoc result before the Phase-1 controls); P10 commit dcbb323 (PH P10_*.csv label "(18 Nov 2026)")

\subsection*{P0. Per-class development recall}
\begin{table}[!htbp]
\caption{Post hoc, rule fixed after first look (first computed 15 Nov 2026). Recall per class on the development images, prototype classifier, mean over $K=1$--8 and 100 draws.}
\label{tab:p0}
\centering\footnotesize
% src: PH P0_per_class_recall.csv lines 2-25, column mean_k1_2_4_8; label: post hoc, rule fixed after first look (first computed 15 Nov 2026), recomputed 2026-11-16
\begin{tabular}{@{}lcccc@{}}
\toprule
Class & MIRAGE & DINOv2 & Inception-v3 & RETFound \\
\midrule
CSC & 0.705 & 0.494 & 0.607 & 0.407 \\
MH & 0.540 & 0.612 & 0.623 & 0.302 \\
MacTel & 0.684 & 0.628 & 0.509 & 0.643 \\
Stargardt\textsuperscript{a} & 0.591 & 0.592 & 0.576 & 0.346 \\
RP\textsuperscript{a} & 0.728 & 0.462 & 0.355 & 0.333 \\
\midrule Macro recall & 0.650 & 0.558 & 0.534 & 0.406 \\
\bottomrule
\end{tabular}

\par\smallskip
\parbox{0.6\textwidth}{\footnotesize \fnRS The folder labels for RP and Stargardt disagree with the counts in the original paper; macro recall is unaffected.}
\end{table}
MIRAGE's macular-hole recall (0.540) was below that of DINOv2 (0.612) and Inception-v3 (0.623).
% src: PH P0_per_class_recall.csv lines 15, 9, 21 (mean_k1_2_4_8 0.5405, 0.6117, 0.6235)

\subsection*{P1. Exact bounds and leave-one-image-out}
\begin{table}[!htbp]
\caption{Post hoc, rule committed 16 Nov 2026. Clopper--Pearson 95\% bounds for nine-way counts of scans called any rare class. OCTDL counts are clustered (332 normal images from 110 patients; 39 MH images from 28 patients), and the exact bound treats images as independent.}
\label{tab:p1}
\centering\footnotesize
% src: PH P1_nine_way_exact_bounds.csv rows count=called_any_rare (k, n, cp95_low, cp95_high); label: post hoc, rule timestamped by commit 1eb3964 (16 Nov 2026), computed 2026-11-16
\begin{tabular}{@{}llcccc@{}}
\toprule
Set & True class & MIRAGE & DINOv2 & Inception-v3 & RETFound \\
\midrule
Kermany patient-disjoint & NORMAL & 0/169 (0.000--0.022) & 6/169 (0.013--0.076) & 4/169 (0.006--0.059) & 1/169 (0.000--0.033) \\
Kermany all test & NORMAL & 0/250 (0.000--0.015) & 10/250 (0.019--0.072) & 7/250 (0.011--0.057) & 1/250 (0.000--0.022) \\
OCTID & NORMAL & 4/206 (0.005--0.049) & 17/206 (0.049--0.129) & 15/206 (0.041--0.117) & 0/206 (0.000--0.018) \\
OCTDL & NORMAL & 30/332 (0.062--0.126) & 141/332 (0.371--0.480) & 11/332 (0.017--0.059) & 18/332 (0.032--0.084) \\
OCTID & CSC & 41/102 (0.306--0.504) & 69/102 (0.577--0.766) & 46/102 (0.352--0.553) & 0/102 (0.000--0.036) \\
OCTID & MH & 11/102 (0.055--0.185) & 85/102 (0.747--0.900) & 44/102 (0.334--0.533) & 0/102 (0.000--0.036) \\
OCTDL & MH & 3/39 (0.016--0.209) & 39/39 (0.910--1.000) & 1/39 (0.001--0.135) & 3/39 (0.016--0.209) \\
\bottomrule
\end{tabular}

\end{table}
For MIRAGE's zero count on the 169 patient-disjoint Kermany normals, the exact 95\% upper bound is 0.022. MIRAGE minus Inception-v3, mean development macro recall over $K=1$--8, was $+0.116$ with all 124 development queries; leaving out one query image at a time gave values from 0.109 (without MH/23.PNG) to 0.122 (without Stargardt/19.PNG). On the hold-out the full difference was $+0.197$; leaving out one image at a time gave 0.173 (without MacTel/18.PNG) to 0.232 (without RP/17.PNG). Per-image changes are in \texttt{P1\_leave\_one\_out.csv}.
% src: PH P1_nine_way_exact_bounds.csv line 116 (cp95_high 0.0216); PH P1_leave_one_out_summary.csv lines 2-3 (0.11577; 0.10941 MH/23.PNG; 0.12219 Stargardt/19.PNG; 0.19737; 0.17288 MacTel/18.PNG; 0.23179 RP/17.PNG)

\subsection*{P2. October's trained heads on OCTID and OCTDL (outside the locked study)}
\begin{table}[!htbp]
\caption{Post hoc, rule committed 16 Nov 2026; outside the locked study (rule 4): October's nine-way heads, trained with real rare images only (none), with basic augmentation (basic) or with CycleGAN images (gan), applied unchanged to the locked Inception-v3 features of OCTID and OCTDL. Counts of images.}
\label{tab:p2}
\centering\footnotesize
% src: PH P2_october_heads_external.csv lines 2-16; gates PH P2_gates.csv; label: post hoc, rule timestamped by commit 1eb3964 (16 Nov 2026), computed 2026-11-16; outside the locked study (rule 4): trained with augmented and GAN images
\begin{tabular}{@{}lllrrrrr@{}}
\toprule
Head & Source & True class & $n$ & Correct & NORMAL & \makecell{CNV, DME\\or DRUSEN} & Any rare \\
\midrule
none & OCTID & NORMAL & 206 & 205 & 205 & 1 & 0 \\
 & OCTID & CSC & 102 & 0 & 14 & 88 & 0 \\
 & OCTID & MH & 102 & 4 & 9 & 89 & 4 \\
 & OCTDL & NORMAL & 332 & 256 & 256 & 76 & 0 \\
 & OCTDL & MH & 39 & 0 & 0 & 39 & 0 \\
\addlinespace[2pt]
basic & OCTID & NORMAL & 206 & 206 & 206 & 0 & 0 \\
 & OCTID & CSC & 102 & 5 & 19 & 76 & 7 \\
 & OCTID & MH & 102 & 25 & 12 & 65 & 25 \\
 & OCTDL & NORMAL & 332 & 269 & 269 & 52 & 11 \\
 & OCTDL & MH & 39 & 1 & 0 & 38 & 1 \\
\addlinespace[2pt]
gan & OCTID & NORMAL & 206 & 175 & 175 & 7 & 24 \\
 & OCTID & CSC & 102 & 18 & 7 & 71 & 24 \\
 & OCTID & MH & 102 & 33 & 2 & 64 & 36 \\
 & OCTDL & NORMAL & 332 & 131 & 131 & 114 & 87 \\
 & OCTDL & MH & 39 & 3 & 0 & 34 & 5 \\
\bottomrule
\end{tabular}

\end{table}
All gates passed: each head's class order matched; the heads reproduced 17, 18 and 19 of the 22 hold-out images (none, basic, gan); and the external lock records the same Inception-v3 recipe as the locked arm. The heads also reproduced all 1,022 October test predictions. A failed gate would have been reported, not fixed.
% src: PH P2_gates.csv lines 2-13; PH P2_october_heads_external.csv lines 2-16; PR l.245

\subsection*{P3. Clinic-to-clinic macular-hole agreement (outside the locked study)}
\begin{table}[!htbp]
\caption{Post hoc, rule committed 16 Nov 2026; outside the locked study (rule 4): OCTDL macular holes used as examples although A1 makes them query-only. Agreement between two macular-hole labels, never transfer. Two-way (NORMAL against MH) macro recall, mean over $K=1$--8 and 100 draws, chance 0.50; OCTDL examples at most one image per patient per class; OCTDL queries scored per image and per patient.}
\label{tab:p3}
\centering\footnotesize
% src: PH P3_clinic_to_clinic_mh.csv rows k=mean_k1_2_4_8 (mean, sd, p05, p95, recall_NORMAL_mean, recall_MH_mean);
%      label: post hoc, rule timestamped by commit 1eb3964 (16 Nov 2026), computed 2026-11-16; outside the locked study (rule 4): OCTDL MH used as examples although A1 makes it query-only
\begin{tabular}{@{}lllccc@{}}
\toprule
Arm & Direction & Unit & \makecell{Two-way macro recall\\(SD; 5th--95th)} & NORMAL recall & MH recall \\
\midrule
RETFound & OCTID examples $\to$ OCTDL queries & image & 0.675 (0.063; 0.564--0.773) & 0.512 & 0.837 \\
 & OCTID examples $\to$ OCTDL queries & patient & 0.676 (0.059; 0.574--0.767) & 0.519 & 0.833 \\
 & OCTDL examples $\to$ OCTID queries & image & 0.699 (0.107; 0.520--0.861) & 0.838 & 0.561 \\
\addlinespace[2pt]
DINOv2 & OCTID examples $\to$ OCTDL queries & image & 0.858 (0.076; 0.710--0.964) & 0.854 & 0.863 \\
 & OCTID examples $\to$ OCTDL queries & patient & 0.860 (0.076; 0.706--0.965) & 0.849 & 0.872 \\
 & OCTDL examples $\to$ OCTID queries & image & 0.787 (0.057; 0.674--0.861) & 0.976 & 0.598 \\
\addlinespace[2pt]
MIRAGE & OCTID examples $\to$ OCTDL queries & image & 0.550 (0.038; 0.507--0.612) & 0.130 & 0.971 \\
 & OCTID examples $\to$ OCTDL queries & patient & 0.559 (0.036; 0.514--0.614) & 0.146 & 0.972 \\
 & OCTDL examples $\to$ OCTID queries & image & 0.691 (0.093; 0.550--0.853) & 0.825 & 0.558 \\
\addlinespace[2pt]
Inception-v3 & OCTID examples $\to$ OCTDL queries & image & 0.500 (0.027; 0.461--0.539) & 0.095 & 0.904 \\
 & OCTID examples $\to$ OCTDL queries & patient & 0.506 (0.028; 0.467--0.549) & 0.106 & 0.905 \\
 & OCTDL examples $\to$ OCTID queries & image & 0.507 (0.076; 0.398--0.623) & 0.466 & 0.548 \\
\bottomrule
\end{tabular}

\end{table}
% src: PH P3_clinic_to_clinic_mh.csv rows k=mean_k1_2_4_8; per-K rows and collapsed draws in the same file; rosters P3_example_rosters.csv (released, not shown); PR l.246, l.379-381

\subsection*{P5. Normal-scan destinations and a ``not sure'' cut-off}
\begin{table}[!htbp]
\caption{Post hoc, rule committed 16 Nov 2026. A ``not sure'' cut-off per model, set at the 9th-smallest top-two score gap of the 169 patient-disjoint Kermany normals (8 lie below it), applied unchanged to OCTID and OCTDL normal scans. ``Sure'': gap at or above the cut-off. The Kermany row is calibration only, not a result. Normal recall is shown beside the false-rare rate, never in its place.}
\label{tab:p5}
\centering\footnotesize
% src: PH P5_not_sure_cutoff.csv lines 2-13; label: post hoc, rule timestamped by commit 1eb3964 (16 Nov 2026), computed 2026-11-16
\begin{tabular}{@{}llrrcccc@{}}
\toprule
Arm & Set & Normals & Not sure & \makecell{Called rare\\(all / sure)} & \makecell{Called NORMAL\\(all / sure)} & \makecell{Called common\\(all / sure)} & Role \\
\midrule
RETFound & Kermany patient-disjoint & 169 & 8 & 1 / 0 & 149 / 145 & 19 / 16 & calibration \\
 & OCTID & 206 & 21 & 0 / 0 & 189 / 179 & 17 / 6 & result \\
 & OCTDL & 332 & 10 & 18 / 17 & 137 / 132 & 177 / 173 & result \\
\addlinespace[2pt]
DINOv2 & Kermany patient-disjoint & 169 & 8 & 6 / 5 & 146 / 143 & 17 / 13 & calibration \\
 & OCTID & 206 & 18 & 17 / 10 & 188 / 177 & 1 / 1 & result \\
 & OCTDL & 332 & 63 & 141 / 113 & 145 / 123 & 46 / 33 & result \\
\addlinespace[2pt]
MIRAGE & Kermany patient-disjoint & 169 & 8 & 0 / 0 & 169 / 161 & 0 / 0 & calibration \\
 & OCTID & 206 & 46 & 4 / 0 & 202 / 160 & 0 / 0 & result \\
 & OCTDL & 332 & 295 & 30 / 14 & 26 / 4 & 276 / 19 & result \\
\addlinespace[2pt]
Inception-v3 & Kermany patient-disjoint & 169 & 8 & 4 / 4 & 150 / 145 & 15 / 12 & calibration \\
 & OCTID & 206 & 12 & 15 / 13 & 182 / 175 & 9 / 6 & result \\
 & OCTDL & 332 & 95 & 11 / 9 & 66 / 42 & 255 / 186 & result \\
\bottomrule
\end{tabular}

\end{table}
The top-two gap distributions per destination, source and arm are in \texttt{P5\_gap\_summary.csv}, and per-scan destinations in \texttt{P5\_normal\_scans.csv} (released, not shown).
% src: PH P5_not_sure_cutoff.csv lines 2-13; PR l.248, l.348-349 (9th-smallest gap); P5_gap_summary.csv, P5_normal_scans.csv

\subsection*{P6. Prototype against logistic probe}
\begin{table}[!htbp]
\caption{Post hoc, computed before the rules were timestamped (first computed 16 Nov 2026). Logistic probe minus prototype, five-way macro recall, paired over the first 25 draws: mean (5th to 95th percentile).}
\label{tab:p6}
\centering\footnotesize
% src: PH P6_probe_vs_prototype.csv (diff_mean, diff_p05, diff_p95; logistic minus prototype, first 25 draws); label: post hoc, computed before the rules were timestamped (first computed 16 Nov 2026), recomputed 2026-11-16
\begin{tabular}{@{}lcccccc@{}}
\toprule
 & Development & \multicolumn{5}{c}{Hold-out} \\ \cmidrule(lr){2-2}\cmidrule(l){3-7}
Arm & $K=4$ & $K=1$ & $K=2$ & $K=4$ & $K=8$ & $K=16$ \\
\midrule
MIRAGE & \makecell{+0.012\\{\scriptsize ($-$0.015 to +0.038)}} & \makecell{+0.006\\{\scriptsize ($-$0.050 to +0.088)}} & \makecell{+0.016\\{\scriptsize ($-$0.074 to +0.100)}} & \makecell{+0.016\\{\scriptsize ($-$0.040 to +0.082)}} & \makecell{+0.014\\{\scriptsize ($-$0.032 to +0.082)}} & \makecell{+0.020\\{\scriptsize (0.000 to +0.050)}} \\
DINOv2 & \makecell{+0.029\\{\scriptsize ($-$0.026 to +0.066)}} & \makecell{+0.019\\{\scriptsize ($-$0.048 to +0.098)}} & \makecell{+0.036\\{\scriptsize ($-$0.042 to +0.128)}} & \makecell{+0.041\\{\scriptsize ($-$0.050 to +0.164)}} & \makecell{+0.064\\{\scriptsize (0.000 to +0.140)}} & \makecell{+0.055\\{\scriptsize ($-$0.008 to +0.138)}} \\
Inception-v3 & \makecell{+0.030\\{\scriptsize (+0.003 to +0.068)}} & \makecell{+0.006\\{\scriptsize ($-$0.050 to +0.074)}} & \makecell{+0.030\\{\scriptsize ($-$0.050 to +0.122)}} & \makecell{+0.044\\{\scriptsize ($-$0.036 to +0.140)}} & \makecell{+0.060\\{\scriptsize ($-$0.032 to +0.148)}} & \makecell{+0.072\\{\scriptsize (0.000 to +0.188)}} \\
RETFound & \makecell{+0.001\\{\scriptsize ($-$0.022 to +0.029)}} & \makecell{$-$0.003\\{\scriptsize ($-$0.098 to +0.132)}} & \makecell{+0.004\\{\scriptsize ($-$0.050 to +0.050)}} & \makecell{+0.011\\{\scriptsize ($-$0.082 to +0.090)}} & \makecell{+0.012\\{\scriptsize ($-$0.050 to +0.098)}} & \makecell{+0.018\\{\scriptsize ($-$0.034 to +0.050)}} \\
\bottomrule
\end{tabular}

\end{table}
The same comparison for five-way output shares on OCTID CSC, OCTID MH and OCTDL MH is in \texttt{P6\_probe\_vs\_prototype.csv}.
% src: PH P6_probe_vs_prototype.csv (development K=4; hold-out K=1-16; external rows); PR l.249, l.331-335

\subsection*{P7. Ensemble (one line)}
Post hoc, rule fixed after first look (first computed 15 Nov 2026). A z-scored MIRAGE + DINOv2 ensemble scored 0.695 against 0.650 for MIRAGE in development and 0.718 against 0.722 on the hold-out (mean over $K=1$--8). It is never called a result: the protocol forbids a classifier added to improve a model.
% src: PH P7_ensemble.csv lines 21, 6, 60, 42 (0.6949, 0.6497, 0.7182, 0.7216); PR l.250

\subsection*{P9. Rare against normal, no cut-off}
\begin{table}[!htbp]
\caption{Post hoc, rule committed 18 Nov 2026. Per-scan AUROC of one prototype score (the largest rare cosine minus the largest common cosine), disease scans against normal scans of the same set, with 95\% intervals from 2,000 bootstrap resamples (patients for OCTDL, images for OCTID).}
\label{tab:p9}
\centering\footnotesize
% src: PH P9_rare_vs_normal_auroc.csv lines 2-21 (auroc, boot95_low, boot95_high); label: post hoc, rule timestamped by commit 39c3d73 (18 Nov 2026), computed 2026-11-18
\begin{tabular}{@{}lcccc@{}}
\toprule
Set & MIRAGE & DINOv2 & Inception-v3 & RETFound \\
\midrule
OCTID CSC against OCTID NORMAL & 0.878 (0.832--0.919) & 0.856 (0.801--0.901) & 0.805 (0.749--0.854) & 0.756 (0.679--0.825) \\
OCTID MH against OCTID NORMAL & 0.684 (0.618--0.748) & 0.918 (0.876--0.954) & 0.784 (0.724--0.838) & 0.550 (0.472--0.637) \\
OCTID CSC+MH against OCTID NORMAL & 0.781 (0.733--0.824) & 0.887 (0.854--0.919) & 0.795 (0.749--0.835) & 0.653 (0.589--0.715) \\
OCTDL MH against OCTDL NORMAL & 0.357 (0.261--0.468) & 0.978 (0.949--0.998) & 0.634 (0.543--0.717) & 0.456 (0.358--0.562) \\
\makecell[l]{Hold-out 22 vs 169 Kermany\\normals (source and disease mixed)} & 1.000 (1.000--1.000) & 0.995 (0.986--1.000) & 0.989 (0.967--1.000) & 0.958 (0.872--1.000) \\
\bottomrule
\end{tabular}

\end{table}
The hold-out row mixes source and disease. An AUROC below 0.5 means this score ranks those scans the wrong way; MIRAGE's OCTDL value (0.357) says nothing about what MIRAGE's features contain. The OCTID values hold under the published input path; with black bars added, the same score ranks OCTID the other way (G1). The share of disease scans above the score that only 5\% of each set's normals exceed is read from the test set itself, so it is an upper bound, never a result; it is released in \texttt{P9\_rare\_vs\_normal\_auroc.csv} with the paired MIRAGE differences in \texttt{P9\_rare\_vs\_normal\_differences.csv}.
% src: PH P9_rare_vs_normal_auroc.csv line 16 (0.3574); PR l.252, l.387-390, l.419-421

\subsection*{G1. Pad-to-square input check}
\begin{table}[!htbp]
\caption{Post hoc sensitivity check outside A2 A-11 (rule committed 16 Nov 2026, computed 18 Nov 2026); it never replaces a locked result. MIRAGE and Inception-v3 re-extracted with every image padded with black to a square around its centre, then given each model's own locked transform; scored with the locked episodes and nine-way roster.}
\label{tab:g1}
\centering\footnotesize
% src: PH G1_pad_to_square.csv columns locked, padded, size_only (development lines 6-10/54-58; hold-out 11-21/59-69);
%      K = 16 omitted from any reading (protocols.md 19 Nov G1 decision). Label: post hoc, rule timestamped by commit 1eb3964 (16 Nov 2026), computed 2026-11-18; sensitivity check outside A2 A-11, never replaces a locked result
\begin{tabular}{@{}llcccccc@{}}
\toprule
 & & \multicolumn{3}{c}{MIRAGE} & \multicolumn{3}{c}{Inception-v3} \\ \cmidrule(lr){3-5}\cmidrule(l){6-8}
Set & $K$ & Locked & Padded & Change & Locked & Padded & Change \\
\midrule
Development & 1 & 0.495 & 0.497 & +0.002 & 0.409 & 0.399 & $-$0.010 \\
 & 2 & 0.615 & 0.611 & $-$0.004 & 0.499 & 0.484 & $-$0.015 \\
 & 4 & 0.712 & 0.703 & $-$0.010 & 0.581 & 0.574 & $-$0.007 \\
 & 8 & 0.777 & 0.774 & $-$0.003 & 0.647 & 0.642 & $-$0.005 \\
 & 1--8 & 0.650 & 0.646 & $-$0.004 & 0.534 & 0.525 & $-$0.009 \\
\addlinespace[2pt]
Hold-out & 1 & 0.536 & 0.518 & $-$0.018 & 0.430 & 0.427 & $-$0.003 \\
 & 2 & 0.685 & 0.662 & $-$0.023 & 0.497 & 0.500 & +0.003 \\
 & 4 & 0.799 & 0.788 & $-$0.011 & 0.556 & 0.559 & +0.003 \\
 & 8 & 0.867 & 0.861 & $-$0.005 & 0.614 & 0.618 & +0.004 \\
 & 1--8 & 0.722 & 0.707 & $-$0.014 & 0.524 & 0.526 & +0.002 \\
\midrule
 & & \multicolumn{6}{c}{Nine-way: scans called NORMAL / called any rare (locked $\to$ padded)} \\
\multicolumn{2}{@{}l}{Kermany patient-disjoint (169)} & \multicolumn{3}{c}{169 $\to$ 166 / 0 $\to$ 2} & \multicolumn{3}{c}{150 $\to$ 157 / 4 $\to$ 4} \\
\multicolumn{2}{@{}l}{OCTID NORMAL (206)} & \multicolumn{3}{c}{202 $\to$ 0 / 4 $\to$ 206} & \multicolumn{3}{c}{182 $\to$ 8 / 15 $\to$ 198} \\
\multicolumn{2}{@{}l}{OCTDL NORMAL (332)} & \multicolumn{3}{c}{26 $\to$ 1 / 30 $\to$ 306} & \multicolumn{3}{c}{66 $\to$ 13 / 11 $\to$ 302} \\
\bottomrule
\end{tabular}

\end{table}
Development and hold-out recall moved by at most 0.023 at $K=1$--8 (mean of 100 draws), a near-null test because the rare inputs are already near square. With black padding, external normal scans went almost entirely to rare classes for both models (MIRAGE: OCTID 202 $\to$ 0 and OCTDL 26 $\to$ 1 called NORMAL). The rare images are letterboxed, and padding gives wide external scans that layout. External nine-way results depend on input geometry. We draw no conclusion about cause, robustness to aspect ratio, a ranking of the two models on padded inputs, or what the features contain, and we read no $K=16$ value. Checks: unpadded controls and 1,032 square images were within $1.43\times10^{-6}$ of the locked features, and repeats were exact; the GPU differed (L4 for the lock, A100 for G1), and the MIRAGE environment differed only in the patch versions of \texttt{filelock} and \texttt{regex}. Full tables, including padded P9: \texttt{G1\_comparison.md}.
% src: PH G1_pad_to_square.csv; G1_comparison.md sections 1-5; PR l.401-422 (allowed and forbidden claims)

\subsection*{P10. Same-picture audit}
Post hoc, rule committed 18 Nov 2026 (commit \texttt{dcbb323}); screen and review sheet computed 19 Nov 2026; the reviewer's marks (SHA256 \texttt{15a9606fce87\ldots}) committed on 21 Nov 2026 before the review key was joined; sensitivity computed 21 Nov 2026. Method in Section~\ref{sec:posthoc-methods}; summary in Section~\ref{sec:res-posthoc}. The SIFT screen scored all 10,585 pairs of the 146 rare images (median 4 inliers, largest 869); the cut-off, the largest score among 2,000 rare--Kermany pairs (1,617 Kermany images), was 18 inliers, and all 146 synthetic copies reached it. Thirty-five pairs reached the cut-off; with 10 synthetic copies and 10 decoys they made the 55 review items. The reviewer, one author who is not a clinician, marked all 10 synthetic copies ``same picture'' and all 10 decoys ``different''. Of the 35 screened pairs, 14 were marked ``same picture'', 1 ``same case, other slice'', 1 ``unsure'' and 19 ``different'' (Table~\ref{tab:p10pairs}). No pair shared a source link: the source list of Yoo et al.\ names no file. Per-pair scores are in \texttt{P10\_screen\_pairs.csv}, the review key and marks in \texttt{P10\_review\_key.csv} and \texttt{P10\_review\_marks.csv} (released; no image is released).
% src: PH P10_screen_summary.csv lines 2-13 (10585; 2000; 1617; cut-off 18; 35; 35; 146; 1.000; median 4.0; max 869; "0 (S1 names no file)"; 55); P10_screen_synthetic.csv (146 of 146 at_or_above_cutoff); marks SHA256 15a9606fce8779b1... (sha256sum of P10_review_marks.csv; label of P10_sensitivity.csv), committed abb1770 (PR@fec827e l.485-486); P10_review_key.csv joined with P10_review_marks.csv (pair: 14 same picture, 1 same case, 1 unsure, 19 different; synthetic 10 same picture; decoy 10 different)

\begin{table}[!htbp]
\caption{Post hoc (P10; rule committed 18 Nov 2026, marks committed 21 Nov 2026). Image pairs marked as a match in the blind review, with the set of each image and the keypoint score (RANSAC inliers; cut-off 18). Images are named by their released folder and file.}
\label{tab:p10pairs}
\centering\footnotesize
% src: PH P10_matched_pairs.csv (pairs, sets); mark and screen inliers: PH P10_review_key.csv joined with
%      PH P10_review_marks.csv on item (synthetic copies excluded, as in fm_posthoc.p10_matched_pairs);
%      label: post hoc, rule timestamped by commit dcbb323 (18 Nov 2026), computed 2026-11-21; marks SHA256 15a9606fce87
\begin{tabular}{@{}rllllr@{}}
\toprule
 & Image & Set & Image & Set & Inliers \\
\midrule
\multicolumn{6}{@{}l}{\emph{Same picture (14)}} \\
1 & MH/01.PNG & development & MH/08.PNG & development & 182 \\
2 & MH/01.PNG & development & MH/14.PNG & development & 193 \\
3 & MH/08.PNG & development & MH/14.PNG & development & 196 \\
4 & MH/12.PNG & development & MH/25.PNG & development & 288 \\
5 & MacTel/01.PNG & development & MacTel/19.PNG & \textbf{hold-out} & 265 \\
6 & MacTel/04.PNG & development & MacTel/06.PNG & development & 869 \\
7 & MacTel/07.PNG & development & MacTel/22.PNG & development & 369 \\
8 & RP/03.PNG & development & RP/20.PNG & \textbf{hold-out} & 203 \\
9 & RP/06.PNG & development & RP/30.PNG & development & 202 \\
10 & RP/11.PNG & development & RP/17.PNG & \textbf{hold-out} & 104 \\
11 & Stargardt/03.PNG & \textbf{hold-out} & Stargardt/08.PNG & development & 208 \\
12 & Stargardt/03.PNG & \textbf{hold-out} & Stargardt/14.PNG & development & 234 \\
13 & Stargardt/08.PNG & development & Stargardt/14.PNG & development & 276 \\
14 & Stargardt/09.PNG & development & Stargardt/18.PNG & development & 342 \\
\multicolumn{6}{@{}l}{\emph{Same case, other slice (1)}} \\
15 & MacTel/09.PNG & development & MacTel/23.PNG & development & 18 \\
\multicolumn{6}{@{}l}{\emph{Unsure (counted only in the second exclusion set) (1)}} \\
16 & CSC/26.PNG & development & CSC/27.PNG & development & 18 \\
\bottomrule
\end{tabular}

\par\smallskip
\parbox{0.75\textwidth}{\footnotesize Hold-out images in bold: MacTel/19, RP/17, RP/20 and Stargardt/03 have a same-picture copy among the development images. RP/20 is an addition in version 2 of the release. The unsure pair carries the reviewer's note ``probably a follow-up visit'' (a non-clinician's reading). The locked duplicate screen (identical files, identical pixels, difference hash at Hamming distance at most 4) flagged none of these pairs. A screen for copies of one picture cannot rule out other slices or visits of the same patient.}
\end{table}
% src: tables/p10_pairs.tex (PH P10_matched_pairs.csv; P10_review_key.csv; P10_review_marks.csv); PR@fec827e l.488-491; unsure note: P10_review_marks.csv item 11

\begin{table}[!htbp]
\caption{Post hoc (P10; rule committed 18 Nov 2026, computed 21 Nov 2026). Five-way macro recall, mean over $K=1$--8 (prototype classifier, locked episodes), locked and recomputed without the marked matches: in development, without each query that shares a mark with a support of its episode; on the hold-out, without each hold-out image marked against any development image. Beside the locked numbers, never replacing them.}
\label{tab:p10}
\centering\footnotesize
\setlength{\tabcolsep}{4pt}
\resizebox{\textwidth}{!}{% src: PH P10_sensitivity.csv rows k=mean_k1_2_4_8 (mean, sd, p05, p95) and column excluded; label: post hoc, rule timestamped by commit dcbb323 (18 Nov 2026), computed 2026-11-21; marks SHA256 15a9606fce87
%      arm rows: mean (SD); paired rows ('<arm> minus Inception'): mean (p05 to p95). Locked columns equal R2 / the locked archive.
\begin{tabular}{@{}lcccccc@{}}
\toprule
 & \multicolumn{3}{c}{Development (124 queries), mean $K$ 1--8} & \multicolumn{3}{c}{Hold-out (22 images), mean $K$ 1--8} \\
\cmidrule(lr){2-4}\cmidrule(l){5-7}
Arm & Locked & \makecell{Without same\\picture or case} & \makecell{Also without\\unsure} & Locked & \makecell{Without same\\picture or case} & \makecell{Also without\\unsure} \\
\midrule
MIRAGE (nominee) & \makecell{0.650\\{\scriptsize (0.028)}} & \makecell{0.643\\{\scriptsize (0.030)}} & \makecell{0.642\\{\scriptsize (0.030)}} & \makecell{0.722\\{\scriptsize (0.063)}} & \makecell{0.751\\{\scriptsize (0.066)}} & \makecell{0.751\\{\scriptsize (0.066)}} \\
DINOv2 & \makecell{0.558\\{\scriptsize (0.028)}} & \makecell{0.551\\{\scriptsize (0.029)}} & \makecell{0.550\\{\scriptsize (0.029)}} & \makecell{0.526\\{\scriptsize (0.065)}} & \makecell{0.511\\{\scriptsize (0.078)}} & \makecell{0.511\\{\scriptsize (0.078)}} \\
Inception-v3 (reference) & \makecell{0.534\\{\scriptsize (0.029)}} & \makecell{0.530\\{\scriptsize (0.029)}} & \makecell{0.529\\{\scriptsize (0.029)}} & \makecell{0.524\\{\scriptsize (0.067)}} & \makecell{0.511\\{\scriptsize (0.090)}} & \makecell{0.511\\{\scriptsize (0.090)}} \\
RETFound & \makecell{0.406\\{\scriptsize (0.024)}} & \makecell{0.400\\{\scriptsize (0.025)}} & \makecell{0.398\\{\scriptsize (0.025)}} & \makecell{0.338\\{\scriptsize (0.065)}} & \makecell{0.300\\{\scriptsize (0.071)}} & \makecell{0.300\\{\scriptsize (0.071)}} \\
\midrule
\makecell[l]{MIRAGE $-$\\Inception-v3 (paired)} & \makecell{+0.116\\{\scriptsize (+0.069 to +0.185)}} & \makecell{+0.113\\{\scriptsize (+0.064 to +0.185)}} & \makecell{+0.113\\{\scriptsize (+0.064 to +0.184)}} & \makecell{+0.197\\{\scriptsize (+0.042 to +0.343)}} & \makecell{+0.240\\{\scriptsize (+0.089 to +0.407)}} & \makecell{+0.240\\{\scriptsize (+0.089 to +0.407)}} \\
\makecell[l]{DINOv2 $-$\\Inception-v3 (paired)} & \makecell{+0.024\\{\scriptsize ($-$0.025 to +0.086)}} & \makecell{+0.021\\{\scriptsize ($-$0.032 to +0.083)}} & \makecell{+0.021\\{\scriptsize ($-$0.032 to +0.083)}} & \makecell{+0.002\\{\scriptsize ($-$0.136 to +0.124)}} & \makecell{0.000\\{\scriptsize ($-$0.149 to +0.162)}} & \makecell{0.000\\{\scriptsize ($-$0.149 to +0.162)}} \\
\makecell[l]{RETFound $-$\\Inception-v3 (paired)} & \makecell{$-$0.128\\{\scriptsize ($-$0.186 to $-$0.070)}} & \makecell{$-$0.130\\{\scriptsize ($-$0.188 to $-$0.070)}} & \makecell{$-$0.131\\{\scriptsize ($-$0.189 to $-$0.073)}} & \makecell{$-$0.186\\{\scriptsize ($-$0.308 to $-$0.052)}} & \makecell{$-$0.211\\{\scriptsize ($-$0.398 to $-$0.060)}} & \makecell{$-$0.211\\{\scriptsize ($-$0.398 to $-$0.060)}} \\
\midrule
Excluded & none & \makecell{1,196 query--episode\\exclusions per arm} & \makecell{1,378 query--episode\\exclusions per arm} & none & \makecell{4 images} & \makecell{4 images} \\
\bottomrule
\end{tabular}
}
\par\smallskip
\parbox{0.9\textwidth}{\footnotesize Arm rows: mean over 100 draws (SD). Paired rows: difference from Inception-v3 on the same draws, mean (5th to 95th percentile). ``Without same picture or case'': the 14 same-picture pairs and the same-case pair; ``also without unsure'': the unsure pair counted as a match too. Locked columns are the locked values (Table~\ref{tab:dev}; hold-out, Table~\ref{tab:holdout}). Spread over support draws, not confidence intervals; the hold-out shows direction only. Per-$K$ rows are in \texttt{P10\_sensitivity.csv}.}
\end{table}
% src: tables/p10_sensitivity.tex (PH P10_sensitivity.csv rows k = mean_k1_2_4_8)

\section*{S-C. Released outside any rule}
The files below were computed outside any committed rule and are released under rule 3; they are not results.
% src: pre_rule/README.md; PR l.393-397, l.435-441

\textbf{Provenance.} \texttt{pre\_rule/2026-11-18/v1\_to\_ours.csv} maps the 119 files of Mendeley version 1 to our 146 by file and pixel hash (118 pixel matches); \texttt{md\_versions.json} is the dataset's version list.
% src: v1_to_ours.csv (119 rows, 118 pixel matches); PR l.59-63, l.435-438

\textbf{Per-image layout (pre-rule count, 18 November 2026; computed outside any rule, not a result).} In a per-image count made outside any rule, 126 of the 146 rare images have a dark flat border covering at least 25\% of the image, mostly bands above and below the scan (\texttt{rare\_layout.csv}, scripts \texttt{geom\_a.py} and \texttt{geom\_b.py}). A figure of 128 quoted in an earlier discussion could not be reproduced and is withdrawn.
% src: PR l.71-74, l.439-441; pre_rule/2026-11-18/rare_layout.csv (not recounted). Placement: agent default (PR HEAD 19 Nov entry: "the pre-rule count of 126 appears only in the supplement with its label"); note pre_rule/README.md l.5-6 says these files may be quoted only to say they exist

% ------------------------------------------------------------------ Phase-1 controls E1-E4
% Every block below is typeset only when make_tables.py has written its table from the E CSVs (\havetable, main.tex).
% Default placement (PR@fec827e l.603-605): E1d, E1g and the E2 summary in the main text (post hoc subsection,
% part "Phase-1 controls E1--E4"); the rest here. The OCTDL locked-feature E2d table, which the E2 rule (l.588)
% fixes as main text, is in main.tex (tab:e2d), placed there by Lam on 21 Nov (agent-loop round 1, finding 3).
\section*{S-D. Phase-1 controls E1--E4 (post hoc; rules committed 20 November 2026)}
Methods in Section~\ref{sec:phase1}. Every E output is post hoc, outside the locked protocol, and never replaces a locked result; failed checks are released. Full per-$K$ rows are in the CSVs.
\havetable{e1_label}\ifgen\else
\todo{no E1--E4 result exists yet; the tables appear here once \texttt{make\_tables.py} has read the E CSVs.}
\fi
% src: PR@fec827e l.516-525 (rules 1-2), l.603-605 (placement)

\havetable{e1_label}\ifgen
\subsection*{E1. Size and framing control}
Label of every E1 row: post hoc, Phase-1 rule E1 timestamped by commit 4b4e883 (20 Nov 2026), computed 2026-11-21; outside the locked protocol; never replaces a locked result
. Limitations stated with every E1 result: the centre square drops lateral retina unequally by source and class; MIRAGE now gets 256 px of detail enlarged to 512; E1-M does not remove effective resolution (content short side 103-496 px), field of view, scanner appearance, compression, intensity, black-clipped background against speckle, or Kermany's white fill wedges.

\fi
\havetable{e1a_checks}\ifgen
\subsubsection*{E1a. Checks}
{\scriptsize\setlength{\tabcolsep}{3pt}% src: E1a_checks.csv in C:/Users/ThinkBook/.t3/worktrees/Few-shot/t3-03f3a2c4/results/fm_v1/posthoc (generated by make_tables.py; formatting only);
% label: post hoc, Phase-1 rule E1 timestamped by commit 4b4e883 (20 Nov 2026), computed 2026-11-21; outside the locked protocol; never replaces a locked result; limitations: the centre square drops lateral retina unequally by source and class; MIRAGE now gets 256 px of detail enlarged to 512; E1-M does not remove effective resolution (content short side 103-496 px), field of view, scanner appearance, compression, intensity, black-clipped background against speckle, or Kermany's white fill wedges
\begin{longtable}{@{}p{3.3cm}lp{1.5cm}p{3.6cm}p{2.2cm}p{1.6cm}p{3.2cm}@{}}
\toprule
Check & Arm & Set & Quantity & Value & Status & Note \\
\midrule
\endfirsthead
\toprule
Check & Arm & Set & Quantity & Value & Status & Note \\
\midrule
\endhead
\bottomrule
\endfoot
\bottomrule
\endlastfoot
rule 1: Phase-1 subsection SHA256 &  &  & at rule commit 4b4e883 & 8b4bb0b22b0b\ldots{} & same & working 8b4bb0b22b0b\ldots{} \\
archive receipt &  &  & phase1-E1-C-DINOv2.zip & 49ed95aaa863\ldots{} & equal &  \\
archive receipt &  &  & phase1-E1-C-Inception.zip & ca5ab96d2bc6\ldots{} & equal &  \\
archive receipt &  &  & phase1-E1-C-MIRAGE.zip & 78fdeba16998\ldots{} & equal &  \\
archive receipt &  &  & phase1-E1-C-RETFound.zip & 0a799bb35684\ldots{} & equal &  \\
archive receipt &  &  & phase1-E1-L-DINOv2.zip & 984f561da2bd\ldots{} & equal &  \\
archive receipt &  &  & phase1-E1-L-Inception.zip & 1d4872d6f4df\ldots{} & equal &  \\
archive receipt &  &  & phase1-E1-L-MIRAGE.zip & d1bcbe96a9ca\ldots{} & equal &  \\
archive receipt &  &  & phase1-E1-L-RETFound.zip & 63bf4fff749a\ldots{} & equal &  \\
archive receipt &  &  & phase1-E1-M-DINOv2.zip & ee6986ba2961\ldots{} & equal &  \\
archive receipt &  &  & phase1-E1-M-Inception.zip & ab5d23be8b44\ldots{} & equal &  \\
archive receipt &  &  & phase1-E1-M-MIRAGE.zip & fe645ceea8ed\ldots{} & equal &  \\
archive receipt &  &  & phase1-E1-M-RETFound.zip & 7782c610a9f6\ldots{} & equal &  \\
archive receipt &  &  & phase1-E3-MIRAGE-E3-ref.zip & 32540b6febf2\ldots{} & equal &  \\
archive receipt &  &  & phase1-E3-MIRAGE-M-bicubic.zip & 413f5731fcd0\ldots{} & equal &  \\
archive receipt &  &  & phase1-E3-MIRAGE-M-crop.zip & 75db7f8db501\ldots{} & equal &  \\
archive receipt &  &  & phase1-E3-MIRAGE-M-nostretch.zip & a7da4833b4be\ldots{} & equal &  \\
archive receipt &  &  & phase1-E3-MIRAGE-M-res256.zip & 8b4612ffcc88\ldots{} & equal &  \\
archive receipt &  &  & phase1-E3-RETFound-E3-ref.zip & 630dba012456\ldots{} & equal &  \\
archive receipt &  &  & phase1-E3-RETFound-R-bilinear.zip & 378c47733e32\ldots{} & equal &  \\
archive receipt &  &  & phase1-E3-RETFound-R-norm05.zip & e532ab091719\ldots{} & equal &  \\
archive receipt &  &  & phase1-E3-RETFound-R-pad.zip & 85cd9607683e\ldots{} & equal &  \\
archive receipt &  &  & phase1-E3-RETFound-R-res128.zip & a92e3b2cf15b\ldots{} & equal &  \\
archive receipt &  &  & phase1-E3-RETFound-R-squash.zip & 18ef66354419\ldots{} & equal &  \\
archive receipt &  &  & phase1-E3-RETFound-R-stretch.zip & c67a23dc6042\ldots{} & equal &  \\
archive receipt &  &  & phase1-render.zip & f28e8cb0f494\ldots{} & equal &  \\
control: locked path against locked features & DINOv2 & E1-C & max\_difference\_ratio & $2.8\times10^{-6}$ & pass & 5927 images; min cosine 1.000000 \\
repeat (recorded, not a gate) & DINOv2 & E1-C & max\_difference\_ratio & 0 &  & 5 images \\
identity & DINOv2 & E1-C & code\_matches\_lock & True &  &  \\
identity & DINOv2 & E1-C & environment\_matches\_lock & False &  &  \\
control: locked path against locked features & Inception-v3 & E1-C & max\_difference\_ratio & $1.5\times10^{-6}$ & pass & 20 images; min cosine 1.000000 \\
repeat (recorded, not a gate) & Inception-v3 & E1-C & max\_difference\_ratio & 0 &  & 5 images \\
identity & Inception-v3 & E1-C & code\_matches\_lock & True &  &  \\
identity & Inception-v3 & E1-C & environment\_matches\_lock & True &  &  \\
control: locked path against locked features & MIRAGE & E1-C & max\_difference\_ratio & $1.4\times10^{-7}$ & pass & 20 images; min cosine 1.000000 \\
repeat (recorded, not a gate) & MIRAGE & E1-C & max\_difference\_ratio & 0 &  & 5 images \\
identity & MIRAGE & E1-C & code\_matches\_lock & True &  &  \\
identity & MIRAGE & E1-C & environment\_matches\_lock & False &  &  \\
control: locked path against locked features & RETFound & E1-C & max\_difference\_ratio & $6.5\times10^{-8}$ & pass & 20 images; min cosine 1.000000 \\
repeat (recorded, not a gate) & RETFound & E1-C & max\_difference\_ratio & 0 &  & 5 images \\
identity & RETFound & E1-C & code\_matches\_lock & True &  &  \\
identity & RETFound & E1-C & environment\_matches\_lock & False &  &  \\
control: locked path against locked features & DINOv2 & E1-L & max\_difference\_ratio & $2.8\times10^{-6}$ & pass & 5927 images; min cosine 1.000000 \\
repeat (recorded, not a gate) & DINOv2 & E1-L & max\_difference\_ratio & 0 &  & 5 images \\
identity & DINOv2 & E1-L & code\_matches\_lock & True &  &  \\
identity & DINOv2 & E1-L & environment\_matches\_lock & False &  &  \\
control: locked path against locked features & Inception-v3 & E1-L & max\_difference\_ratio & $1.7\times10^{-6}$ & pass & 20 images; min cosine 1.000000 \\
repeat (recorded, not a gate) & Inception-v3 & E1-L & max\_difference\_ratio & 0 &  & 5 images \\
identity & Inception-v3 & E1-L & code\_matches\_lock & True &  &  \\
identity & Inception-v3 & E1-L & environment\_matches\_lock & True &  &  \\
control: locked path against locked features & MIRAGE & E1-L & max\_difference\_ratio & $1.4\times10^{-7}$ & pass & 20 images; min cosine 1.000000 \\
repeat (recorded, not a gate) & MIRAGE & E1-L & max\_difference\_ratio & 0 &  & 5 images \\
identity & MIRAGE & E1-L & code\_matches\_lock & True &  &  \\
identity & MIRAGE & E1-L & environment\_matches\_lock & False &  &  \\
control: locked path against locked features & RETFound & E1-L & max\_difference\_ratio & $6.5\times10^{-8}$ & pass & 20 images; min cosine 1.000000 \\
repeat (recorded, not a gate) & RETFound & E1-L & max\_difference\_ratio & 0 &  & 5 images \\
identity & RETFound & E1-L & code\_matches\_lock & True &  &  \\
identity & RETFound & E1-L & environment\_matches\_lock & False &  &  \\
control: locked path against locked features & DINOv2 & E1-M & max\_difference\_ratio & $2.8\times10^{-6}$ & pass & 5927 images; min cosine 1.000000 \\
repeat (recorded, not a gate) & DINOv2 & E1-M & max\_difference\_ratio & 0 &  & 5 images \\
identity & DINOv2 & E1-M & code\_matches\_lock & True &  &  \\
identity & DINOv2 & E1-M & environment\_matches\_lock & False &  &  \\
control: locked path against locked features & Inception-v3 & E1-M & max\_difference\_ratio & $1.0\times10^{-6}$ & pass & 20 images; min cosine 1.000000 \\
repeat (recorded, not a gate) & Inception-v3 & E1-M & max\_difference\_ratio & 0 &  & 5 images \\
identity & Inception-v3 & E1-M & code\_matches\_lock & True &  &  \\
identity & Inception-v3 & E1-M & environment\_matches\_lock & True &  &  \\
control: locked path against locked features & MIRAGE & E1-M & max\_difference\_ratio & $1.4\times10^{-7}$ & pass & 20 images; min cosine 1.000000 \\
repeat (recorded, not a gate) & MIRAGE & E1-M & max\_difference\_ratio & 0 &  & 5 images \\
identity & MIRAGE & E1-M & code\_matches\_lock & True &  &  \\
identity & MIRAGE & E1-M & environment\_matches\_lock & False &  &  \\
control: locked path against locked features & RETFound & E1-M & max\_difference\_ratio & $6.5\times10^{-8}$ & pass & 20 images; min cosine 1.000000 \\
repeat (recorded, not a gate) & RETFound & E1-M & max\_difference\_ratio & 0 &  & 5 images \\
identity & RETFound & E1-M & code\_matches\_lock & True &  &  \\
identity & RETFound & E1-M & environment\_matches\_lock & False &  &  \\
E3-ref full control & MIRAGE &  & max\_difference\_ratio & $3.0\times10^{-7}$ & pass & 5927 images \\
E3-ref full control & RETFound &  & max\_difference\_ratio & $1.9\times10^{-7}$ & pass & 5927 images \\
Inception-v3 decoder change (PIL-decoded PNG originals) & Inception-v3 & E1-M & Kermany & 0.025 &  & 20 images; min cosine 0.999642 \\
Inception-v3 decoder change (PIL-decoded PNG originals) & Inception-v3 & E1-M & OCTID & 0.020 &  & 10 images; min cosine 0.999682 \\
Inception-v3 decoder change (PIL-decoded PNG originals) & Inception-v3 & E1-M & OCTDL & 0.046 &  & 10 images; min cosine 0.997780 \\
Inception-v3 resize identity & Inception-v3 & E1-M & tf.image.resize & \{"arrays": 6, "failures": 0, "unchanged": true\} &  &  \\
render &  &  & images rendered (twice; saved PNG pixel hashes verified) & 5927 &  &  \\
render &  &  & content box fallback to the full frame (flagged) & 0 &  &  \\
render &  &  & files with alpha below 255 (dropped by convert('RGB')) & 0 &  &  \\
render &  &  & mode L, format JPEG & 5000 &  &  \\
render &  &  & mode RGB, format JPEG & 708 &  &  \\
render &  &  & mode RGB, format PNG & 11 &  &  \\
render &  &  & mode RGBA, format PNG & 208 &  &  \\
render & Kermany &  & content short-side range (px) & 240--512 &  & 5000 images \\
render & OCTDL &  & content short-side range (px) & 228--536 &  & 371 images \\
render & OCTID &  & content short-side range (px) & 488--586 &  & 410 images \\
render & Yoo rare &  & content short-side range (px) & 103--471 &  & 146 images \\
render & CSC & E1-L & content short side below 128 px & 2 &  &  \\
render & MH & E1-L & content short side below 128 px & 4 &  &  \\
render & Stargardt\textsuperscript{a} & E1-L & content short side below 128 px & 1 &  &  \\
framing residual & CSC & original & mean A2 border share & 0.452 &  & n=32 \\
framing residual & CSC & original & share with A2 border share > 0 & 1.000 &  & n=32 \\
framing residual & CSC & original & share with near-black edge run > 5\% of side & 1.000 &  & n=32 \\
framing residual & CSC & original & share with near-black edge run > 10\% of side & 0.969 &  & n=32 \\
framing residual & CSC & original & mean near-white share & $4.9\times10^{-4}$ &  & n=32 \\
framing residual & CSC & E1-M & mean A2 border share & 0.044 &  & n=32 \\
framing residual & CSC & E1-M & share with A2 border share > 0 & 0.500 &  & n=32 \\
framing residual & CSC & E1-M & share with near-black edge run > 5\% of side & 0.469 &  & n=32 \\
framing residual & CSC & E1-M & share with near-black edge run > 10\% of side & 0.313 &  & n=32 \\
framing residual & CSC & E1-M & mean near-white share & $1.4\times10^{-4}$ &  & n=32 \\
framing residual & CSC & E1-L & mean A2 border share & 0.052 &  & n=32 \\
framing residual & CSC & E1-L & share with A2 border share > 0 & 0.594 &  & n=32 \\
framing residual & CSC & E1-L & share with near-black edge run > 5\% of side & 0.500 &  & n=32 \\
framing residual & CSC & E1-L & share with near-black edge run > 10\% of side & 0.375 &  & n=32 \\
framing residual & CSC & E1-L & mean near-white share & $7.2\times10^{-5}$ &  & n=32 \\
framing residual & CSC & E1-C & mean A2 border share & 0 &  & n=32 \\
framing residual & CSC & E1-C & share with A2 border share > 0 & 0 &  & n=32 \\
framing residual & CSC & E1-C & share with near-black edge run > 5\% of side & 0.375 &  & n=32 \\
framing residual & CSC & E1-C & share with near-black edge run > 10\% of side & 0.125 &  & n=32 \\
framing residual & CSC & E1-C & mean near-white share & $6.6\times10^{-4}$ &  & n=32 \\
framing residual & Kermany & original & mean A2 border share & 0.023 &  & n=5000 \\
framing residual & Kermany & original & share with A2 border share > 0 & 0.465 &  & n=5000 \\
framing residual & Kermany & original & share with near-black edge run > 5\% of side & 0.005 &  & n=5000 \\
framing residual & Kermany & original & share with near-black edge run > 10\% of side & 0.002 &  & n=5000 \\
framing residual & Kermany & original & mean near-white share & 0.042 &  & n=5000 \\
framing residual & Kermany & E1-M & mean A2 border share & 0.003 &  & n=5000 \\
framing residual & Kermany & E1-M & share with A2 border share > 0 & 0.306 &  & n=5000 \\
framing residual & Kermany & E1-M & share with near-black edge run > 5\% of side & 0.007 &  & n=5000 \\
framing residual & Kermany & E1-M & share with near-black edge run > 10\% of side & 0.003 &  & n=5000 \\
framing residual & Kermany & E1-M & mean near-white share & 0.018 &  & n=5000 \\
framing residual & Kermany & E1-L & mean A2 border share & 0.005 &  & n=5000 \\
framing residual & Kermany & E1-L & share with A2 border share > 0 & 0.313 &  & n=5000 \\
framing residual & Kermany & E1-L & share with near-black edge run > 5\% of side & 0.025 &  & n=5000 \\
framing residual & Kermany & E1-L & share with near-black edge run > 10\% of side & 0.011 &  & n=5000 \\
framing residual & Kermany & E1-L & mean near-white share & 0.015 &  & n=5000 \\
framing residual & Kermany & E1-C & mean A2 border share & 0 &  & n=5000 \\
framing residual & Kermany & E1-C & share with A2 border share > 0 & 0 &  & n=5000 \\
framing residual & Kermany & E1-C & share with near-black edge run > 5\% of side & 0.003 &  & n=5000 \\
framing residual & Kermany & E1-C & share with near-black edge run > 10\% of side & 0.001 &  & n=5000 \\
framing residual & Kermany & E1-C & mean near-white share & 0.021 &  & n=5000 \\
framing residual & MH & original & mean A2 border share & 0.421 &  & n=31 \\
framing residual & MH & original & share with A2 border share > 0 & 1.000 &  & n=31 \\
framing residual & MH & original & share with near-black edge run > 5\% of side & 1.000 &  & n=31 \\
framing residual & MH & original & share with near-black edge run > 10\% of side & 0.968 &  & n=31 \\
framing residual & MH & original & mean near-white share & 0.001 &  & n=31 \\
framing residual & MH & E1-M & mean A2 border share & 0.018 &  & n=31 \\
framing residual & MH & E1-M & share with A2 border share > 0 & 0.323 &  & n=31 \\
framing residual & MH & E1-M & share with near-black edge run > 5\% of side & 0.161 &  & n=31 \\
framing residual & MH & E1-M & share with near-black edge run > 10\% of side & 0.097 &  & n=31 \\
framing residual & MH & E1-M & mean near-white share & 0.001 &  & n=31 \\
framing residual & MH & E1-L & mean A2 border share & 0.024 &  & n=31 \\
framing residual & MH & E1-L & share with A2 border share > 0 & 0.387 &  & n=31 \\
framing residual & MH & E1-L & share with near-black edge run > 5\% of side & 0.161 &  & n=31 \\
framing residual & MH & E1-L & share with near-black edge run > 10\% of side & 0.097 &  & n=31 \\
framing residual & MH & E1-L & mean near-white share & $7.6\times10^{-4}$ &  & n=31 \\
framing residual & MH & E1-C & mean A2 border share & 0 &  & n=31 \\
framing residual & MH & E1-C & share with A2 border share > 0 & 0 &  & n=31 \\
framing residual & MH & E1-C & share with near-black edge run > 5\% of side & 0.129 &  & n=31 \\
framing residual & MH & E1-C & share with near-black edge run > 10\% of side & 0.065 &  & n=31 \\
framing residual & MH & E1-C & mean near-white share & 0.002 &  & n=31 \\
framing residual & MacTel & original & mean A2 border share & 0.413 &  & n=29 \\
framing residual & MacTel & original & share with A2 border share > 0 & 1.000 &  & n=29 \\
framing residual & MacTel & original & share with near-black edge run > 5\% of side & 1.000 &  & n=29 \\
framing residual & MacTel & original & share with near-black edge run > 10\% of side & 1.000 &  & n=29 \\
framing residual & MacTel & original & mean near-white share & $2.5\times10^{-4}$ &  & n=29 \\
framing residual & MacTel & E1-M & mean A2 border share & 0.016 &  & n=29 \\
framing residual & MacTel & E1-M & share with A2 border share > 0 & 0.172 &  & n=29 \\
framing residual & MacTel & E1-M & share with near-black edge run > 5\% of side & 0.276 &  & n=29 \\
framing residual & MacTel & E1-M & share with near-black edge run > 10\% of side & 0.241 &  & n=29 \\
framing residual & MacTel & E1-M & mean near-white share & $2.1\times10^{-4}$ &  & n=29 \\
framing residual & MacTel & E1-L & mean A2 border share & 0.029 &  & n=29 \\
framing residual & MacTel & E1-L & share with A2 border share > 0 & 0.276 &  & n=29 \\
framing residual & MacTel & E1-L & share with near-black edge run > 5\% of side & 0.310 &  & n=29 \\
framing residual & MacTel & E1-L & share with near-black edge run > 10\% of side & 0.276 &  & n=29 \\
framing residual & MacTel & E1-L & mean near-white share & $7.5\times10^{-5}$ &  & n=29 \\
framing residual & MacTel & E1-C & mean A2 border share & 0 &  & n=29 \\
framing residual & MacTel & E1-C & share with A2 border share > 0 & 0 &  & n=29 \\
framing residual & MacTel & E1-C & share with near-black edge run > 5\% of side & 0.241 &  & n=29 \\
framing residual & MacTel & E1-C & share with near-black edge run > 10\% of side & 0.172 &  & n=29 \\
framing residual & MacTel & E1-C & mean near-white share & $4.1\times10^{-4}$ &  & n=29 \\
framing residual & OCTDL & original & mean A2 border share & 0 &  & n=371 \\
framing residual & OCTDL & original & share with A2 border share > 0 & 0 &  & n=371 \\
framing residual & OCTDL & original & share with near-black edge run > 5\% of side & 0 &  & n=371 \\
framing residual & OCTDL & original & share with near-black edge run > 10\% of side & 0 &  & n=371 \\
framing residual & OCTDL & original & mean near-white share & $5.7\times10^{-4}$ &  & n=371 \\
framing residual & OCTDL & E1-M & mean A2 border share & 0 &  & n=371 \\
framing residual & OCTDL & E1-M & share with A2 border share > 0 & 0 &  & n=371 \\
framing residual & OCTDL & E1-M & share with near-black edge run > 5\% of side & 0 &  & n=371 \\
framing residual & OCTDL & E1-M & share with near-black edge run > 10\% of side & 0 &  & n=371 \\
framing residual & OCTDL & E1-M & mean near-white share & $5.8\times10^{-4}$ &  & n=371 \\
framing residual & OCTDL & E1-L & mean A2 border share & $6.1\times10^{-4}$ &  & n=371 \\
framing residual & OCTDL & E1-L & share with A2 border share > 0 & 0.024 &  & n=371 \\
framing residual & OCTDL & E1-L & share with near-black edge run > 5\% of side & 0 &  & n=371 \\
framing residual & OCTDL & E1-L & share with near-black edge run > 10\% of side & 0 &  & n=371 \\
framing residual & OCTDL & E1-L & mean near-white share & $5.4\times10^{-5}$ &  & n=371 \\
framing residual & OCTDL & E1-C & mean A2 border share & 0 &  & n=371 \\
framing residual & OCTDL & E1-C & share with A2 border share > 0 & 0 &  & n=371 \\
framing residual & OCTDL & E1-C & share with near-black edge run > 5\% of side & 0 &  & n=371 \\
framing residual & OCTDL & E1-C & share with near-black edge run > 10\% of side & 0 &  & n=371 \\
framing residual & OCTDL & E1-C & mean near-white share & $5.7\times10^{-4}$ &  & n=371 \\
framing residual & OCTID & original & mean A2 border share & $3.5\times10^{-4}$ &  & n=410 \\
framing residual & OCTID & original & share with A2 border share > 0 & 0.061 &  & n=410 \\
framing residual & OCTID & original & share with near-black edge run > 5\% of side & 0 &  & n=410 \\
framing residual & OCTID & original & share with near-black edge run > 10\% of side & 0 &  & n=410 \\
framing residual & OCTID & original & mean near-white share & 0.003 &  & n=410 \\
framing residual & OCTID & E1-M & mean A2 border share & $3.1\times10^{-4}$ &  & n=410 \\
framing residual & OCTID & E1-M & share with A2 border share > 0 & 0.061 &  & n=410 \\
framing residual & OCTID & E1-M & share with near-black edge run > 5\% of side & 0 &  & n=410 \\
framing residual & OCTID & E1-M & share with near-black edge run > 10\% of side & 0 &  & n=410 \\
framing residual & OCTID & E1-M & mean near-white share & $2.6\times10^{-4}$ &  & n=410 \\
framing residual & OCTID & E1-L & mean A2 border share & $1.6\times10^{-4}$ &  & n=410 \\
framing residual & OCTID & E1-L & share with A2 border share > 0 & 0.022 &  & n=410 \\
framing residual & OCTID & E1-L & share with near-black edge run > 5\% of side & 0 &  & n=410 \\
framing residual & OCTID & E1-L & share with near-black edge run > 10\% of side & 0 &  & n=410 \\
framing residual & OCTID & E1-L & mean near-white share & $3.9\times10^{-6}$ &  & n=410 \\
framing residual & OCTID & E1-C & mean A2 border share & 0 &  & n=410 \\
framing residual & OCTID & E1-C & share with A2 border share > 0 & 0 &  & n=410 \\
framing residual & OCTID & E1-C & share with near-black edge run > 5\% of side & 0 &  & n=410 \\
framing residual & OCTID & E1-C & share with near-black edge run > 10\% of side & 0 &  & n=410 \\
framing residual & OCTID & E1-C & mean near-white share & 0.003 &  & n=410 \\
framing residual & RP\textsuperscript{a} & original & mean A2 border share & 0.301 &  & n=31 \\
framing residual & RP\textsuperscript{a} & original & share with A2 border share > 0 & 0.903 &  & n=31 \\
framing residual & RP\textsuperscript{a} & original & share with near-black edge run > 5\% of side & 0.871 &  & n=31 \\
framing residual & RP\textsuperscript{a} & original & share with near-black edge run > 10\% of side & 0.871 &  & n=31 \\
framing residual & RP\textsuperscript{a} & original & mean near-white share & $2.1\times10^{-4}$ &  & n=31 \\
framing residual & RP\textsuperscript{a} & E1-M & mean A2 border share & 0.024 &  & n=31 \\
framing residual & RP\textsuperscript{a} & E1-M & share with A2 border share > 0 & 0.290 &  & n=31 \\
framing residual & RP\textsuperscript{a} & E1-M & share with near-black edge run > 5\% of side & 0.161 &  & n=31 \\
framing residual & RP\textsuperscript{a} & E1-M & share with near-black edge run > 10\% of side & 0.065 &  & n=31 \\
framing residual & RP\textsuperscript{a} & E1-M & mean near-white share & $2.8\times10^{-4}$ &  & n=31 \\
framing residual & RP\textsuperscript{a} & E1-L & mean A2 border share & 0.043 &  & n=31 \\
framing residual & RP\textsuperscript{a} & E1-L & share with A2 border share > 0 & 0.419 &  & n=31 \\
framing residual & RP\textsuperscript{a} & E1-L & share with near-black edge run > 5\% of side & 0.194 &  & n=31 \\
framing residual & RP\textsuperscript{a} & E1-L & share with near-black edge run > 10\% of side & 0.065 &  & n=31 \\
framing residual & RP\textsuperscript{a} & E1-L & mean near-white share & $1.3\times10^{-4}$ &  & n=31 \\
framing residual & RP\textsuperscript{a} & E1-C & mean A2 border share & 0 &  & n=31 \\
framing residual & RP\textsuperscript{a} & E1-C & share with A2 border share > 0 & 0 &  & n=31 \\
framing residual & RP\textsuperscript{a} & E1-C & share with near-black edge run > 5\% of side & 0.129 &  & n=31 \\
framing residual & RP\textsuperscript{a} & E1-C & share with near-black edge run > 10\% of side & 0.032 &  & n=31 \\
framing residual & RP\textsuperscript{a} & E1-C & mean near-white share & $3.1\times10^{-4}$ &  & n=31 \\
framing residual & Stargardt\textsuperscript{a} & original & mean A2 border share & 0.353 &  & n=23 \\
framing residual & Stargardt\textsuperscript{a} & original & share with A2 border share > 0 & 0.957 &  & n=23 \\
framing residual & Stargardt\textsuperscript{a} & original & share with near-black edge run > 5\% of side & 0.957 &  & n=23 \\
framing residual & Stargardt\textsuperscript{a} & original & share with near-black edge run > 10\% of side & 0.913 &  & n=23 \\
framing residual & Stargardt\textsuperscript{a} & original & mean near-white share & $8.8\times10^{-5}$ &  & n=23 \\
framing residual & Stargardt\textsuperscript{a} & E1-M & mean A2 border share & 0.032 &  & n=23 \\
framing residual & Stargardt\textsuperscript{a} & E1-M & share with A2 border share > 0 & 0.391 &  & n=23 \\
framing residual & Stargardt\textsuperscript{a} & E1-M & share with near-black edge run > 5\% of side & 0.087 &  & n=23 \\
framing residual & Stargardt\textsuperscript{a} & E1-M & share with near-black edge run > 10\% of side & 0 &  & n=23 \\
framing residual & Stargardt\textsuperscript{a} & E1-M & mean near-white share & $1.3\times10^{-4}$ &  & n=23 \\
framing residual & Stargardt\textsuperscript{a} & E1-L & mean A2 border share & 0.042 &  & n=23 \\
framing residual & Stargardt\textsuperscript{a} & E1-L & share with A2 border share > 0 & 0.522 &  & n=23 \\
framing residual & Stargardt\textsuperscript{a} & E1-L & share with near-black edge run > 5\% of side & 0.043 &  & n=23 \\
framing residual & Stargardt\textsuperscript{a} & E1-L & share with near-black edge run > 10\% of side & 0 &  & n=23 \\
framing residual & Stargardt\textsuperscript{a} & E1-L & mean near-white share & $3.1\times10^{-5}$ &  & n=23 \\
framing residual & Stargardt\textsuperscript{a} & E1-C & mean A2 border share & 0 &  & n=23 \\
framing residual & Stargardt\textsuperscript{a} & E1-C & share with A2 border share > 0 & 0 &  & n=23 \\
framing residual & Stargardt\textsuperscript{a} & E1-C & share with near-black edge run > 5\% of side & 0.217 &  & n=23 \\
framing residual & Stargardt\textsuperscript{a} & E1-C & share with near-black edge run > 10\% of side & 0 &  & n=23 \\
framing residual & Stargardt\textsuperscript{a} & E1-C & mean near-white share & $1.3\times10^{-4}$ &  & n=23 \\
laptop re-render agreement &  & E1-M & images compared & 5146 &  & a report, never a stop \\
laptop re-render agreement &  & E1-M & images differing & 0 &  & a report, never a stop \\
laptop re-render agreement &  & E1-M & largest grey difference & 0 &  & a report, never a stop \\
laptop re-render agreement &  & E1-M & originals missing or differing from locked SHA256 & 781 &  & a report, never a stop \\
laptop re-render agreement &  & E1-M & originals not readable on this laptop & 0 &  & a report, never a stop \\
laptop re-render agreement &  & E1-L & images compared & 5146 &  & a report, never a stop \\
laptop re-render agreement &  & E1-L & images differing & 0 &  & a report, never a stop \\
laptop re-render agreement &  & E1-L & largest grey difference & 0 &  & a report, never a stop \\
laptop re-render agreement &  & E1-L & originals missing or differing from locked SHA256 & 781 &  & a report, never a stop \\
laptop re-render agreement &  & E1-L & originals not readable on this laptop & 0 &  & a report, never a stop \\
laptop re-render agreement &  & E1-C & images compared & 5146 &  & a report, never a stop \\
laptop re-render agreement &  & E1-C & images differing & 0 &  & a report, never a stop \\
laptop re-render agreement &  & E1-C & largest grey difference & 0 &  & a report, never a stop \\
laptop re-render agreement &  & E1-C & originals missing or differing from locked SHA256 & 781 &  & a report, never a stop \\
laptop re-render agreement &  & E1-C & originals not readable on this laptop & 0 &  & a report, never a stop \\
\end{longtable}
}
Hashes are shortened to 12 characters; the full values are in the CSV.

\fi
\havetable{e1b_development}\ifgen
\subsubsection*{E1b. Development macro recall per model and feature set, references beside}
Mean over 100 draws (SD); mean over $K$ 1--8: mean (SD; 5th--95th percentile). Spread over support draws, not a confidence interval.
{\scriptsize\setlength{\tabcolsep}{3pt}% src: E1b_development.csv in C:/Users/ThinkBook/.t3/worktrees/Few-shot/t3-03f3a2c4/results/fm_v1/posthoc (generated by make_tables.py; formatting only);
% label: post hoc, Phase-1 rule E1 timestamped by commit 4b4e883 (20 Nov 2026), computed 2026-11-21; outside the locked protocol; never replaces a locked result; limitations: the centre square drops lateral retina unequally by source and class; MIRAGE now gets 256 px of detail enlarged to 512; E1-M does not remove effective resolution (content short side 103-496 px), field of view, scanner appearance, compression, intensity, black-clipped background against speckle, or Kermany's white fill wedges
\begin{longtable}{@{}p{3.6cm}lcccccp{2.6cm}@{}}
\toprule
Model & Set & $K=1$ & $K=2$ & $K=4$ & $K=8$ & Mean $K$ 1--8 & Status \\
\midrule
\endfirsthead
\toprule
Model & Set & $K=1$ & $K=2$ & $K=4$ & $K=8$ & Mean $K$ 1--8 & Status \\
\midrule
\endhead
\bottomrule
\endfoot
\bottomrule
\endlastfoot
MIRAGE & locked & \makecell{0.495\\{\scriptsize (0.046)}} & \makecell{0.615\\{\scriptsize (0.048)}} & \makecell{0.712\\{\scriptsize (0.040)}} & \makecell{0.777\\{\scriptsize (0.030)}} & \makecell{0.650\\{\scriptsize (0.028; 0.603--0.688)}} & interpreted \\
 & E1-M & \makecell{0.486\\{\scriptsize (0.045)}} & \makecell{0.614\\{\scriptsize (0.046)}} & \makecell{0.715\\{\scriptsize (0.035)}} & \makecell{0.784\\{\scriptsize (0.025)}} & \makecell{0.649\\{\scriptsize (0.024; 0.611--0.692)}} & interpreted \\
 & E1-L & \makecell{0.494\\{\scriptsize (0.046)}} & \makecell{0.620\\{\scriptsize (0.045)}} & \makecell{0.719\\{\scriptsize (0.037)}} & \makecell{0.786\\{\scriptsize (0.026)}} & \makecell{0.655\\{\scriptsize (0.024; 0.618--0.694)}} & interpreted \\
 & E1-C & \makecell{0.471\\{\scriptsize (0.045)}} & \makecell{0.585\\{\scriptsize (0.044)}} & \makecell{0.684\\{\scriptsize (0.031)}} & \makecell{0.750\\{\scriptsize (0.028)}} & \makecell{0.622\\{\scriptsize (0.024; 0.585--0.660)}} & interpreted \\
\addlinespace[2pt]
DINOv2 & locked & \makecell{0.407\\{\scriptsize (0.046)}} & \makecell{0.497\\{\scriptsize (0.043)}} & \makecell{0.616\\{\scriptsize (0.043)}} & \makecell{0.711\\{\scriptsize (0.031)}} & \makecell{0.558\\{\scriptsize (0.028; 0.511--0.606)}} & interpreted \\
 & E1-M & \makecell{0.403\\{\scriptsize (0.043)}} & \makecell{0.501\\{\scriptsize (0.045)}} & \makecell{0.597\\{\scriptsize (0.045)}} & \makecell{0.668\\{\scriptsize (0.035)}} & \makecell{0.542\\{\scriptsize (0.028; 0.499--0.587)}} & interpreted \\
 & E1-L & \makecell{0.431\\{\scriptsize (0.042)}} & \makecell{0.527\\{\scriptsize (0.048)}} & \makecell{0.632\\{\scriptsize (0.043)}} & \makecell{0.713\\{\scriptsize (0.035)}} & \makecell{0.576\\{\scriptsize (0.030; 0.529--0.622)}} & interpreted \\
 & E1-C & \makecell{0.399\\{\scriptsize (0.044)}} & \makecell{0.501\\{\scriptsize (0.045)}} & \makecell{0.597\\{\scriptsize (0.045)}} & \makecell{0.663\\{\scriptsize (0.037)}} & \makecell{0.540\\{\scriptsize (0.028; 0.498--0.588)}} & interpreted \\
\addlinespace[2pt]
Inception-v3 & locked & \makecell{0.409\\{\scriptsize (0.046)}} & \makecell{0.499\\{\scriptsize (0.043)}} & \makecell{0.581\\{\scriptsize (0.039)}} & \makecell{0.647\\{\scriptsize (0.033)}} & \makecell{0.534\\{\scriptsize (0.029; 0.494--0.583)}} & interpreted \\
 & E1-M & \makecell{0.454\\{\scriptsize (0.049)}} & \makecell{0.547\\{\scriptsize (0.042)}} & \makecell{0.635\\{\scriptsize (0.038)}} & \makecell{0.709\\{\scriptsize (0.031)}} & \makecell{0.586\\{\scriptsize (0.027; 0.540--0.635)}} & interpreted \\
 & E1-L & \makecell{0.441\\{\scriptsize (0.048)}} & \makecell{0.536\\{\scriptsize (0.041)}} & \makecell{0.630\\{\scriptsize (0.037)}} & \makecell{0.705\\{\scriptsize (0.033)}} & \makecell{0.578\\{\scriptsize (0.025; 0.537--0.614)}} & interpreted \\
 & E1-C & \makecell{0.423\\{\scriptsize (0.051)}} & \makecell{0.514\\{\scriptsize (0.042)}} & \makecell{0.606\\{\scriptsize (0.039)}} & \makecell{0.675\\{\scriptsize (0.035)}} & \makecell{0.555\\{\scriptsize (0.031; 0.510--0.608)}} & interpreted \\
\addlinespace[2pt]
RETFound & locked & \makecell{0.318\\{\scriptsize (0.040)}} & \makecell{0.375\\{\scriptsize (0.041)}} & \makecell{0.442\\{\scriptsize (0.039)}} & \makecell{0.489\\{\scriptsize (0.031)}} & \makecell{0.406\\{\scriptsize (0.024; 0.367--0.441)}} & interpreted \\
 & E1-M & \makecell{0.317\\{\scriptsize (0.042)}} & \makecell{0.367\\{\scriptsize (0.046)}} & \makecell{0.446\\{\scriptsize (0.048)}} & \makecell{0.517\\{\scriptsize (0.042)}} & \makecell{0.412\\{\scriptsize (0.029; 0.364--0.458)}} & interpreted \\
 & E1-L & \makecell{0.340\\{\scriptsize (0.042)}} & \makecell{0.411\\{\scriptsize (0.048)}} & \makecell{0.483\\{\scriptsize (0.048)}} & \makecell{0.565\\{\scriptsize (0.044)}} & \makecell{0.450\\{\scriptsize (0.032; 0.400--0.501)}} & interpreted \\
 & E1-C & \makecell{0.316\\{\scriptsize (0.042)}} & \makecell{0.365\\{\scriptsize (0.046)}} & \makecell{0.447\\{\scriptsize (0.049)}} & \makecell{0.519\\{\scriptsize (0.042)}} & \makecell{0.412\\{\scriptsize (0.029; 0.363--0.456)}} & interpreted \\
\addlinespace[2pt]
Size-only & raw & \makecell{0.578\\{\scriptsize (0.058)}} & \makecell{0.550\\{\scriptsize (0.052)}} & \makecell{0.560\\{\scriptsize (0.051)}} & \makecell{0.589\\{\scriptsize (0.049)}} & \makecell{0.569\\{\scriptsize (0.036; 0.508--0.619)}} & reference without a model \\
\addlinespace[2pt]
Content-size & raw content box & \makecell{0.404\\{\scriptsize (0.058)}} & \makecell{0.432\\{\scriptsize (0.049)}} & \makecell{0.463\\{\scriptsize (0.037)}} & \makecell{0.490\\{\scriptsize (0.026)}} & \makecell{0.447\\{\scriptsize (0.028; 0.395--0.495)}} & reference without a model \\
\addlinespace[2pt]
Thumbnail & E1-M & \makecell{0.305\\{\scriptsize (0.041)}} & \makecell{0.335\\{\scriptsize (0.041)}} & \makecell{0.381\\{\scriptsize (0.042)}} & \makecell{0.419\\{\scriptsize (0.037)}} & \makecell{0.360\\{\scriptsize (0.025; 0.325--0.399)}} & reference without a model \\
 & E1-L & \makecell{0.305\\{\scriptsize (0.040)}} & \makecell{0.336\\{\scriptsize (0.042)}} & \makecell{0.382\\{\scriptsize (0.042)}} & \makecell{0.419\\{\scriptsize (0.037)}} & \makecell{0.360\\{\scriptsize (0.025; 0.325--0.399)}} & reference without a model \\
 & E1-C & \makecell{0.273\\{\scriptsize (0.035)}} & \makecell{0.306\\{\scriptsize (0.039)}} & \makecell{0.349\\{\scriptsize (0.041)}} & \makecell{0.393\\{\scriptsize (0.036)}} & \makecell{0.331\\{\scriptsize (0.024; 0.292--0.364)}} & reference without a model \\
 & raw & \makecell{0.256\\{\scriptsize (0.038)}} & \makecell{0.272\\{\scriptsize (0.035)}} & \makecell{0.300\\{\scriptsize (0.042)}} & \makecell{0.325\\{\scriptsize (0.038)}} & \makecell{0.288\\{\scriptsize (0.025; 0.246--0.325)}} & reference without a model \\
\end{longtable}
}
\fi
\havetable{e1c_paired}\ifgen
\subsubsection*{E1c. Paired differences}
Each foundation model minus Inception-v3 within a feature set, and each feature set minus the same model's locked features, on the same draws: mean (5th to 95th percentile).
{\tiny\let\scriptsize\tiny\setlength{\tabcolsep}{2pt}% src: E1c_paired.csv in C:/Users/ThinkBook/.t3/worktrees/Few-shot/t3-03f3a2c4/results/fm_v1/posthoc (generated by make_tables.py; formatting only);
% label: post hoc, Phase-1 rule E1 timestamped by commit 4b4e883 (20 Nov 2026), computed 2026-11-21; outside the locked protocol; never replaces a locked result; limitations: the centre square drops lateral retina unequally by source and class; MIRAGE now gets 256 px of detail enlarged to 512; E1-M does not remove effective resolution (content short side 103-496 px), field of view, scanner appearance, compression, intensity, black-clipped background against speckle, or Kermany's white fill wedges
\begin{longtable}{@{}p{3.6cm}lcccccp{2.6cm}@{}}
\toprule
Comparison & Set & $K=1$ & $K=2$ & $K=4$ & $K=8$ & Mean $K$ 1--8 & Status \\
\midrule
\endfirsthead
\toprule
Comparison & Set & $K=1$ & $K=2$ & $K=4$ & $K=8$ & Mean $K$ 1--8 & Status \\
\midrule
\endhead
\bottomrule
\endfoot
\bottomrule
\endlastfoot
MIRAGE E1-C $-$ MIRAGE locked & E1-C & \makecell{$-$0.024\\{\scriptsize ($-$0.103 to +0.072)}} & \makecell{$-$0.029\\{\scriptsize ($-$0.122 to +0.037)}} & \makecell{$-$0.028\\{\scriptsize ($-$0.098 to +0.041)}} & \makecell{$-$0.027\\{\scriptsize ($-$0.095 to +0.025)}} & \makecell{$-$0.027\\{\scriptsize (0.030; $-$0.077 to +0.021)}} & interpreted \\
\addlinespace[2pt]
MIRAGE E1-L $-$ MIRAGE locked & E1-L & \makecell{$-$0.001\\{\scriptsize ($-$0.094 to +0.077)}} & \makecell{+0.005\\{\scriptsize ($-$0.101 to +0.086)}} & \makecell{+0.007\\{\scriptsize ($-$0.056 to +0.080)}} & \makecell{+0.009\\{\scriptsize ($-$0.041 to +0.062)}} & \makecell{+0.005\\{\scriptsize (0.030; $-$0.046 to +0.048)}} & interpreted \\
\addlinespace[2pt]
MIRAGE E1-M $-$ MIRAGE locked & E1-M & \makecell{$-$0.009\\{\scriptsize ($-$0.096 to +0.077)}} & \makecell{$-$0.001\\{\scriptsize ($-$0.105 to +0.091)}} & \makecell{+0.002\\{\scriptsize ($-$0.053 to +0.075)}} & \makecell{+0.007\\{\scriptsize ($-$0.033 to +0.053)}} & \makecell{$-$0.000\\{\scriptsize (0.029; $-$0.051 to +0.043)}} & interpreted \\
\addlinespace[2pt]
MIRAGE $-$ Inception-v3 & locked & \makecell{+0.086\\{\scriptsize (+0.001 to +0.184)}} & \makecell{+0.116\\{\scriptsize (+0.015 to +0.230)}} & \makecell{+0.132\\{\scriptsize (+0.052 to +0.222)}} & \makecell{+0.130\\{\scriptsize (+0.071 to +0.210)}} & \makecell{+0.116\\{\scriptsize (0.037; +0.069 to +0.185)}} & interpreted \\
 & E1-M & \makecell{+0.031\\{\scriptsize ($-$0.068 to +0.122)}} & \makecell{+0.067\\{\scriptsize ($-$0.034 to +0.160)}} & \makecell{+0.080\\{\scriptsize (+0.015 to +0.168)}} & \makecell{+0.075\\{\scriptsize (+0.016 to +0.142)}} & \makecell{+0.063\\{\scriptsize (0.036; +0.003 to +0.117)}} & interpreted \\
 & E1-L & \makecell{+0.053\\{\scriptsize ($-$0.041 to +0.139)}} & \makecell{+0.084\\{\scriptsize (+0.005 to +0.162)}} & \makecell{+0.090\\{\scriptsize (+0.012 to +0.163)}} & \makecell{+0.081\\{\scriptsize (+0.008 to +0.152)}} & \makecell{+0.077\\{\scriptsize (0.034; +0.023 to +0.132)}} & interpreted \\
 & E1-C & \makecell{+0.048\\{\scriptsize ($-$0.049 to +0.154)}} & \makecell{+0.072\\{\scriptsize ($-$0.013 to +0.156)}} & \makecell{+0.078\\{\scriptsize (+0.002 to +0.146)}} & \makecell{+0.074\\{\scriptsize (+0.021 to +0.152)}} & \makecell{+0.068\\{\scriptsize (0.033; +0.017 to +0.123)}} & interpreted \\
\addlinespace[2pt]
DINOv2 E1-C $-$ DINOv2 locked & E1-C & \makecell{$-$0.008\\{\scriptsize ($-$0.068 to +0.081)}} & \makecell{+0.005\\{\scriptsize ($-$0.077 to +0.095)}} & \makecell{$-$0.019\\{\scriptsize ($-$0.096 to +0.060)}} & \makecell{$-$0.048\\{\scriptsize ($-$0.098 to +0.011)}} & \makecell{$-$0.018\\{\scriptsize (0.030; $-$0.066 to +0.040)}} & interpreted \\
\addlinespace[2pt]
DINOv2 E1-L $-$ DINOv2 locked & E1-L & \makecell{+0.024\\{\scriptsize ($-$0.044 to +0.109)}} & \makecell{+0.030\\{\scriptsize ($-$0.054 to +0.133)}} & \makecell{+0.015\\{\scriptsize ($-$0.052 to +0.092)}} & \makecell{+0.002\\{\scriptsize ($-$0.050 to +0.059)}} & \makecell{+0.018\\{\scriptsize (0.030; $-$0.025 to +0.064)}} & interpreted \\
\addlinespace[2pt]
DINOv2 E1-M $-$ DINOv2 locked & E1-M & \makecell{$-$0.003\\{\scriptsize ($-$0.074 to +0.088)}} & \makecell{+0.004\\{\scriptsize ($-$0.075 to +0.088)}} & \makecell{$-$0.019\\{\scriptsize ($-$0.097 to +0.064)}} & \makecell{$-$0.043\\{\scriptsize ($-$0.087 to +0.017)}} & \makecell{$-$0.015\\{\scriptsize (0.030; $-$0.058 to +0.034)}} & interpreted \\
\addlinespace[2pt]
DINOv2 $-$ Inception-v3 & locked & \makecell{$-$0.002\\{\scriptsize ($-$0.089 to +0.077)}} & \makecell{$-$0.002\\{\scriptsize ($-$0.094 to +0.087)}} & \makecell{+0.036\\{\scriptsize ($-$0.046 to +0.120)}} & \makecell{+0.064\\{\scriptsize ($-$0.008 to +0.136)}} & \makecell{+0.024\\{\scriptsize (0.034; $-$0.025 to +0.086)}} & interpreted \\
 & E1-M & \makecell{$-$0.051\\{\scriptsize ($-$0.124 to +0.024)}} & \makecell{$-$0.046\\{\scriptsize ($-$0.143 to +0.033)}} & \makecell{$-$0.038\\{\scriptsize ($-$0.124 to +0.036)}} & \makecell{$-$0.041\\{\scriptsize ($-$0.110 to +0.029)}} & \makecell{$-$0.044\\{\scriptsize (0.032; $-$0.099 to +0.004)}} & interpreted \\
 & E1-L & \makecell{$-$0.010\\{\scriptsize ($-$0.093 to +0.077)}} & \makecell{$-$0.009\\{\scriptsize ($-$0.096 to +0.086)}} & \makecell{+0.002\\{\scriptsize ($-$0.085 to +0.081)}} & \makecell{+0.008\\{\scriptsize ($-$0.061 to +0.079)}} & \makecell{$-$0.002\\{\scriptsize (0.035; $-$0.054 to +0.051)}} & interpreted \\
 & E1-C & \makecell{$-$0.025\\{\scriptsize ($-$0.130 to +0.064)}} & \makecell{$-$0.012\\{\scriptsize ($-$0.105 to +0.068)}} & \makecell{$-$0.009\\{\scriptsize ($-$0.099 to +0.079)}} & \makecell{$-$0.012\\{\scriptsize ($-$0.085 to +0.068)}} & \makecell{$-$0.014\\{\scriptsize (0.038; $-$0.072 to +0.052)}} & interpreted \\
\addlinespace[2pt]
Inception-v3 E1-C $-$ Inception-v3 locked & E1-C & \makecell{+0.015\\{\scriptsize ($-$0.061 to +0.113)}} & \makecell{+0.014\\{\scriptsize ($-$0.057 to +0.083)}} & \makecell{+0.025\\{\scriptsize ($-$0.032 to +0.089)}} & \makecell{+0.028\\{\scriptsize ($-$0.023 to +0.081)}} & \makecell{+0.021\\{\scriptsize (0.024; $-$0.018 to +0.057)}} & interpreted \\
\addlinespace[2pt]
Inception-v3 E1-L $-$ Inception-v3 locked & E1-L & \makecell{+0.032\\{\scriptsize ($-$0.064 to +0.135)}} & \makecell{+0.037\\{\scriptsize ($-$0.045 to +0.105)}} & \makecell{+0.049\\{\scriptsize ($-$0.016 to +0.107)}} & \makecell{+0.058\\{\scriptsize ($-$0.003 to +0.126)}} & \makecell{+0.044\\{\scriptsize (0.028; $-$0.004 to +0.093)}} & interpreted \\
\addlinespace[2pt]
Inception-v3 E1-M $-$ Inception-v3 locked & E1-M & \makecell{+0.046\\{\scriptsize ($-$0.027 to +0.143)}} & \makecell{+0.048\\{\scriptsize ($-$0.033 to +0.122)}} & \makecell{+0.055\\{\scriptsize ($-$0.018 to +0.127)}} & \makecell{+0.062\\{\scriptsize (+0.004 to +0.130)}} & \makecell{+0.052\\{\scriptsize (0.028; +0.006 to +0.093)}} & interpreted \\
\addlinespace[2pt]
RETFound E1-C $-$ RETFound locked & E1-C & \makecell{$-$0.002\\{\scriptsize ($-$0.067 to +0.067)}} & \makecell{$-$0.010\\{\scriptsize ($-$0.074 to +0.060)}} & \makecell{+0.005\\{\scriptsize ($-$0.061 to +0.081)}} & \makecell{+0.030\\{\scriptsize ($-$0.043 to +0.100)}} & \makecell{+0.006\\{\scriptsize (0.028; $-$0.033 to +0.052)}} & interpreted \\
\addlinespace[2pt]
RETFound E1-L $-$ RETFound locked & E1-L & \makecell{+0.021\\{\scriptsize ($-$0.060 to +0.096)}} & \makecell{+0.036\\{\scriptsize ($-$0.039 to +0.121)}} & \makecell{+0.041\\{\scriptsize ($-$0.022 to +0.113)}} & \makecell{+0.076\\{\scriptsize (+0.004 to +0.151)}} & \makecell{+0.044\\{\scriptsize (0.031; $-$0.000 to +0.100)}} & interpreted \\
\addlinespace[2pt]
RETFound E1-M $-$ RETFound locked & E1-M & \makecell{$-$0.002\\{\scriptsize ($-$0.068 to +0.063)}} & \makecell{$-$0.009\\{\scriptsize ($-$0.068 to +0.059)}} & \makecell{+0.004\\{\scriptsize ($-$0.060 to +0.081)}} & \makecell{+0.028\\{\scriptsize ($-$0.039 to +0.106)}} & \makecell{+0.005\\{\scriptsize (0.027; $-$0.037 to +0.051)}} & interpreted \\
\addlinespace[2pt]
RETFound $-$ Inception-v3 & locked & \makecell{$-$0.090\\{\scriptsize ($-$0.182 to $-$0.000)}} & \makecell{$-$0.124\\{\scriptsize ($-$0.231 to $-$0.029)}} & \makecell{$-$0.139\\{\scriptsize ($-$0.217 to $-$0.060)}} & \makecell{$-$0.159\\{\scriptsize ($-$0.232 to $-$0.086)}} & \makecell{$-$0.128\\{\scriptsize (0.035; $-$0.186 to $-$0.070)}} & interpreted \\
 & E1-M & \makecell{$-$0.138\\{\scriptsize ($-$0.245 to $-$0.044)}} & \makecell{$-$0.180\\{\scriptsize ($-$0.274 to $-$0.079)}} & \makecell{$-$0.189\\{\scriptsize ($-$0.269 to $-$0.106)}} & \makecell{$-$0.192\\{\scriptsize ($-$0.271 to $-$0.117)}} & \makecell{$-$0.175\\{\scriptsize (0.035; $-$0.236 to $-$0.124)}} & interpreted \\
 & E1-L & \makecell{$-$0.101\\{\scriptsize ($-$0.188 to $-$0.014)}} & \makecell{$-$0.125\\{\scriptsize ($-$0.217 to $-$0.043)}} & \makecell{$-$0.147\\{\scriptsize ($-$0.238 to $-$0.069)}} & \makecell{$-$0.140\\{\scriptsize ($-$0.211 to $-$0.065)}} & \makecell{$-$0.128\\{\scriptsize (0.037; $-$0.187 to $-$0.071)}} & interpreted \\
 & E1-C & \makecell{$-$0.107\\{\scriptsize ($-$0.200 to $-$0.014)}} & \makecell{$-$0.148\\{\scriptsize ($-$0.258 to $-$0.057)}} & \makecell{$-$0.159\\{\scriptsize ($-$0.252 to $-$0.077)}} & \makecell{$-$0.157\\{\scriptsize ($-$0.237 to $-$0.078)}} & \makecell{$-$0.143\\{\scriptsize (0.037; $-$0.201 to $-$0.085)}} & interpreted \\
\end{longtable}
}
\fi
\havetable{e1e_matched}\ifgen
\subsubsection*{E1e. Matched subset}
Matched subset: L=164 px, U=296 px, 92 development queries; supports are not matched.

{\tiny\let\scriptsize\tiny\setlength{\tabcolsep}{2pt}% src: E1e_matched.csv in C:/Users/ThinkBook/.t3/worktrees/Few-shot/t3-03f3a2c4/results/fm_v1/posthoc (generated by make_tables.py; formatting only);
% label: post hoc, Phase-1 rule E1 timestamped by commit 4b4e883 (20 Nov 2026), computed 2026-11-21; outside the locked protocol; never replaces a locked result; limitations: the centre square drops lateral retina unequally by source and class; MIRAGE now gets 256 px of detail enlarged to 512; E1-M does not remove effective resolution (content short side 103-496 px), field of view, scanner appearance, compression, intensity, black-clipped background against speckle, or Kermany's white fill wedges
\begin{longtable}{@{}p{3.2cm}lcccccp{2.4cm}p{2.2cm}@{}}
\toprule
Model & Set & $K=1$ & $K=2$ & $K=4$ & $K=8$ & Mean $K$ 1--8 & Classes present & Status \\
\midrule
\endfirsthead
\toprule
Model & Set & $K=1$ & $K=2$ & $K=4$ & $K=8$ & Mean $K$ 1--8 & Classes present & Status \\
\midrule
\endhead
\bottomrule
\endfoot
\bottomrule
\endlastfoot
MIRAGE & locked & \makecell{0.490\\{\scriptsize (0.054)}} & \makecell{0.614\\{\scriptsize (0.054)}} & \makecell{0.718\\{\scriptsize (0.046)}} & \makecell{0.791\\{\scriptsize (0.035)}} & \makecell{0.653\\{\scriptsize (0.033; 0.599--0.700)}} & CSC MH MacTel Stargardt RP & interpreted \\
 & E1-M & \makecell{0.487\\{\scriptsize (0.051)}} & \makecell{0.621\\{\scriptsize (0.053)}} & \makecell{0.727\\{\scriptsize (0.044)}} & \makecell{0.793\\{\scriptsize (0.030)}} & \makecell{0.657\\{\scriptsize (0.030; 0.606--0.696)}} & CSC MH MacTel Stargardt RP & interpreted \\
 & E1-L & \makecell{0.494\\{\scriptsize (0.054)}} & \makecell{0.624\\{\scriptsize (0.052)}} & \makecell{0.724\\{\scriptsize (0.043)}} & \makecell{0.787\\{\scriptsize (0.033)}} & \makecell{0.657\\{\scriptsize (0.030; 0.616--0.701)}} & CSC MH MacTel Stargardt RP & interpreted \\
 & E1-C & \makecell{0.482\\{\scriptsize (0.056)}} & \makecell{0.604\\{\scriptsize (0.051)}} & \makecell{0.713\\{\scriptsize (0.039)}} & \makecell{0.781\\{\scriptsize (0.031)}} & \makecell{0.645\\{\scriptsize (0.030; 0.594--0.690)}} & CSC MH MacTel Stargardt RP & interpreted \\
\addlinespace[2pt]
MIRAGE $-$ Inception-v3 & locked & \makecell{+0.081\\{\scriptsize ($-$0.022 to +0.174)}} & \makecell{+0.118\\{\scriptsize (+0.008 to +0.240)}} & \makecell{+0.148\\{\scriptsize (+0.058 to +0.244)}} & \makecell{+0.161\\{\scriptsize (+0.089 to +0.245)}} & \makecell{+0.127\\{\scriptsize (0.041; +0.067 to +0.196)}} & -- & interpreted \\
 & E1-M & \makecell{+0.024\\{\scriptsize ($-$0.104 to +0.129)}} & \makecell{+0.064\\{\scriptsize ($-$0.048 to +0.180)}} & \makecell{+0.072\\{\scriptsize ($-$0.016 to +0.151)}} & \makecell{+0.066\\{\scriptsize ($-$0.003 to +0.128)}} & \makecell{+0.057\\{\scriptsize (0.039; $-$0.012 to +0.120)}} & -- & interpreted \\
 & E1-L & \makecell{+0.044\\{\scriptsize ($-$0.059 to +0.143)}} & \makecell{+0.080\\{\scriptsize ($-$0.019 to +0.177)}} & \makecell{+0.081\\{\scriptsize ($-$0.010 to +0.166)}} & \makecell{+0.067\\{\scriptsize ($-$0.018 to +0.140)}} & \makecell{+0.068\\{\scriptsize (0.039; +0.010 to +0.134)}} & -- & interpreted \\
 & E1-C & \makecell{+0.062\\{\scriptsize ($-$0.047 to +0.166)}} & \makecell{+0.089\\{\scriptsize (+0.001 to +0.179)}} & \makecell{+0.097\\{\scriptsize (+0.021 to +0.187)}} & \makecell{+0.090\\{\scriptsize (+0.028 to +0.171)}} & \makecell{+0.084\\{\scriptsize (0.037; +0.027 to +0.140)}} & -- & interpreted \\
\addlinespace[2pt]
DINOv2 & locked & \makecell{0.405\\{\scriptsize (0.057)}} & \makecell{0.490\\{\scriptsize (0.055)}} & \makecell{0.613\\{\scriptsize (0.055)}} & \makecell{0.717\\{\scriptsize (0.035)}} & \makecell{0.556\\{\scriptsize (0.036; 0.511--0.620)}} & CSC MH MacTel Stargardt RP & interpreted \\
 & E1-M & \makecell{0.394\\{\scriptsize (0.054)}} & \makecell{0.490\\{\scriptsize (0.055)}} & \makecell{0.586\\{\scriptsize (0.053)}} & \makecell{0.665\\{\scriptsize (0.040)}} & \makecell{0.534\\{\scriptsize (0.034; 0.478--0.589)}} & CSC MH MacTel Stargardt RP & interpreted \\
 & E1-L & \makecell{0.422\\{\scriptsize (0.052)}} & \makecell{0.520\\{\scriptsize (0.057)}} & \makecell{0.629\\{\scriptsize (0.049)}} & \makecell{0.719\\{\scriptsize (0.036)}} & \makecell{0.572\\{\scriptsize (0.034; 0.517--0.632)}} & CSC MH MacTel Stargardt RP & interpreted \\
 & E1-C & \makecell{0.390\\{\scriptsize (0.054)}} & \makecell{0.488\\{\scriptsize (0.056)}} & \makecell{0.585\\{\scriptsize (0.053)}} & \makecell{0.657\\{\scriptsize (0.041)}} & \makecell{0.530\\{\scriptsize (0.035; 0.477--0.580)}} & CSC MH MacTel Stargardt RP & interpreted \\
\addlinespace[2pt]
Inception-v3 & locked & \makecell{0.409\\{\scriptsize (0.053)}} & \makecell{0.496\\{\scriptsize (0.050)}} & \makecell{0.570\\{\scriptsize (0.041)}} & \makecell{0.630\\{\scriptsize (0.035)}} & \makecell{0.526\\{\scriptsize (0.032; 0.478--0.579)}} & CSC MH MacTel Stargardt RP & interpreted \\
 & E1-M & \makecell{0.463\\{\scriptsize (0.057)}} & \makecell{0.557\\{\scriptsize (0.049)}} & \makecell{0.655\\{\scriptsize (0.043)}} & \makecell{0.727\\{\scriptsize (0.029)}} & \makecell{0.600\\{\scriptsize (0.029; 0.550--0.641)}} & CSC MH MacTel Stargardt RP & interpreted \\
 & E1-L & \makecell{0.450\\{\scriptsize (0.057)}} & \makecell{0.544\\{\scriptsize (0.048)}} & \makecell{0.644\\{\scriptsize (0.043)}} & \makecell{0.720\\{\scriptsize (0.033)}} & \makecell{0.589\\{\scriptsize (0.029; 0.538--0.633)}} & CSC MH MacTel Stargardt RP & interpreted \\
 & E1-C & \makecell{0.420\\{\scriptsize (0.054)}} & \makecell{0.515\\{\scriptsize (0.047)}} & \makecell{0.616\\{\scriptsize (0.045)}} & \makecell{0.691\\{\scriptsize (0.036)}} & \makecell{0.561\\{\scriptsize (0.033; 0.514--0.617)}} & CSC MH MacTel Stargardt RP & interpreted \\
\addlinespace[2pt]
RETFound & locked & \makecell{0.317\\{\scriptsize (0.049)}} & \makecell{0.377\\{\scriptsize (0.048)}} & \makecell{0.446\\{\scriptsize (0.047)}} & \makecell{0.493\\{\scriptsize (0.041)}} & \makecell{0.408\\{\scriptsize (0.029; 0.361--0.453)}} & CSC MH MacTel Stargardt RP & interpreted \\
 & E1-M & \makecell{0.315\\{\scriptsize (0.048)}} & \makecell{0.370\\{\scriptsize (0.049)}} & \makecell{0.457\\{\scriptsize (0.053)}} & \makecell{0.524\\{\scriptsize (0.045)}} & \makecell{0.417\\{\scriptsize (0.031; 0.364--0.469)}} & CSC MH MacTel Stargardt RP & interpreted \\
 & E1-L & \makecell{0.341\\{\scriptsize (0.051)}} & \makecell{0.411\\{\scriptsize (0.050)}} & \makecell{0.490\\{\scriptsize (0.054)}} & \makecell{0.567\\{\scriptsize (0.052)}} & \makecell{0.452\\{\scriptsize (0.035; 0.400--0.514)}} & CSC MH MacTel Stargardt RP & interpreted \\
 & E1-C & \makecell{0.314\\{\scriptsize (0.048)}} & \makecell{0.369\\{\scriptsize (0.048)}} & \makecell{0.459\\{\scriptsize (0.053)}} & \makecell{0.525\\{\scriptsize (0.045)}} & \makecell{0.417\\{\scriptsize (0.031; 0.358--0.458)}} & CSC MH MacTel Stargardt RP & interpreted \\
\end{longtable}
}
\fi
\havetable{e1f_holdout}\ifgen
\subsubsection*{E1f. Hold-out of 22}
The hold-out is read for direction only: the sign of the mean paired MIRAGE $-$ Inception-v3 difference, per $K$ and in the mean over $K=1$--8. $K=16$ is shown, never read.

{\tiny\let\scriptsize\tiny\setlength{\tabcolsep}{2pt}% src: E1f_holdout.csv in results/fm_v1/posthoc (generated by make_tables.py; formatting only);
% label: post hoc, Phase-1 rule E1 timestamped by commit 4b4e883 (20 Nov 2026), computed 2026-11-21; outside the locked protocol; never replaces a locked result; limitations: the centre square drops lateral retina unequally by source and class; MIRAGE now gets 256 px of detail enlarged to 512; E1-M does not remove effective resolution (content short side 103-496 px), field of view, scanner appearance, compression, intensity, black-clipped background against speckle, or Kermany's white fill wedges
\begin{longtable}{@{}p{3.0cm}lcccccclp{2.2cm}@{}}
\toprule
Model & Set & $K=1$ & $K=2$ & $K=4$ & $K=8$ & \makecell{$K=16$\\(not read)} & Mean $K$ 1--8 & Direction & Status \\
\midrule
\endfirsthead
\toprule
Model & Set & $K=1$ & $K=2$ & $K=4$ & $K=8$ & \makecell{$K=16$\\(not read)} & Mean $K$ 1--8 & Direction & Status \\
\midrule
\endhead
\bottomrule
\endfoot
\bottomrule
\endlastfoot
\multicolumn{10}{@{}l}{\emph{Five-way macro recall}} \\
MIRAGE & locked & 0.536 & 0.685 & 0.799 & 0.867 & 0.894 & \makecell{0.722\\{\scriptsize (0.063; 0.615--0.830)}} &  & interpreted \\
 & E1-M & 0.528 & 0.633 & 0.713 & 0.797 & 0.853 & \makecell{0.667\\{\scriptsize (0.071; 0.532--0.763)}} &  & interpreted \\
 & E1-L & 0.556 & 0.645 & 0.745 & 0.823 & 0.855 & \makecell{0.692\\{\scriptsize (0.065; 0.562--0.788)}} &  & interpreted \\
 & E1-C & 0.527 & 0.630 & 0.732 & 0.830 & 0.889 & \makecell{0.680\\{\scriptsize (0.071; 0.537--0.788)}} &  & interpreted \\
DINOv2 & locked & 0.395 & 0.477 & 0.577 & 0.657 & 0.680 & \makecell{0.526\\{\scriptsize (0.065; 0.410--0.638)}} &  & interpreted \\
 & E1-M & 0.422 & 0.525 & 0.607 & 0.657 & 0.661 & \makecell{0.553\\{\scriptsize (0.061; 0.455--0.653)}} &  & interpreted \\
 & E1-L & 0.415 & 0.505 & 0.575 & 0.607 & 0.620 & \makecell{0.526\\{\scriptsize (0.062; 0.438--0.640)}} &  & interpreted \\
 & E1-C & 0.412 & 0.519 & 0.603 & 0.649 & 0.645 & \makecell{0.546\\{\scriptsize (0.062; 0.447--0.628)}} &  & interpreted \\
Inception-v3 & locked & 0.430 & 0.497 & 0.556 & 0.614 & 0.621 & \makecell{0.524\\{\scriptsize (0.067; 0.427--0.628)}} &  & interpreted \\
 & E1-M & 0.442 & 0.531 & 0.620 & 0.691 & 0.722 & \makecell{0.571\\{\scriptsize (0.059; 0.472--0.668)}} &  & interpreted \\
 & E1-L & 0.439 & 0.521 & 0.608 & 0.692 & 0.745 & \makecell{0.565\\{\scriptsize (0.069; 0.445--0.657)}} &  & interpreted \\
 & E1-C & 0.446 & 0.528 & 0.621 & 0.695 & 0.729 & \makecell{0.572\\{\scriptsize (0.065; 0.465--0.688)}} &  & interpreted \\
RETFound & locked & 0.255 & 0.308 & 0.368 & 0.422 & 0.476 & \makecell{0.338\\{\scriptsize (0.065; 0.225--0.420)}} &  & interpreted \\
 & E1-M & 0.282 & 0.316 & 0.357 & 0.408 & 0.455 & \makecell{0.341\\{\scriptsize (0.063; 0.252--0.438)}} &  & interpreted \\
 & E1-L & 0.329 & 0.390 & 0.450 & 0.538 & 0.646 & \makecell{0.427\\{\scriptsize (0.068; 0.332--0.546)}} &  & interpreted \\
 & E1-C & 0.281 & 0.314 & 0.362 & 0.406 & 0.460 & \makecell{0.341\\{\scriptsize (0.061; 0.252--0.438)}} &  & interpreted \\
Size-only & raw & 0.645 & 0.590 & 0.592 & 0.645 & 0.727 & \makecell{0.618\\{\scriptsize (0.101; 0.424--0.751)}} &  & reference without a model \\
Content-size & raw content box & 0.405 & 0.426 & 0.480 & 0.509 & 0.521 & \makecell{0.455\\{\scriptsize (0.057; 0.363--0.545)}} &  & reference without a model \\
Thumbnail & E1-M & 0.272 & 0.262 & 0.279 & 0.287 & 0.287 & \makecell{0.275\\{\scriptsize (0.064; 0.179--0.373)}} &  & reference without a model \\
 & E1-L & 0.269 & 0.263 & 0.277 & 0.285 & 0.289 & \makecell{0.273\\{\scriptsize (0.065; 0.170--0.373)}} &  & reference without a model \\
 & E1-C & 0.251 & 0.246 & 0.272 & 0.299 & 0.321 & \makecell{0.267\\{\scriptsize (0.068; 0.150--0.373)}} &  & reference without a model \\
 & raw & 0.230 & 0.227 & 0.244 & 0.246 & 0.234 & \makecell{0.237\\{\scriptsize (0.061; 0.137--0.331)}} &  & reference without a model \\
\multicolumn{10}{@{}l}{\emph{Mean number named correctly (of 22)}} \\
MIRAGE & locked & 11.880 & 15.090 & 17.660 & 19.230 & 19.870 & -- &  & interpreted \\
 & E1-M & 11.620 & 13.860 & 15.580 & 17.370 & 18.590 & -- &  & interpreted \\
 & E1-L & 12.240 & 14.140 & 16.270 & 17.950 & 18.650 & -- &  & interpreted \\
 & E1-C & 11.460 & 13.700 & 15.880 & 18.020 & 19.290 & -- &  & interpreted \\
DINOv2 & locked & 8.660 & 10.490 & 12.770 & 14.580 & 15.080 & -- &  & interpreted \\
 & E1-M & 9.440 & 11.760 & 13.720 & 14.860 & 15.030 & -- &  & interpreted \\
 & E1-L & 9.260 & 11.290 & 12.910 & 13.620 & 13.970 & -- &  & interpreted \\
 & E1-C & 9.230 & 11.630 & 13.650 & 14.730 & 14.710 & -- &  & interpreted \\
Inception-v3 & locked & 9.690 & 11.210 & 12.670 & 14.000 & 14.220 & -- &  & interpreted \\
 & E1-M & 9.970 & 11.970 & 13.960 & 15.500 & 16.180 & -- &  & interpreted \\
 & E1-L & 9.900 & 11.710 & 13.660 & 15.530 & 16.640 & -- &  & interpreted \\
 & E1-C & 9.980 & 11.920 & 14.110 & 15.790 & 16.530 & -- &  & interpreted \\
RETFound & locked & 5.570 & 6.590 & 7.870 & 8.940 & 10.020 & -- &  & interpreted \\
 & E1-M & 6.300 & 6.990 & 7.940 & 9.050 & 10.020 & -- &  & interpreted \\
 & E1-L & 7.280 & 8.600 & 9.940 & 11.850 & 14.080 & -- &  & interpreted \\
 & E1-C & 6.270 & 6.960 & 8.040 & 9.020 & 10.130 & -- &  & interpreted \\
\multicolumn{10}{@{}l}{\emph{MIRAGE minus Inception-v3, paired (direction only)}} \\
MIRAGE $-$ Inception-v3 & locked & \makecell{+0.106\\{\scriptsize ($-$0.121 to +0.350)}} & \makecell{+0.188\\{\scriptsize ($-$0.020 to +0.440)}} & \makecell{+0.243\\{\scriptsize (+0.030 to +0.440)}} & \makecell{+0.253\\{\scriptsize (+0.060 to +0.390)}} & \makecell{+0.273\\{\scriptsize (+0.189 to +0.342)}} & \makecell{+0.197\\{\scriptsize (0.085; +0.042 to +0.343)}} & positive & interpreted \\
 & E1-M & \makecell{+0.086\\{\scriptsize ($-$0.121 to +0.322)}} & \makecell{+0.101\\{\scriptsize ($-$0.132 to +0.341)}} & \makecell{+0.093\\{\scriptsize ($-$0.071 to +0.270)}} & \makecell{+0.106\\{\scriptsize ($-$0.061 to +0.210)}} & \makecell{+0.131\\{\scriptsize (+0.010 to +0.210)}} & \makecell{+0.096\\{\scriptsize (0.074; $-$0.026 to +0.218)}} & positive & interpreted \\
 & E1-L & \makecell{+0.117\\{\scriptsize ($-$0.080 to +0.320)}} & \makecell{+0.124\\{\scriptsize ($-$0.093 to +0.341)}} & \makecell{+0.137\\{\scriptsize ($-$0.070 to +0.311)}} & \makecell{+0.131\\{\scriptsize ($-$0.021 to +0.262)}} & \makecell{+0.111\\{\scriptsize ($-$0.030 to +0.250)}} & \makecell{+0.127\\{\scriptsize (0.077; $-$0.000 to +0.250)}} & positive & interpreted \\
 & E1-C & \makecell{+0.081\\{\scriptsize ($-$0.090 to +0.271)}} & \makecell{+0.102\\{\scriptsize ($-$0.131 to +0.301)}} & \makecell{+0.111\\{\scriptsize ($-$0.080 to +0.311)}} & \makecell{+0.135\\{\scriptsize (+0.010 to +0.261)}} & \makecell{+0.161\\{\scriptsize (+0.030 to +0.270)}} & \makecell{+0.107\\{\scriptsize (0.072; $-$0.008 to +0.225)}} & positive & interpreted \\
\end{longtable}
}
\fi
\havetable{e1h_leave_one_out}\ifgen
\subsubsection*{E1h. Each development query left out}
Smallest and largest MIRAGE minus Inception-v3 mean over $K$ 1--8 with one development query left out (as P1). locked: locked = P1, shown for reference.

\begin{center}\footnotesize% src: E1h_leave_one_out.csv in results/fm_v1/posthoc (generated by make_tables.py; formatting only);
% label: post hoc, Phase-1 rule E1 timestamped by commit 4b4e883 (20 Nov 2026), computed 2026-11-21; outside the locked protocol; never replaces a locked result; limitations: the centre square drops lateral retina unequally by source and class; MIRAGE now gets 256 px of detail enlarged to 512; E1-M does not remove effective resolution (content short side 103-496 px), field of view, scanner appearance, compression, intensity, black-clipped background against speckle, or Kermany's white fill wedges
\begin{tabular}{@{}llrcllll@{}}
\toprule
Set & Status & Left out & Full & Smallest & (image) & Largest & (image) \\
\midrule
locked & interpreted & 124 & +0.116 & +0.109 & MH/23.PNG & +0.122 & Stargardt/19.PNG \\
E1-M & interpreted & 124 & +0.063 & +0.058 & MacTel/29.PNG & +0.068 & Stargardt/10.PNG \\
\bottomrule
\end{tabular}
\end{center}
\fi

\havetable{e2_label}\ifgen
\subsection*{E2. External error analysis}
Label of every E2 row: post hoc, Phase-1 rule E2 timestamped by commit 4b4e883 (20 Nov 2026), computed 2026-11-21; outside the locked protocol; never replaces a locked result
. A failure is a normal scan given any class but NORMAL; a false rare, one given a rare class. Descriptive only: no cause, no test, no clinical reading.
\fi
\havetable{e2a_destinations}\ifgen
\subsubsection*{E2a. Destinations}
Counts in each of the nine classes and the three groups. External disease images are shown here only. 

{\scriptsize\setlength{\tabcolsep}{2.5pt}% src: E2a_destinations.csv in results/fm_v1/posthoc (generated by make_tables.py; formatting only);
% label: post hoc, Phase-1 rule E2 timestamped by commit 4b4e883 (20 Nov 2026), computed 2026-11-21; outside the locked protocol; never replaces a locked result
\begin{longtable}{@{}lll rrrrrrrrrrc@{}}
\toprule
Arm & Set & Query set & $n$ & NOR & CNV & DME & DRU & CSC & MH & MacT & Starg\textsuperscript{a} & RP\textsuperscript{a} & \makecell{NORMAL / common\\/ rare} \\
\midrule
\endfirsthead
\toprule
Arm & Set & Query set & $n$ & NOR & CNV & DME & DRU & CSC & MH & MacT & Starg\textsuperscript{a} & RP\textsuperscript{a} & \makecell{NORMAL / common\\/ rare} \\
\midrule
\endhead
\bottomrule
\endfoot
\bottomrule
\endlastfoot
MIRAGE & locked & Kermany patient-disjoint & 169 & 169 & 0 & 0 & 0 & 0 & 0 & 0 & 0 & 0 & 169 / 0 / 0 \\
 &  & OCTID NORMAL & 206 & 202 & 0 & 0 & 0 & 0 & 0 & 4 & 0 & 0 & 202 / 0 / 4 \\
 &  & OCTDL NORMAL & 332 & 26 & 148 & 118 & 10 & 0 & 6 & 24 & 0 & 0 & 26 / 276 / 30 \\
 &  & OCTID CSC & 102 & 8 & 28 & 23 & 2 & 23 & 1 & 1 & 15 & 1 & 8 / 53 / 41 \\
 &  & OCTID MH & 102 & 8 & 10 & 72 & 1 & 0 & 10 & 0 & 1 & 0 & 8 / 83 / 11 \\
 &  & OCTDL MH & 39 & 0 & 6 & 30 & 0 & 0 & 3 & 0 & 0 & 0 & 0 / 36 / 3 \\
\addlinespace[1pt]
MIRAGE & E1-M & Kermany patient-disjoint & 169 & 169 & 0 & 0 & 0 & 0 & 0 & 0 & 0 & 0 & 169 / 0 / 0 \\
 &  & OCTID NORMAL & 206 & 206 & 0 & 0 & 0 & 0 & 0 & 0 & 0 & 0 & 206 / 0 / 0 \\
 &  & OCTDL NORMAL & 332 & 65 & 2 & 2 & 2 & 0 & 30 & 130 & 8 & 93 & 65 / 6 / 261 \\
 &  & OCTID CSC & 102 & 17 & 32 & 19 & 0 & 6 & 4 & 0 & 24 & 0 & 17 / 51 / 34 \\
 &  & OCTID MH & 102 & 8 & 5 & 76 & 0 & 0 & 11 & 0 & 1 & 1 & 8 / 81 / 13 \\
 &  & OCTDL MH & 39 & 0 & 0 & 0 & 0 & 0 & 36 & 0 & 2 & 1 & 0 / 0 / 39 \\
\addlinespace[1pt]
MIRAGE & E1-L & Kermany patient-disjoint & 169 & 165 & 0 & 0 & 4 & 0 & 0 & 0 & 0 & 0 & 165 / 4 / 0 \\
 &  & OCTID NORMAL & 206 & 206 & 0 & 0 & 0 & 0 & 0 & 0 & 0 & 0 & 206 / 0 / 0 \\
 &  & OCTDL NORMAL & 332 & 38 & 2 & 0 & 1 & 0 & 21 & 107 & 8 & 155 & 38 / 3 / 291 \\
 &  & OCTID CSC & 102 & 23 & 32 & 17 & 0 & 8 & 2 & 0 & 19 & 1 & 23 / 49 / 30 \\
 &  & OCTID MH & 102 & 13 & 10 & 65 & 0 & 0 & 12 & 0 & 1 & 1 & 13 / 75 / 14 \\
 &  & OCTDL MH & 39 & 0 & 0 & 0 & 0 & 0 & 36 & 0 & 2 & 1 & 0 / 0 / 39 \\
\addlinespace[1pt]
MIRAGE & E1-C & Kermany patient-disjoint & 169 & 169 & 0 & 0 & 0 & 0 & 0 & 0 & 0 & 0 & 169 / 0 / 0 \\
 &  & OCTID NORMAL & 206 & 206 & 0 & 0 & 0 & 0 & 0 & 0 & 0 & 0 & 206 / 0 / 0 \\
 &  & OCTDL NORMAL & 332 & 29 & 55 & 12 & 8 & 1 & 196 & 27 & 4 & 0 & 29 / 75 / 228 \\
 &  & OCTID CSC & 102 & 12 & 40 & 34 & 2 & 3 & 2 & 0 & 9 & 0 & 12 / 76 / 14 \\
 &  & OCTID MH & 102 & 10 & 11 & 77 & 1 & 0 & 2 & 0 & 1 & 0 & 10 / 89 / 3 \\
 &  & OCTDL MH & 39 & 0 & 2 & 7 & 0 & 0 & 30 & 0 & 0 & 0 & 0 / 9 / 30 \\
\addlinespace[1pt]
DINOv2 & locked & Kermany patient-disjoint & 169 & 146 & 5 & 2 & 10 & 2 & 0 & 0 & 0 & 4 & 146 / 17 / 6 \\
 &  & OCTID NORMAL & 206 & 188 & 1 & 0 & 0 & 0 & 0 & 3 & 3 & 11 & 188 / 1 / 17 \\
 &  & OCTDL NORMAL & 332 & 145 & 29 & 1 & 16 & 6 & 0 & 10 & 58 & 67 & 145 / 46 / 141 \\
 &  & OCTID CSC & 102 & 11 & 19 & 1 & 2 & 30 & 0 & 3 & 30 & 6 & 11 / 22 / 69 \\
 &  & OCTID MH & 102 & 4 & 2 & 6 & 5 & 2 & 48 & 23 & 2 & 10 & 4 / 13 / 85 \\
 &  & OCTDL MH & 39 & 0 & 0 & 0 & 0 & 1 & 26 & 8 & 2 & 2 & 0 / 0 / 39 \\
\addlinespace[1pt]
DINOv2 & E1-M & Kermany patient-disjoint & 169 & 153 & 4 & 0 & 11 & 1 & 0 & 0 & 0 & 0 & 153 / 15 / 1 \\
 &  & OCTID NORMAL & 206 & 203 & 1 & 0 & 0 & 0 & 0 & 1 & 0 & 1 & 203 / 1 / 2 \\
 &  & OCTDL NORMAL & 332 & 162 & 48 & 6 & 36 & 1 & 0 & 21 & 1 & 57 & 162 / 90 / 80 \\
 &  & OCTID CSC & 102 & 21 & 38 & 8 & 3 & 12 & 0 & 7 & 9 & 4 & 21 / 49 / 32 \\
 &  & OCTID MH & 102 & 7 & 2 & 16 & 5 & 2 & 36 & 26 & 0 & 8 & 7 / 23 / 72 \\
 &  & OCTDL MH & 39 & 0 & 2 & 0 & 0 & 0 & 27 & 9 & 0 & 1 & 0 / 2 / 37 \\
\addlinespace[1pt]
DINOv2 & E1-L & Kermany patient-disjoint & 169 & 149 & 2 & 3 & 15 & 0 & 0 & 0 & 0 & 0 & 149 / 20 / 0 \\
 &  & OCTID NORMAL & 206 & 56 & 0 & 0 & 0 & 0 & 0 & 120 & 1 & 29 & 56 / 0 / 150 \\
 &  & OCTDL NORMAL & 332 & 128 & 33 & 6 & 58 & 4 & 0 & 20 & 15 & 68 & 128 / 97 / 107 \\
 &  & OCTID CSC & 102 & 3 & 13 & 3 & 0 & 28 & 0 & 21 & 30 & 4 & 3 / 16 / 83 \\
 &  & OCTID MH & 102 & 1 & 3 & 8 & 1 & 5 & 50 & 22 & 3 & 9 & 1 / 12 / 89 \\
 &  & OCTDL MH & 39 & 1 & 1 & 1 & 0 & 3 & 29 & 3 & 0 & 1 & 1 / 2 / 36 \\
\addlinespace[1pt]
DINOv2 & E1-C & Kermany patient-disjoint & 169 & 152 & 4 & 1 & 11 & 1 & 0 & 0 & 0 & 0 & 152 / 16 / 1 \\
 &  & OCTID NORMAL & 206 & 203 & 1 & 0 & 0 & 0 & 0 & 2 & 0 & 0 & 203 / 1 / 2 \\
 &  & OCTDL NORMAL & 332 & 180 & 48 & 10 & 31 & 1 & 0 & 20 & 0 & 42 & 180 / 89 / 63 \\
 &  & OCTID CSC & 102 & 21 & 39 & 8 & 3 & 11 & 0 & 7 & 10 & 3 & 21 / 50 / 31 \\
 &  & OCTID MH & 102 & 8 & 2 & 16 & 5 & 3 & 32 & 28 & 1 & 7 & 8 / 23 / 71 \\
 &  & OCTDL MH & 39 & 0 & 2 & 0 & 0 & 0 & 27 & 9 & 0 & 1 & 0 / 2 / 37 \\
\addlinespace[1pt]
Inception-v3 & locked & Kermany patient-disjoint & 169 & 150 & 0 & 6 & 9 & 0 & 0 & 3 & 1 & 0 & 150 / 15 / 4 \\
 &  & OCTID NORMAL & 206 & 182 & 0 & 9 & 0 & 12 & 0 & 3 & 0 & 0 & 182 / 9 / 15 \\
 &  & OCTDL NORMAL & 332 & 66 & 196 & 21 & 38 & 0 & 0 & 0 & 0 & 11 & 66 / 255 / 11 \\
 &  & OCTID CSC & 102 & 3 & 16 & 26 & 11 & 22 & 0 & 2 & 14 & 8 & 3 / 53 / 46 \\
 &  & OCTID MH & 102 & 6 & 5 & 45 & 2 & 2 & 27 & 0 & 1 & 14 & 6 / 52 / 44 \\
 &  & OCTDL MH & 39 & 2 & 25 & 11 & 0 & 1 & 0 & 0 & 0 & 0 & 2 / 36 / 1 \\
\addlinespace[1pt]
Inception-v3 & E1-M & Kermany patient-disjoint & 169 & 160 & 0 & 0 & 7 & 1 & 0 & 0 & 0 & 1 & 160 / 7 / 2 \\
 &  & OCTID NORMAL & 206 & 125 & 0 & 1 & 4 & 14 & 0 & 38 & 6 & 18 & 125 / 5 / 76 \\
 &  & OCTDL NORMAL & 332 & 115 & 3 & 52 & 68 & 2 & 0 & 18 & 3 & 71 & 115 / 123 / 94 \\
 &  & OCTID CSC & 102 & 4 & 20 & 7 & 33 & 14 & 1 & 1 & 17 & 5 & 4 / 60 / 38 \\
 &  & OCTID MH & 102 & 4 & 10 & 43 & 22 & 2 & 16 & 0 & 2 & 3 & 4 / 75 / 23 \\
 &  & OCTDL MH & 39 & 0 & 1 & 6 & 2 & 1 & 29 & 0 & 0 & 0 & 0 / 9 / 30 \\
\addlinespace[1pt]
Inception-v3 & E1-L & Kermany patient-disjoint & 169 & 155 & 0 & 0 & 12 & 0 & 0 & 0 & 0 & 2 & 155 / 12 / 2 \\
 &  & OCTID NORMAL & 206 & 123 & 0 & 0 & 15 & 15 & 0 & 0 & 22 & 31 & 123 / 15 / 68 \\
 &  & OCTDL NORMAL & 332 & 152 & 12 & 34 & 64 & 0 & 0 & 15 & 4 & 51 & 152 / 110 / 70 \\
 &  & OCTID CSC & 102 & 4 & 16 & 9 & 32 & 16 & 0 & 1 & 21 & 3 & 4 / 57 / 41 \\
 &  & OCTID MH & 102 & 5 & 11 & 39 & 17 & 2 & 19 & 0 & 4 & 5 & 5 / 67 / 30 \\
 &  & OCTDL MH & 39 & 0 & 3 & 5 & 0 & 0 & 31 & 0 & 0 & 0 & 0 / 8 / 31 \\
\addlinespace[1pt]
Inception-v3 & E1-C & Kermany patient-disjoint & 169 & 148 & 0 & 11 & 6 & 0 & 0 & 1 & 1 & 2 & 148 / 17 / 4 \\
 &  & OCTID NORMAL & 206 & 146 & 0 & 9 & 0 & 2 & 0 & 47 & 1 & 1 & 146 / 9 / 51 \\
 &  & OCTDL NORMAL & 332 & 34 & 120 & 7 & 8 & 24 & 21 & 29 & 5 & 84 & 34 / 135 / 163 \\
 &  & OCTID CSC & 102 & 6 & 14 & 32 & 3 & 15 & 0 & 2 & 27 & 3 & 6 / 49 / 47 \\
 &  & OCTID MH & 102 & 5 & 3 & 48 & 1 & 2 & 24 & 0 & 2 & 17 & 5 / 52 / 45 \\
 &  & OCTDL MH & 39 & 0 & 3 & 4 & 0 & 6 & 23 & 2 & 0 & 1 & 0 / 7 / 32 \\
\addlinespace[1pt]
RETFound & locked & Kermany patient-disjoint & 169 & 149 & 1 & 1 & 17 & 0 & 0 & 1 & 0 & 0 & 149 / 19 / 1 \\
 &  & OCTID NORMAL & 206 & 189 & 0 & 3 & 14 & 0 & 0 & 0 & 0 & 0 & 189 / 17 / 0 \\
 &  & OCTDL NORMAL & 332 & 137 & 8 & 167 & 2 & 0 & 0 & 12 & 1 & 5 & 137 / 177 / 18 \\
 &  & OCTID CSC & 102 & 13 & 12 & 72 & 5 & 0 & 0 & 0 & 0 & 0 & 13 / 89 / 0 \\
 &  & OCTID MH & 102 & 8 & 3 & 80 & 11 & 0 & 0 & 0 & 0 & 0 & 8 / 94 / 0 \\
 &  & OCTDL MH & 39 & 0 & 0 & 36 & 0 & 0 & 0 & 0 & 0 & 3 & 0 / 36 / 3 \\
\addlinespace[1pt]
RETFound & E1-M & Kermany patient-disjoint & 169 & 149 & 0 & 0 & 19 & 0 & 0 & 0 & 1 & 0 & 149 / 19 / 1 \\
 &  & OCTID NORMAL & 206 & 198 & 0 & 8 & 0 & 0 & 0 & 0 & 0 & 0 & 198 / 8 / 0 \\
 &  & OCTDL NORMAL & 332 & 92 & 9 & 99 & 1 & 0 & 0 & 16 & 115 & 0 & 92 / 109 / 131 \\
 &  & OCTID CSC & 102 & 12 & 9 & 78 & 1 & 0 & 0 & 0 & 2 & 0 & 12 / 88 / 2 \\
 &  & OCTID MH & 102 & 9 & 3 & 87 & 3 & 0 & 0 & 0 & 0 & 0 & 9 / 93 / 0 \\
 &  & OCTDL MH & 39 & 0 & 0 & 23 & 0 & 0 & 3 & 0 & 13 & 0 & 0 / 23 / 16 \\
\addlinespace[1pt]
RETFound & E1-L & Kermany patient-disjoint & 169 & 154 & 0 & 4 & 11 & 0 & 0 & 0 & 0 & 0 & 154 / 15 / 0 \\
 &  & OCTID NORMAL & 206 & 157 & 0 & 0 & 8 & 0 & 0 & 0 & 39 & 2 & 157 / 8 / 41 \\
 &  & OCTDL NORMAL & 332 & 31 & 1 & 0 & 1 & 0 & 0 & 237 & 36 & 26 & 31 / 2 / 299 \\
 &  & OCTID CSC & 102 & 36 & 7 & 20 & 9 & 0 & 0 & 1 & 27 & 2 & 36 / 36 / 30 \\
 &  & OCTID MH & 102 & 17 & 1 & 32 & 6 & 0 & 42 & 0 & 3 & 1 & 17 / 39 / 46 \\
 &  & OCTDL MH & 39 & 0 & 0 & 0 & 0 & 0 & 34 & 3 & 2 & 0 & 0 / 0 / 39 \\
\addlinespace[1pt]
RETFound & E1-C & Kermany patient-disjoint & 169 & 149 & 0 & 0 & 19 & 0 & 0 & 0 & 1 & 0 & 149 / 19 / 1 \\
 &  & OCTID NORMAL & 206 & 199 & 0 & 7 & 0 & 0 & 0 & 0 & 0 & 0 & 199 / 7 / 0 \\
 &  & OCTDL NORMAL & 332 & 92 & 9 & 102 & 1 & 0 & 0 & 13 & 115 & 0 & 92 / 112 / 128 \\
 &  & OCTID CSC & 102 & 12 & 10 & 77 & 1 & 0 & 0 & 0 & 2 & 0 & 12 / 88 / 2 \\
 &  & OCTID MH & 102 & 9 & 3 & 87 & 3 & 0 & 0 & 0 & 0 & 0 & 9 / 93 / 0 \\
 &  & OCTDL MH & 39 & 0 & 0 & 23 & 0 & 0 & 3 & 0 & 13 & 0 & 0 / 23 / 16 \\
\end{longtable}
}
\fi
\havetable{e2b_pairs}\ifgen
\subsubsection*{E2b. Overlap between pairs of models}
Jaccard index of the failure sets and of the false-rare sets, and the share of scans given the same class. --- : both sets empty (no Jaccard index). 

{\scriptsize% src: E2b_overlap.csv in results/fm_v1/posthoc (generated by make_tables.py; formatting only);
% label: post hoc, Phase-1 rule E2 timestamped by commit 4b4e883 (20 Nov 2026), computed 2026-11-21; outside the locked protocol; never replaces a locked result
\begin{longtable}{@{}lllccc@{}}
\toprule
Normal scans & Set & Arms & Failure Jaccard & False-rare Jaccard & Same class \\
\midrule
\endfirsthead
\toprule
Normal scans & Set & Arms & Failure Jaccard & False-rare Jaccard & Same class \\
\midrule
\endhead
\bottomrule
\endfoot
\bottomrule
\endlastfoot
Kermany patient-disjoint & locked & MIRAGE / Inception-v3 & 0 & 0 & 0.888 \\
 &  & DINOv2 / MIRAGE & 0 & 0 & 0.864 \\
 &  & DINOv2 / Inception-v3 & 0.024 & 0 & 0.757 \\
 &  & RETFound / MIRAGE & 0 & 0 & 0.882 \\
 &  & RETFound / DINOv2 & 0.049 & 0 & 0.763 \\
 &  & RETFound / Inception-v3 & 0.054 & 0 & 0.787 \\
\addlinespace[1pt]
Kermany patient-disjoint & E1-M & MIRAGE / Inception-v3 & 0 & 0 & 0.947 \\
 &  & DINOv2 / MIRAGE & 0 & 0 & 0.905 \\
 &  & DINOv2 / Inception-v3 & 0 & 0 & 0.852 \\
 &  & RETFound / MIRAGE & 0 & 0 & 0.882 \\
 &  & RETFound / DINOv2 & 0.059 & 0 & 0.811 \\
 &  & RETFound / Inception-v3 & 0.036 & 0 & 0.840 \\
\addlinespace[1pt]
Kermany patient-disjoint & E1-L & MIRAGE / Inception-v3 & 0.125 & 0 & 0.917 \\
 &  & DINOv2 / MIRAGE & 0.043 & --- & 0.870 \\
 &  & DINOv2 / Inception-v3 & 0.097 & 0 & 0.834 \\
 &  & RETFound / MIRAGE & 0.118 & --- & 0.899 \\
 &  & RETFound / DINOv2 & 0.094 & --- & 0.822 \\
 &  & RETFound / Inception-v3 & 0.115 & 0 & 0.858 \\
\addlinespace[1pt]
Kermany patient-disjoint & E1-C & MIRAGE / Inception-v3 & 0 & 0 & 0.876 \\
 &  & DINOv2 / MIRAGE & 0 & 0 & 0.899 \\
 &  & DINOv2 / Inception-v3 & 0.027 & 0 & 0.787 \\
 &  & RETFound / MIRAGE & 0 & 0 & 0.882 \\
 &  & RETFound / DINOv2 & 0.088 & 0 & 0.817 \\
 &  & RETFound / Inception-v3 & 0.051 & 0 & 0.775 \\
\addlinespace[1pt]
OCTID NORMAL & locked & MIRAGE / Inception-v3 & 0.077 & 0.118 & 0.874 \\
 &  & DINOv2 / MIRAGE & 0.048 & 0.050 & 0.898 \\
 &  & DINOv2 / Inception-v3 & 0.077 & 0.032 & 0.811 \\
 &  & RETFound / MIRAGE & 0 & 0 & 0.898 \\
 &  & RETFound / DINOv2 & 0.029 & 0 & 0.835 \\
 &  & RETFound / Inception-v3 & 0.051 & 0 & 0.816 \\
\addlinespace[1pt]
OCTID NORMAL & E1-M & MIRAGE / Inception-v3 & 0 & 0 & 0.607 \\
 &  & DINOv2 / MIRAGE & 0 & 0 & 0.985 \\
 &  & DINOv2 / Inception-v3 & 0.012 & 0.013 & 0.602 \\
 &  & RETFound / MIRAGE & 0 & --- & 0.961 \\
 &  & RETFound / DINOv2 & 0 & 0 & 0.947 \\
 &  & RETFound / Inception-v3 & 0.047 & 0 & 0.587 \\
\addlinespace[1pt]
OCTID NORMAL & E1-L & MIRAGE / Inception-v3 & 0 & 0 & 0.597 \\
 &  & DINOv2 / MIRAGE & 0 & 0 & 0.272 \\
 &  & DINOv2 / Inception-v3 & 0.395 & 0.313 & 0.257 \\
 &  & RETFound / MIRAGE & 0 & 0 & 0.762 \\
 &  & RETFound / DINOv2 & 0.221 & 0.179 & 0.214 \\
 &  & RETFound / Inception-v3 & 0.179 & 0.101 & 0.490 \\
\addlinespace[1pt]
OCTID NORMAL & E1-C & MIRAGE / Inception-v3 & 0 & 0 & 0.709 \\
 &  & DINOv2 / MIRAGE & 0 & 0 & 0.985 \\
 &  & DINOv2 / Inception-v3 & 0.016 & 0.019 & 0.704 \\
 &  & RETFound / MIRAGE & 0 & --- & 0.966 \\
 &  & RETFound / DINOv2 & 0 & 0 & 0.951 \\
 &  & RETFound / Inception-v3 & 0.015 & 0 & 0.684 \\
\addlinespace[1pt]
OCTDL NORMAL & locked & MIRAGE / Inception-v3 & 0.749 & 0.079 & 0.349 \\
 &  & DINOv2 / MIRAGE & 0.541 & 0.075 & 0.081 \\
 &  & DINOv2 / Inception-v3 & 0.551 & 0.041 & 0.223 \\
 &  & RETFound / MIRAGE & 0.596 & 0.091 & 0.289 \\
 &  & RETFound / DINOv2 & 0.431 & 0.053 & 0.199 \\
 &  & RETFound / Inception-v3 & 0.579 & 0.036 & 0.166 \\
\addlinespace[1pt]
OCTDL NORMAL & E1-M & MIRAGE / Inception-v3 & 0.603 & 0.300 & 0.175 \\
 &  & DINOv2 / MIRAGE & 0.491 & 0.231 & 0.196 \\
 &  & DINOv2 / Inception-v3 & 0.455 & 0.270 & 0.328 \\
 &  & RETFound / MIRAGE & 0.730 & 0.446 & 0.148 \\
 &  & RETFound / DINOv2 & 0.439 & 0.172 & 0.157 \\
 &  & RETFound / Inception-v3 & 0.592 & 0.216 & 0.205 \\
\addlinespace[1pt]
OCTDL NORMAL & E1-L & MIRAGE / Inception-v3 & 0.495 & 0.224 & 0.145 \\
 &  & DINOv2 / MIRAGE & 0.596 & 0.318 & 0.199 \\
 &  & DINOv2 / Inception-v3 & 0.548 & 0.221 & 0.377 \\
 &  & RETFound / MIRAGE & 0.865 & 0.855 & 0.416 \\
 &  & RETFound / DINOv2 & 0.588 & 0.318 & 0.099 \\
 &  & RETFound / Inception-v3 & 0.513 & 0.206 & 0.096 \\
\addlinespace[1pt]
OCTDL NORMAL & E1-C & MIRAGE / Inception-v3 & 0.827 & 0.487 & 0.157 \\
 &  & DINOv2 / MIRAGE & 0.472 & 0.169 & 0.117 \\
 &  & DINOv2 / Inception-v3 & 0.461 & 0.183 & 0.199 \\
 &  & RETFound / MIRAGE & 0.724 & 0.441 & 0.063 \\
 &  & RETFound / DINOv2 & 0.400 & 0.124 & 0.172 \\
 &  & RETFound / Inception-v3 & 0.687 & 0.335 & 0.057 \\
\end{longtable}
}
\fi
\havetable{e2c_transitions}\ifgen
\subsubsection*{E2c. Locked to E1-M transitions}
Locked $\to$ E1-M, per scan: N, called NORMAL; C, a common disease (CNV, DME or DRUSEN); R, any rare class. 

{\scriptsize% src: E2c_transitions.csv in C:/Users/ThinkBook/.t3/worktrees/Few-shot/t3-03f3a2c4/results/fm_v1/posthoc (generated by make_tables.py; formatting only);
% label: post hoc, Phase-1 rule E2 timestamped by commit 4b4e883 (20 Nov 2026), computed 2026-11-21; outside the locked protocol; never replaces a locked result
\begin{longtable}{@{}llrrrrrrrrr@{}}
\toprule
Arm & Scans & N$\to$N & N$\to$C & N$\to$R & C$\to$N & C$\to$C & C$\to$R & R$\to$N & R$\to$C & R$\to$R \\
\midrule
\endfirsthead
\toprule
Arm & Scans & N$\to$N & N$\to$C & N$\to$R & C$\to$N & C$\to$C & C$\to$R & R$\to$N & R$\to$C & R$\to$R \\
\midrule
\endhead
\bottomrule
\endfoot
\bottomrule
\endlastfoot
MIRAGE & Kermany patient-disjoint & 169 & 0 & 0 & 0 & 0 & 0 & 0 & 0 & 0 \\
 & OCTID NORMAL & 202 & 0 & 0 & 0 & 0 & 0 & 4 & 0 & 0 \\
 & OCTDL NORMAL & 18 & 0 & 8 & 42 & 6 & 228 & 5 & 0 & 25 \\
 & OCTID CSC & 6 & 2 & 0 & 6 & 37 & 10 & 5 & 12 & 24 \\
 & OCTID MH & 5 & 3 & 0 & 3 & 71 & 9 & 0 & 7 & 4 \\
 & OCTDL MH & 0 & 0 & 0 & 0 & 0 & 36 & 0 & 0 & 3 \\
\addlinespace[1pt]
DINOv2 & Kermany patient-disjoint & 145 & 1 & 0 & 3 & 14 & 0 & 5 & 0 & 1 \\
 & OCTID NORMAL & 187 & 0 & 1 & 0 & 1 & 0 & 16 & 0 & 1 \\
 & OCTDL NORMAL & 125 & 3 & 17 & 0 & 45 & 1 & 37 & 42 & 62 \\
 & OCTID CSC & 11 & 0 & 0 & 0 & 22 & 0 & 10 & 27 & 32 \\
 & OCTID MH & 4 & 0 & 0 & 0 & 13 & 0 & 3 & 10 & 72 \\
 & OCTDL MH & 0 & 0 & 0 & 0 & 0 & 0 & 0 & 2 & 37 \\
\addlinespace[1pt]
Inception-v3 & Kermany patient-disjoint & 148 & 1 & 1 & 9 & 5 & 1 & 3 & 1 & 0 \\
 & OCTID NORMAL & 117 & 3 & 62 & 4 & 2 & 3 & 4 & 0 & 11 \\
 & OCTDL NORMAL & 39 & 9 & 18 & 74 & 109 & 72 & 2 & 5 & 4 \\
 & OCTID CSC & 1 & 1 & 1 & 2 & 36 & 15 & 1 & 23 & 22 \\
 & OCTID MH & 1 & 1 & 4 & 2 & 44 & 6 & 1 & 30 & 13 \\
 & OCTDL MH & 0 & 0 & 2 & 0 & 9 & 27 & 0 & 0 & 1 \\
\addlinespace[1pt]
RETFound & Kermany patient-disjoint & 148 & 1 & 0 & 1 & 18 & 0 & 0 & 0 & 1 \\
 & OCTID NORMAL & 188 & 1 & 0 & 10 & 7 & 0 & 0 & 0 & 0 \\
 & OCTDL NORMAL & 88 & 0 & 49 & 4 & 109 & 64 & 0 & 0 & 18 \\
 & OCTID CSC & 10 & 3 & 0 & 2 & 85 & 2 & 0 & 0 & 0 \\
 & OCTID MH & 7 & 1 & 0 & 2 & 92 & 0 & 0 & 0 & 0 \\
 & OCTDL MH & 0 & 0 & 0 & 0 & 23 & 13 & 0 & 0 & 3 \\
\end{longtable}
}
\fi
\havetable{e2d_properties}\ifgen
\subsubsection*{E2d. Image properties, all sets}
45 of 133 reported E2d intervals exclude 0.5 (all arms, sets, feature sets and properties). 

{\tiny\setlength{\tabcolsep}{2pt}% src: E2d_properties.csv in C:/Users/ThinkBook/.t3/worktrees/Few-shot/t3-03f3a2c4/results/fm_v1/posthoc (generated by make_tables.py; formatting only);
% label: post hoc, Phase-1 rule E2 timestamped by commit 4b4e883 (20 Nov 2026), computed 2026-11-21; outside the locked protocol; never replaces a locked result
\begin{longtable}{@{}llllccccc@{}}
\toprule
Arm & Normals & Set & Property & \makecell{Failed /\\not failed} & \makecell{Failed: median\\(q1--q3)} & \makecell{Not failed: median\\(q1--q3)} & \makecell{Failure share\\by tertile} & \makecell{AUROC failed vs not\\(95\% interval)} \\
\midrule
\endfirsthead
\toprule
Arm & Normals & Set & Property & \makecell{Failed /\\not failed} & \makecell{Failed: median\\(q1--q3)} & \makecell{Not failed: median\\(q1--q3)} & \makecell{Failure share\\by tertile} & \makecell{AUROC failed vs not\\(95\% interval)} \\
\midrule
\endhead
\bottomrule
\endfoot
\bottomrule
\endlastfoot
RETFound & Kermany patient-disjoint & locked & Width (px) & 20 / 149 & 512.000 (512.000--512.000) & 512.000 (512.000--512.000) & 0.119 / -- / 0.111 & 0.498 (0.460--0.560) \\
 &  &  & Height (px) & 20 / 149 & 496.000 (496.000--496.000) & 496.000 (496.000--496.000) & 0.118 / -- / -- & 0.500 (0.500--0.500) \\
 &  &  & Aspect ratio & 20 / 149 & 1.032 (1.032--1.032) & 1.032 (1.032--1.032) & 0.119 / -- / 0.111 & 0.498 (0.460--0.558) \\
 &  &  & Border share & 20 / 149 & 0.016 (0--0.052) & 0 (0--0.016) & 0.090 / 0 / 0.196 & 0.634 (0.496--0.763) \\
 &  &  & Near-white share & 20 / 149 & 0.044 (0.010--0.071) & 0.021 (0.013--0.037) & 0.105 / 0.036 / 0.214 & 0.632 (0.467--0.790) \\
 &  &  & Grey mean & 20 / 149 & 54.110 (45.969--65.362) & 43.122 (38.148--47.294) & 0.035 / 0.054 / 0.268 & 0.785 (0.647--0.905) \\
 &  &  & Grey SD & 20 / 149 & 57.347 (41.565--65.584) & 48.554 (42.948--55.710) & 0.105 / 0.054 / 0.196 & 0.628 (0.468--0.775) \\
\addlinespace[1pt]
RETFound & Kermany patient-disjoint & E1-M & Width (px) & 20 / 149 & 512.000 (512.000--512.000) & 512.000 (512.000--512.000) & 0.119 / -- / 0.111 & 0.498 (0.460--0.555) \\
 &  &  & Height (px) & 20 / 149 & 496.000 (496.000--496.000) & 496.000 (496.000--496.000) & 0.118 / -- / -- & 0.500 (0.500--0.500) \\
 &  &  & Aspect ratio & 20 / 149 & 1.032 (1.032--1.032) & 1.032 (1.032--1.032) & 0.119 / -- / 0.111 & 0.498 (0.459--0.556) \\
 &  &  & Border share & 20 / 149 & 0.011 (0--0.048) & 0 (0--0.016) & 0.090 / 0.077 / 0.179 & 0.614 (0.483--0.742) \\
 &  &  & Near-white share & 20 / 149 & 0.037 (0.010--0.066) & 0.022 (0.013--0.038) & 0.123 / 0.036 / 0.196 & 0.589 (0.422--0.744) \\
 &  &  & Grey mean & 20 / 149 & 53.383 (44.177--62.137) & 43.122 (38.148--47.370) & 0.035 / 0.071 / 0.250 & 0.759 (0.625--0.876) \\
 &  &  & Grey SD & 20 / 149 & 54.599 (41.565--64.512) & 48.611 (43.522--55.856) & 0.123 / 0.054 / 0.179 & 0.585 (0.430--0.733) \\
\addlinespace[1pt]
RETFound & OCTID NORMAL & locked & Width (px) & 17 / 189 & 750.000 (750.000--750.000) & 750.000 (750.000--750.000) & 0.084 / -- / 0 & 0.489 (0.476--0.497) \\
 &  &  & Height (px) & 17 / 189 & 500.000 (500.000--500.000) & 500.000 (500.000--500.000) & 0.084 / -- / 0 & 0.489 (0.479--0.497) \\
 &  &  & Aspect ratio & 17 / 189 & 1.500 (1.500--1.500) & 1.500 (1.500--1.500) & 0.083 / -- / -- & 0.500 (0.500--0.500) \\
 &  &  & Border share & 17 / 189 & 0 (0--0) & 0 (0--0) & 0.083 / -- / -- & 0.500 (0.500--0.500) \\
 &  &  & Near-white share & 17 / 189 & 0.002 (0.002--0.003) & 0.004 (0.003--0.005) & 0.188 / 0.058 / 0 & 0.157 (0.084--0.238) \\
 &  &  & Grey mean & 17 / 189 & 59.654 (57.668--60.776) & 60.956 (59.196--62.034) & 0.130 / 0.074 / 0.043 & 0.321 (0.197--0.456) \\
 &  &  & Grey SD & 17 / 189 & 37.998 (35.219--40.150) & 42.760 (40.466--44.142) & 0.217 / 0.015 / 0.014 & 0.151 (0.073--0.247) \\
\addlinespace[1pt]
RETFound & OCTID NORMAL & E1-M & Width (px) & 8 / 198 & 750.000 (750.000--750.000) & 750.000 (750.000--750.000) & 0.040 / -- / 0 & 0.490 (0.480--0.497) \\
 &  &  & Height (px) & 8 / 198 & 500.000 (500.000--500.000) & 500.000 (500.000--500.000) & 0.040 / -- / 0 & 0.490 (0.480--0.497) \\
 &  &  & Aspect ratio & 8 / 198 & 1.500 (1.500--1.500) & 1.500 (1.500--1.500) & 0.039 / -- / -- & 0.500 (0.500--0.500) \\
 &  &  & Border share & 8 / 198 & 0 (0--0) & 0 (0--0) & 0.039 / -- / -- & 0.500 (0.500--0.500) \\
 &  &  & Near-white share & 8 / 198 & 0.002 (0.002--0.004) & 0.004 (0.003--0.005) & 0.072 / 0.043 / 0 & 0.213 (0.074--0.381) \\
 &  &  & Grey mean & 8 / 198 & 59.778 (56.794--61.619) & 60.878 (59.141--61.978) & 0.058 / 0.015 / 0.043 & 0.386 (0.162--0.612) \\
 &  &  & Grey SD & 8 / 198 & 37.090 (35.043--42.181) & 42.427 (40.114--43.980) & 0.072 / 0.015 / 0.029 & 0.266 (0.059--0.497) \\
\addlinespace[1pt]
RETFound & OCTDL NORMAL & locked & Width (px) & 195 / 137 & 882.000 (879.000--882.000) & 882.000 (880.000--1102.000) & 0.604 / 0.726 / 0.452 & 0.421 (0.337--0.509) \\
 &  &  & Height (px) & 195 / 137 & 321.000 (313.000--322.000) & 321.000 (315.000--389.000) & 0.598 / 0.664 / 0.438 & 0.459 (0.388--0.538) \\
 &  &  & Aspect ratio & 195 / 137 & 2.745 (2.739--2.807) & 2.782 (2.739--2.938) & 0.692 / 0.569 / 0.491 & 0.375 (0.304--0.449) \\
 &  &  & Border share & 195 / 137 & 0 (0--0) & 0 (0--0) & 0.587 / -- / -- & 0.500 (0.500--0.500) \\
 &  &  & Near-white share & 195 / 137 & $6.5\times10^{-4}$ ($5.0\times10^{-4}$--$7.8\times10^{-4}$) & $5.3\times10^{-4}$ ($3.7\times10^{-4}$--$6.8\times10^{-4}$) & 0.477 / 0.527 / 0.757 & 0.637 (0.552--0.716) \\
 &  &  & Grey mean & 195 / 137 & 58.054 (54.798--62.486) & 52.599 (49.812--55.456) & 0.270 / 0.627 / 0.865 & 0.807 (0.749--0.859) \\
 &  &  & Grey SD & 195 / 137 & 46.712 (44.614--48.523) & 45.038 (43.611--46.433) & 0.441 / 0.545 / 0.775 & 0.669 (0.596--0.738) \\
\addlinespace[1pt]
RETFound & OCTDL NORMAL & E1-M & Width (px) & 240 / 92 & 882.000 (879.000--912.000) & 882.000 (881.000--1102.000) & 0.778 / 0.726 / 0.644 & 0.393 (0.312--0.478) \\
 &  &  & Height (px) & 240 / 92 & 320.000 (313.000--322.000) & 322.000 (318.000--389.250) & 0.813 / 0.750 / 0.550 & 0.365 (0.286--0.438) \\
 &  &  & Aspect ratio & 240 / 92 & 2.752 (2.739--2.815) & 2.756 (2.739--2.818) & 0.717 / 0.696 / 0.755 & 0.501 (0.423--0.577) \\
 &  &  & Border share & 240 / 92 & 0 (0--0) & 0 (0--0) & 0.723 / -- / -- & 0.500 (0.500--0.500) \\
 &  &  & Near-white share & 240 / 92 & $6.1\times10^{-4}$ ($3.8\times10^{-4}$--$7.7\times10^{-4}$) & $5.8\times10^{-4}$ ($4.7\times10^{-4}$--$6.9\times10^{-4}$) & 0.730 / 0.627 / 0.811 & 0.536 (0.452--0.622) \\
 &  &  & Grey mean & 240 / 92 & 57.288 (54.292--61.942) & 51.542 (48.973--53.810) & 0.414 / 0.791 / 0.964 & 0.841 (0.786--0.889) \\
 &  &  & Grey SD & 240 / 92 & 46.411 (44.653--48.358) & 44.378 (43.202--46.106) & 0.532 / 0.773 / 0.865 & 0.691 (0.619--0.764) \\
\addlinespace[1pt]
DINOv2 & Kermany patient-disjoint & locked & Width (px) & 23 / 146 & 512.000 (512.000--512.000) & 512.000 (512.000--512.000) & 0.144 / -- / 0 & 0.469 (0.449--0.487) \\
 &  &  & Height (px) & 23 / 146 & 496.000 (496.000--496.000) & 496.000 (496.000--496.000) & 0.136 / -- / -- & 0.500 (0.500--0.500) \\
 &  &  & Aspect ratio & 23 / 146 & 1.032 (1.032--1.032) & 1.032 (1.032--1.032) & 0.144 / -- / 0 & 0.469 (0.449--0.486) \\
 &  &  & Border share & 23 / 146 & 0 (0--0.040) & 0 (0--0.016) & 0.130 / 0 / 0.179 & 0.551 (0.427--0.683) \\
 &  &  & Near-white share & 23 / 146 & 0.041 (0.018--0.071) & 0.021 (0.012--0.036) & 0.070 / 0.071 / 0.268 & 0.666 (0.537--0.790) \\
 &  &  & Grey mean & 23 / 146 & 45.194 (43.348--53.473) & 43.146 (38.254--48.601) & 0.070 / 0.143 / 0.196 & 0.638 (0.515--0.760) \\
 &  &  & Grey SD & 23 / 146 & 55.856 (46.552--65.757) & 48.582 (42.591--55.653) & 0.070 / 0.089 / 0.250 & 0.661 (0.531--0.780) \\
\addlinespace[1pt]
DINOv2 & Kermany patient-disjoint & E1-M & Width (px) & 16 / 153 & 512.000 (512.000--512.000) & 512.000 (512.000--512.000) & 0.100 / -- / 0 & 0.471 (0.451--0.487) \\
 &  &  & Height (px) & 16 / 153 & 496.000 (496.000--496.000) & 496.000 (496.000--496.000) & 0.095 / -- / -- & 0.500 (0.500--0.500) \\
 &  &  & Aspect ratio & 16 / 153 & 1.032 (1.032--1.032) & 1.032 (1.032--1.032) & 0.100 / -- / 0 & 0.471 (0.450--0.487) \\
 &  &  & Border share & 16 / 153 & 0 (0--0.016) & 0 (0--0.016) & 0.110 / 0 / 0.089 & 0.453 (0.330--0.595) \\
 &  &  & Near-white share & 16 / 153 & 0.035 (0.019--0.044) & 0.022 (0.012--0.040) & 0.053 / 0.071 / 0.161 & 0.589 (0.450--0.729) \\
 &  &  & Grey mean & 16 / 153 & 44.239 (43.165--53.416) & 43.544 (38.452--48.872) & 0.053 / 0.125 / 0.107 & 0.595 (0.462--0.737) \\
 &  &  & Grey SD & 16 / 153 & 54.433 (46.749--58.490) & 48.850 (42.603--56.094) & 0.053 / 0.089 / 0.143 & 0.594 (0.463--0.733) \\
\addlinespace[1pt]
DINOv2 & OCTID NORMAL & locked & Width (px) & 18 / 188 & 750.000 (750.000--750.000) & 750.000 (750.000--750.000) & 0.084 / -- / 0.250 & 0.520 (0.487--0.581) \\
 &  &  & Height (px) & 18 / 188 & 500.000 (500.000--500.000) & 500.000 (500.000--500.000) & 0.084 / -- / 0.250 & 0.520 (0.484--0.581) \\
 &  &  & Aspect ratio & 18 / 188 & 1.500 (1.500--1.500) & 1.500 (1.500--1.500) & 0.087 / -- / -- & 0.500 (0.500--0.500) \\
 &  &  & Border share & 18 / 188 & 0 (0--0) & 0 (0--0) & 0.087 / -- / -- & 0.500 (0.500--0.500) \\
 &  &  & Near-white share & 18 / 188 & 0.004 (0.003--0.005) & 0.004 (0.003--0.005) & 0.072 / 0.130 / 0.059 & 0.466 (0.344--0.592) \\
 &  &  & Grey mean & 18 / 188 & 60.889 (59.608--62.138) & 60.824 (59.055--61.894) & 0.072 / 0.088 / 0.101 & 0.528 (0.384--0.668) \\
 &  &  & Grey SD & 18 / 188 & 42.811 (39.894--44.407) & 42.302 (39.892--43.834) & 0.087 / 0.059 / 0.116 & 0.554 (0.412--0.697) \\
\addlinespace[1pt]
DINOv2 & OCTID NORMAL & E1-M & Width (px) & 3 / 203 & 750.000 & 750.000 & -- / -- / -- & -- \\
 &  &  & Height (px) & 3 / 203 & 500.000 & 500.000 & -- / -- / -- & -- \\
 &  &  & Aspect ratio & 3 / 203 & 1.500 & 1.500 & -- / -- / -- & -- \\
 &  &  & Border share & 3 / 203 & 0 & 0 & -- / -- / -- & -- \\
 &  &  & Near-white share & 3 / 203 & 0.004 & 0.004 & -- / -- / -- & -- \\
 &  &  & Grey mean & 3 / 203 & 62.758 & 60.831 & -- / -- / -- & -- \\
 &  &  & Grey SD & 3 / 203 & 42.715 & 42.309 & -- / -- / -- & -- \\
\addlinespace[1pt]
DINOv2 & OCTDL NORMAL & locked & Width (px) & 187 / 145 & 882.000 (880.000--914.000) & 882.000 (879.000--1101.000) & 0.569 / 0.607 / 0.519 & 0.501 (0.421--0.585) \\
 &  &  & Height (px) & 187 / 145 & 321.000 (315.000--322.000) & 320.000 (314.000--387.000) & 0.554 / 0.593 / 0.525 & 0.519 (0.449--0.591) \\
 &  &  & Aspect ratio & 187 / 145 & 2.750 (2.739--2.813) & 2.759 (2.739--2.826) & 0.592 / 0.559 / 0.536 & 0.478 (0.406--0.549) \\
 &  &  & Border share & 187 / 145 & 0 (0--0) & 0 (0--0) & 0.563 / -- / -- & 0.500 (0.500--0.500) \\
 &  &  & Near-white share & 187 / 145 & $6.1\times10^{-4}$ ($4.4\times10^{-4}$--$7.4\times10^{-4}$) & $6.0\times10^{-4}$ ($3.9\times10^{-4}$--$7.5\times10^{-4}$) & 0.532 / 0.609 / 0.550 & 0.509 (0.431--0.591) \\
 &  &  & Grey mean & 187 / 145 & 56.468 (53.502--60.498) & 54.348 (51.243--60.225) & 0.423 / 0.664 / 0.604 & 0.586 (0.510--0.657) \\
 &  &  & Grey SD & 187 / 145 & 46.173 (44.542--47.969) & 45.418 (43.798--47.547) & 0.441 / 0.600 / 0.649 & 0.575 (0.505--0.642) \\
\addlinespace[1pt]
DINOv2 & OCTDL NORMAL & E1-M & Width (px) & 170 / 162 & 882.000 (880.000--1099.000) & 882.000 (879.000--912.750) & 0.507 / 0.488 / 0.538 & 0.534 (0.452--0.618) \\
 &  &  & Height (px) & 170 / 162 & 321.000 (315.000--388.000) & 320.000 (313.000--322.000) & 0.420 / 0.543 / 0.588 & 0.579 (0.508--0.652) \\
 &  &  & Aspect ratio & 170 / 162 & 2.747 (2.739--2.815) & 2.761 (2.739--2.820) & 0.558 / 0.461 / 0.509 & 0.458 (0.392--0.523) \\
 &  &  & Border share & 170 / 162 & 0 (0--0) & 0 (0--0) & 0.512 / -- / -- & 0.500 (0.500--0.500) \\
 &  &  & Near-white share & 170 / 162 & $6.1\times10^{-4}$ ($4.6\times10^{-4}$--$7.5\times10^{-4}$) & $6.0\times10^{-4}$ ($3.8\times10^{-4}$--$7.5\times10^{-4}$) & 0.486 / 0.536 / 0.514 & 0.537 (0.452--0.617) \\
 &  &  & Grey mean & 170 / 162 & 56.704 (53.578--60.593) & 54.594 (51.561--59.331) & 0.378 / 0.555 / 0.604 & 0.596 (0.525--0.667) \\
 &  &  & Grey SD & 170 / 162 & 46.090 (44.334--47.945) & 45.713 (44.041--47.833) & 0.477 / 0.482 / 0.577 & 0.537 (0.467--0.607) \\
\addlinespace[1pt]
MIRAGE & Kermany patient-disjoint & locked & Width (px) & 0 / 169 & -- & 512.000 & -- / -- / -- & -- \\
 &  &  & Height (px) & 0 / 169 & -- & 496.000 & -- / -- / -- & -- \\
 &  &  & Aspect ratio & 0 / 169 & -- & 1.032 & -- / -- / -- & -- \\
 &  &  & Border share & 0 / 169 & -- & 0 & -- / -- / -- & -- \\
 &  &  & Near-white share & 0 / 169 & -- & 0.023 & -- / -- / -- & -- \\
 &  &  & Grey mean & 0 / 169 & -- & 43.833 & -- / -- / -- & -- \\
 &  &  & Grey SD & 0 / 169 & -- & 49.120 & -- / -- / -- & -- \\
\addlinespace[1pt]
MIRAGE & Kermany patient-disjoint & E1-M & Width (px) & 0 / 169 & -- & 512.000 & -- / -- / -- & -- \\
 &  &  & Height (px) & 0 / 169 & -- & 496.000 & -- / -- / -- & -- \\
 &  &  & Aspect ratio & 0 / 169 & -- & 1.032 & -- / -- / -- & -- \\
 &  &  & Border share & 0 / 169 & -- & 0 & -- / -- / -- & -- \\
 &  &  & Near-white share & 0 / 169 & -- & 0.023 & -- / -- / -- & -- \\
 &  &  & Grey mean & 0 / 169 & -- & 43.833 & -- / -- / -- & -- \\
 &  &  & Grey SD & 0 / 169 & -- & 49.120 & -- / -- / -- & -- \\
\addlinespace[1pt]
MIRAGE & OCTID NORMAL & locked & Width (px) & 4 / 202 & 750.000 & 750.000 & -- / -- / -- & -- \\
 &  &  & Height (px) & 4 / 202 & 500.000 & 500.000 & -- / -- / -- & -- \\
 &  &  & Aspect ratio & 4 / 202 & 1.500 & 1.500 & -- / -- / -- & -- \\
 &  &  & Border share & 4 / 202 & 0 & 0 & -- / -- / -- & -- \\
 &  &  & Near-white share & 4 / 202 & 0.003 & 0.004 & -- / -- / -- & -- \\
 &  &  & Grey mean & 4 / 202 & 58.443 & 60.878 & -- / -- / -- & -- \\
 &  &  & Grey SD & 4 / 202 & 41.432 & 42.419 & -- / -- / -- & -- \\
\addlinespace[1pt]
MIRAGE & OCTID NORMAL & E1-M & Width (px) & 0 / 206 & -- & 750.000 & -- / -- / -- & -- \\
 &  &  & Height (px) & 0 / 206 & -- & 500.000 & -- / -- / -- & -- \\
 &  &  & Aspect ratio & 0 / 206 & -- & 1.500 & -- / -- / -- & -- \\
 &  &  & Border share & 0 / 206 & -- & 0 & -- / -- / -- & -- \\
 &  &  & Near-white share & 0 / 206 & -- & 0.004 & -- / -- / -- & -- \\
 &  &  & Grey mean & 0 / 206 & -- & 60.843 & -- / -- / -- & -- \\
 &  &  & Grey SD & 0 / 206 & -- & 42.361 & -- / -- / -- & -- \\
\addlinespace[1pt]
MIRAGE & OCTDL NORMAL & locked & Width (px) & 306 / 26 & 882.000 (880.000--914.000) & 882.000 (879.250--1099.000) & 0.917 / 0.929 / 0.923 & 0.515 (0.382--0.659) \\
 &  &  & Height (px) & 306 / 26 & 321.000 (314.000--322.000) & 321.000 (315.250--403.250) & 0.929 / 0.929 / 0.900 & 0.449 (0.327--0.581) \\
 &  &  & Aspect ratio & 306 / 26 & 2.755 (2.739--2.821) & 2.739 (2.687--2.779) & 0.883 / 0.912 / 0.973 & 0.642 (0.546--0.736) \\
 &  &  & Border share & 306 / 26 & 0 (0--0) & 0 (0--0) & 0.922 / -- / -- & 0.500 (0.500--0.500) \\
 &  &  & Near-white share & 306 / 26 & $6.0\times10^{-4}$ ($3.9\times10^{-4}$--$7.5\times10^{-4}$) & $6.7\times10^{-4}$ ($5.2\times10^{-4}$--$7.5\times10^{-4}$) & 0.946 / 0.918 / 0.901 & 0.416 (0.317--0.534) \\
 &  &  & Grey mean & 306 / 26 & 55.994 (52.648--60.445) & 52.680 (49.631--57.293) & 0.874 / 0.936 / 0.955 & 0.635 (0.510--0.752) \\
 &  &  & Grey SD & 306 / 26 & 46.002 (44.236--48.070) & 43.722 (41.018--45.456) & 0.847 / 0.945 / 0.973 & 0.726 (0.622--0.827) \\
\addlinespace[1pt]
MIRAGE & OCTDL NORMAL & E1-M & Width (px) & 267 / 65 & 882.000 (879.000--912.000) & 882.000 (881.000--1102.000) & 0.840 / 0.833 / 0.731 & 0.423 (0.326--0.518) \\
 &  &  & Height (px) & 267 / 65 & 321.000 (314.000--322.000) & 321.000 (314.000--389.000) & 0.759 / 0.914 / 0.675 & 0.463 (0.365--0.560) \\
 &  &  & Aspect ratio & 267 / 65 & 2.750 (2.739--2.815) & 2.787 (2.693--2.815) & 0.792 / 0.853 / 0.773 & 0.499 (0.402--0.597) \\
 &  &  & Border share & 267 / 65 & 0 (0--0) & 0 (0--0) & 0.804 / -- / -- & 0.500 (0.500--0.500) \\
 &  &  & Near-white share & 267 / 65 & $6.1\times10^{-4}$ ($3.9\times10^{-4}$--$7.5\times10^{-4}$) & $5.9\times10^{-4}$ ($4.6\times10^{-4}$--$7.3\times10^{-4}$) & 0.802 / 0.791 / 0.820 & 0.496 (0.402--0.595) \\
 &  &  & Grey mean & 267 / 65 & 56.596 (53.433--60.638) & 51.890 (48.918--55.831) & 0.613 / 0.891 / 0.910 & 0.725 (0.627--0.811) \\
 &  &  & Grey SD & 267 / 65 & 46.143 (44.479--48.203) & 43.812 (40.605--46.248) & 0.658 / 0.827 / 0.928 & 0.723 (0.641--0.794) \\
\addlinespace[1pt]
Inception-v3 & Kermany patient-disjoint & locked & Width (px) & 19 / 150 & 512.000 (512.000--768.000) & 512.000 (512.000--512.000) & 0.069 / -- / 0.889 & 0.707 (0.602--0.816) \\
 &  &  & Height (px) & 19 / 150 & 496.000 (496.000--496.000) & 496.000 (496.000--496.000) & 0.112 / -- / -- & 0.500 (0.500--0.500) \\
 &  &  & Aspect ratio & 19 / 150 & 1.032 (1.032--1.548) & 1.032 (1.032--1.032) & 0.069 / -- / 0.889 & 0.707 (0.599--0.816) \\
 &  &  & Border share & 19 / 150 & 0 (0--0.010) & 0 (0--0.016) & 0.120 / 0.154 / 0.089 & 0.478 (0.359--0.604) \\
 &  &  & Near-white share & 19 / 150 & 0.009 ($4.3\times10^{-4}$--0.021) & 0.024 (0.014--0.041) & 0.211 / 0.054 / 0.071 & 0.281 (0.139--0.452) \\
 &  &  & Grey mean & 19 / 150 & 36.841 (32.820--44.874) & 44.132 (39.262--49.365) & 0.211 / 0.054 / 0.071 & 0.327 (0.179--0.488) \\
 &  &  & Grey SD & 19 / 150 & 42.459 (37.774--48.379) & 49.608 (44.498--56.530) & 0.211 / 0.054 / 0.071 & 0.312 (0.170--0.471) \\
\addlinespace[1pt]
Inception-v3 & Kermany patient-disjoint & E1-M & Width (px) & 9 / 160 & 512.000 (512.000--768.000) & 512.000 (512.000--512.000) & 0.031 / -- / 0.444 & 0.707 (0.541--0.871) \\
 &  &  & Height (px) & 9 / 160 & 496.000 (496.000--496.000) & 496.000 (496.000--496.000) & 0.053 / -- / -- & 0.500 (0.500--0.500) \\
 &  &  & Aspect ratio & 9 / 160 & 1.032 (1.032--1.548) & 1.032 (1.032--1.032) & 0.031 / -- / 0.444 & 0.707 (0.544--0.874) \\
 &  &  & Border share & 9 / 160 & 0 (0--0) & 0 (0--0.016) & 0.070 / 0 / 0.036 & 0.405 (0.279--0.563) \\
 &  &  & Near-white share & 9 / 160 & 0.010 ($5.1\times10^{-4}$--0.025) & 0.023 (0.013--0.041) & 0.088 / 0.054 / 0.018 & 0.299 (0.112--0.506) \\
 &  &  & Grey mean & 9 / 160 & 36.841 (33.802--43.413) & 44.063 (39.116--49.513) & 0.105 / 0.036 / 0.018 & 0.293 (0.127--0.487) \\
 &  &  & Grey SD & 9 / 160 & 46.389 (39.878--46.958) & 49.402 (43.378--56.660) & 0.070 / 0.071 / 0.018 & 0.340 (0.176--0.514) \\
\addlinespace[1pt]
Inception-v3 & OCTID NORMAL & locked & Width (px) & 24 / 182 & 750.000 (750.000--750.000) & 750.000 (750.000--750.000) & 0.114 / -- / 0.250 & 0.513 (0.484--0.563) \\
 &  &  & Height (px) & 24 / 182 & 500.000 (500.000--500.000) & 500.000 (500.000--500.000) & 0.114 / -- / 0.250 & 0.513 (0.484--0.560) \\
 &  &  & Aspect ratio & 24 / 182 & 1.500 (1.500--1.500) & 1.500 (1.500--1.500) & 0.117 / -- / -- & 0.500 (0.500--0.500) \\
 &  &  & Border share & 24 / 182 & 0 (0--0) & 0 (0--0) & 0.117 / -- / -- & 0.500 (0.500--0.500) \\
 &  &  & Near-white share & 24 / 182 & 0.003 (0.002--0.005) & 0.004 (0.003--0.005) & 0.174 / 0.087 / 0.088 & 0.411 (0.282--0.552) \\
 &  &  & Grey mean & 24 / 182 & 59.939 (57.251--61.230) & 60.881 (59.215--62.040) & 0.159 / 0.118 / 0.072 & 0.361 (0.238--0.481) \\
 &  &  & Grey SD & 24 / 182 & 41.686 (37.908--43.891) & 42.419 (40.169--43.940) & 0.145 / 0.103 / 0.101 & 0.443 (0.306--0.582) \\
\addlinespace[1pt]
Inception-v3 & OCTID NORMAL & E1-M & Width (px) & 81 / 125 & 750.000 (750.000--750.000) & 750.000 (750.000--750.000) & 0.391 / -- / 0.500 & 0.504 (0.484--0.525) \\
 &  &  & Height (px) & 81 / 125 & 500.000 (500.000--500.000) & 500.000 (500.000--500.000) & 0.391 / -- / 0.500 & 0.504 (0.486--0.527) \\
 &  &  & Aspect ratio & 81 / 125 & 1.500 (1.500--1.500) & 1.500 (1.500--1.500) & 0.393 / -- / -- & 0.500 (0.500--0.500) \\
 &  &  & Border share & 81 / 125 & 0 (0--0) & 0 (0--0) & 0.393 / -- / -- & 0.500 (0.500--0.500) \\
 &  &  & Near-white share & 81 / 125 & 0.004 (0.003--0.005) & 0.004 (0.003--0.005) & 0.391 / 0.420 / 0.368 & 0.503 (0.421--0.584) \\
 &  &  & Grey mean & 81 / 125 & 60.687 (58.791--61.969) & 60.924 (59.396--61.871) & 0.464 / 0.294 / 0.420 & 0.472 (0.390--0.554) \\
 &  &  & Grey SD & 81 / 125 & 42.715 (39.076--44.255) & 42.296 (39.919--43.837) & 0.377 / 0.397 / 0.406 & 0.501 (0.417--0.584) \\
\addlinespace[1pt]
Inception-v3 & OCTDL NORMAL & locked & Width (px) & 266 / 66 & 882.000 (880.000--913.000) & 882.000 (879.250--1102.000) & 0.847 / 0.833 / 0.712 & 0.436 (0.347--0.530) \\
 &  &  & Height (px) & 266 / 66 & 321.000 (315.000--322.000) & 320.000 (312.000--389.000) & 0.768 / 0.871 / 0.725 & 0.516 (0.426--0.605) \\
 &  &  & Aspect ratio & 266 / 66 & 2.750 (2.739--2.809) & 2.786 (2.739--2.940) & 0.808 / 0.882 / 0.718 & 0.429 (0.339--0.527) \\
 &  &  & Border share & 266 / 66 & 0 (0--0) & 0 (0--0) & 0.801 / -- / -- & 0.500 (0.500--0.500) \\
 &  &  & Near-white share & 266 / 66 & $6.0\times10^{-4}$ ($4.0\times10^{-4}$--$7.5\times10^{-4}$) & $6.2\times10^{-4}$ ($4.3\times10^{-4}$--$7.5\times10^{-4}$) & 0.811 / 0.782 / 0.811 & 0.500 (0.410--0.587) \\
 &  &  & Grey mean & 266 / 66 & 56.149 (52.870--60.480) & 53.741 (51.105--58.337) & 0.712 / 0.855 / 0.838 & 0.584 (0.496--0.670) \\
 &  &  & Grey SD & 266 / 66 & 45.816 (44.174--47.898) & 45.838 (43.942--47.488) & 0.757 / 0.845 / 0.802 & 0.535 (0.450--0.621) \\
\addlinespace[1pt]
Inception-v3 & OCTDL NORMAL & E1-M & Width (px) & 217 / 115 & 882.000 (879.000--1097.000) & 882.000 (880.000--912.500) & 0.674 / 0.548 / 0.712 & 0.502 (0.412--0.583) \\
 &  &  & Height (px) & 217 / 115 & 320.000 (313.000--388.000) & 321.000 (315.500--322.000) & 0.679 / 0.607 / 0.700 & 0.482 (0.412--0.553) \\
 &  &  & Aspect ratio & 217 / 115 & 2.753 (2.739--2.817) & 2.759 (2.739--2.811) & 0.658 / 0.627 / 0.673 & 0.478 (0.405--0.550) \\
 &  &  & Border share & 217 / 115 & 0 (0--0) & 0 (0--0) & 0.654 / -- / -- & 0.500 (0.500--0.500) \\
 &  &  & Near-white share & 217 / 115 & $6.1\times10^{-4}$ ($4.5\times10^{-4}$--$7.5\times10^{-4}$) & $6.0\times10^{-4}$ ($3.8\times10^{-4}$--$7.5\times10^{-4}$) & 0.631 / 0.682 / 0.649 & 0.530 (0.453--0.617) \\
 &  &  & Grey mean & 217 / 115 & 56.370 (52.694--60.592) & 54.762 (51.694--58.712) & 0.613 / 0.618 / 0.730 & 0.578 (0.504--0.651) \\
 &  &  & Grey SD & 217 / 115 & 45.952 (44.131--47.829) & 45.635 (44.041--47.856) & 0.604 / 0.691 / 0.667 & 0.515 (0.443--0.581) \\
\end{longtable}
}
\fi
\havetable{e2d_categorical}\ifgen
\subsubsection*{E2d. OCTDL file mode and container format}
{\scriptsize% src: E2d_categorical.csv in results/fm_v1/posthoc (generated by make_tables.py; formatting only);
% label: post hoc, Phase-1 rule E2 timestamped by commit 4b4e883 (20 Nov 2026), computed 2026-11-21; outside the locked protocol; never replaces a locked result
\begin{longtable}{@{}lllllrrrrc@{}}
\toprule
Arm & Normals & Set & Property & Value & $n$ & Failed & Common & Rare & Failure share \\
\midrule
\endfirsthead
\toprule
Arm & Normals & Set & Property & Value & $n$ & Failed & Common & Rare & Failure share \\
\midrule
\endhead
\bottomrule
\endfoot
\bottomrule
\endlastfoot
RETFound & OCTDL NORMAL & locked & mode & RGB & 276 & 162 & 151 & 11 & 0.587 \\
RETFound & OCTDL NORMAL & locked & mode & RGBA & 56 & 33 & 26 & 7 & 0.589 \\
RETFound & OCTDL NORMAL & locked & format & JPEG & 270 & 157 & 147 & 10 & 0.581 \\
RETFound & OCTDL NORMAL & locked & format & PNG & 62 & 38 & 30 & 8 & 0.613 \\
RETFound & OCTDL NORMAL & E1-M & mode & RGB & 276 & 198 & 104 & 94 & 0.717 \\
RETFound & OCTDL NORMAL & E1-M & mode & RGBA & 56 & 42 & 5 & 37 & 0.750 \\
RETFound & OCTDL NORMAL & E1-M & format & JPEG & 270 & 193 & 102 & 91 & 0.715 \\
RETFound & OCTDL NORMAL & E1-M & format & PNG & 62 & 47 & 7 & 40 & 0.758 \\
DINOv2 & OCTDL NORMAL & locked & mode & RGB & 276 & 162 & 39 & 123 & 0.587 \\
DINOv2 & OCTDL NORMAL & locked & mode & RGBA & 56 & 25 & 7 & 18 & 0.446 \\
DINOv2 & OCTDL NORMAL & locked & format & JPEG & 270 & 159 & 38 & 121 & 0.589 \\
DINOv2 & OCTDL NORMAL & locked & format & PNG & 62 & 28 & 8 & 20 & 0.452 \\
DINOv2 & OCTDL NORMAL & E1-M & mode & RGB & 276 & 148 & 76 & 72 & 0.536 \\
DINOv2 & OCTDL NORMAL & E1-M & mode & RGBA & 56 & 22 & 14 & 8 & 0.393 \\
DINOv2 & OCTDL NORMAL & E1-M & format & JPEG & 270 & 146 & 75 & 71 & 0.541 \\
DINOv2 & OCTDL NORMAL & E1-M & format & PNG & 62 & 24 & 15 & 9 & 0.387 \\
MIRAGE & OCTDL NORMAL & locked & mode & RGB & 276 & 254 & 227 & 27 & 0.920 \\
MIRAGE & OCTDL NORMAL & locked & mode & RGBA & 56 & 52 & 49 & 3 & 0.929 \\
MIRAGE & OCTDL NORMAL & locked & format & JPEG & 270 & 248 & 225 & 23 & 0.919 \\
MIRAGE & OCTDL NORMAL & locked & format & PNG & 62 & 58 & 51 & 7 & 0.935 \\
MIRAGE & OCTDL NORMAL & E1-M & mode & RGB & 276 & 219 & 6 & 213 & 0.793 \\
MIRAGE & OCTDL NORMAL & E1-M & mode & RGBA & 56 & 48 & 0 & 48 & 0.857 \\
MIRAGE & OCTDL NORMAL & E1-M & format & JPEG & 270 & 214 & 5 & 209 & 0.793 \\
MIRAGE & OCTDL NORMAL & E1-M & format & PNG & 62 & 53 & 1 & 52 & 0.855 \\
Inception-v3 & OCTDL NORMAL & locked & mode & RGB & 276 & 220 & 211 & 9 & 0.797 \\
Inception-v3 & OCTDL NORMAL & locked & mode & RGBA & 56 & 46 & 44 & 2 & 0.821 \\
Inception-v3 & OCTDL NORMAL & locked & format & JPEG & 270 & 215 & 209 & 6 & 0.796 \\
Inception-v3 & OCTDL NORMAL & locked & format & PNG & 62 & 51 & 46 & 5 & 0.823 \\
Inception-v3 & OCTDL NORMAL & E1-M & mode & RGB & 276 & 181 & 103 & 78 & 0.656 \\
Inception-v3 & OCTDL NORMAL & E1-M & mode & RGBA & 56 & 36 & 20 & 16 & 0.643 \\
Inception-v3 & OCTDL NORMAL & E1-M & format & JPEG & 270 & 176 & 101 & 75 & 0.652 \\
Inception-v3 & OCTDL NORMAL & E1-M & format & PNG & 62 & 41 & 22 & 19 & 0.661 \\
\end{longtable}
}
\fi
\havetable{e2e_patients}\ifgen
\subsubsection*{E2e. OCTDL normal scans per patient}


\begin{center}\scriptsize% src: E2e_octdl_patients.csv in C:/Users/ThinkBook/.t3/worktrees/Few-shot/t3-03f3a2c4/results/fm_v1/posthoc (generated by make_tables.py; formatting only);
% label: post hoc, Phase-1 rule E2 timestamped by commit 4b4e883 (20 Nov 2026), computed 2026-11-21; outside the locked protocol; never replaces a locked result
\begin{tabular}{@{}llrcccccc@{}}
\toprule
Arm & Set & Patients & \makecell{All / some / no\\scans failed} & \makecell{Failure share\\per patient} & \makecell{Common-disease\\share per patient} & \makecell{False-rare share\\per patient} & \makecell{Scans per patient:\\mean, median (min--max)} \\
\midrule
MIRAGE & locked & 110 & 87 / 18 / 5 & 0.902 & 0.786 & 0.115 & 3.018, 2.000 (1--9) \\
MIRAGE & E1-M & 110 & 63 / 30 / 17 & 0.730 & 0.026 & 0.704 & 3.018, 2.000 (1--9) \\
MIRAGE & E1-L & 110 & 82 / 19 / 9 & 0.845 & 0.012 & 0.833 & 3.018, 2.000 (1--9) \\
MIRAGE & E1-C & 110 & 84 / 19 / 7 & 0.878 & 0.281 & 0.598 & 3.018, 2.000 (1--9) \\
DINOv2 & locked & 110 & 30 / 64 / 16 & 0.572 & 0.140 & 0.433 & 3.018, 2.000 (1--9) \\
DINOv2 & E1-M & 110 & 27 / 57 / 26 & 0.509 & 0.266 & 0.243 & 3.018, 2.000 (1--9) \\
DINOv2 & E1-L & 110 & 39 / 51 / 20 & 0.618 & 0.278 & 0.340 & 3.018, 2.000 (1--9) \\
DINOv2 & E1-C & 110 & 23 / 57 / 30 & 0.460 & 0.263 & 0.198 & 3.018, 2.000 (1--9) \\
Inception-v3 & locked & 110 & 57 / 41 / 12 & 0.762 & 0.719 & 0.043 & 3.018, 2.000 (1--9) \\
Inception-v3 & E1-M & 110 & 39 / 59 / 12 & 0.660 & 0.366 & 0.294 & 3.018, 2.000 (1--9) \\
Inception-v3 & E1-L & 110 & 34 / 59 / 17 & 0.586 & 0.366 & 0.219 & 3.018, 2.000 (1--9) \\
Inception-v3 & E1-C & 110 & 78 / 28 / 4 & 0.881 & 0.425 & 0.456 & 3.018, 2.000 (1--9) \\
RETFound & locked & 110 & 41 / 51 / 18 & 0.604 & 0.545 & 0.059 & 3.018, 2.000 (1--9) \\
RETFound & E1-M & 110 & 54 / 46 / 10 & 0.732 & 0.361 & 0.371 & 3.018, 2.000 (1--9) \\
RETFound & E1-L & 110 & 82 / 24 / 4 & 0.895 & 0.012 & 0.882 & 3.018, 2.000 (1--9) \\
RETFound & E1-C & 110 & 53 / 47 / 10 & 0.732 & 0.368 & 0.364 & 3.018, 2.000 (1--9) \\
\bottomrule
\end{tabular}
\end{center}
\fi
\havetable{e2f_pass1}\ifgen
\subsubsection*{E2f. Blinded contact sheet, pass 1 (E1-M squares)}
\input{tables/e2f_pass1_note.tex}
{\scriptsize\input{tables/e2f_pass1.tex}}
\fi
\havetable{e2f_pass2}\ifgen
\subsubsection*{E2f. Blinded contact sheet, pass 2 (originals)}
\input{tables/e2f_pass2_note.tex}
{\scriptsize\input{tables/e2f_pass2.tex}}
\fi

\havetable{e3_label}\ifgen
\subsection*{E3. Preprocessing sensitivity, MIRAGE and RETFound}
Label of every E3 row: post hoc, Phase-1 rule E3 timestamped by commit 4b4e883 (20 Nov 2026), computed 2026-11-21; outside the locked protocol; never replaces a locked result; input preprocessing only; pooling and token choice are not varied
.
\fi
\havetable{e3_criteria}\ifgen
\subsubsection*{E3. Criteria (1)--(3)}
{\scriptsize% src: E3_criteria.csv in results/fm_v1/posthoc (generated by make_tables.py; formatting only);
% label: post hoc, Phase-1 rule E3 timestamped by commit 4b4e883 (20 Nov 2026), computed 2026-11-21; outside the locked protocol; never replaces a locked result; input preprocessing only; pooling and token choice are not varied
\begin{longtable}{@{}p{1.6cm}lp{1.5cm}p{4.2cm}rcp{5.2cm}@{}}
\toprule
Criterion & Arm & Variant & Statistic & $n$ & Value & Result \\
\midrule
\endfirsthead
\toprule
Criterion & Arm & Variant & Statistic & $n$ & Value & Result \\
\midrule
\endhead
\bottomrule
\endfoot
\bottomrule
\endlastfoot
(1) & MIRAGE & M-crop & p05 of K mean (variant - Inception-v3 locked) & 100 & 0.071 & met \\
(1) & MIRAGE & M-bicubic & p05 of K mean (variant - Inception-v3 locked) & 100 & 0.075 & met \\
(1) & MIRAGE & M-nostretch & p05 of K mean (variant - Inception-v3 locked) & 100 & 0.067 & met \\
(1) & MIRAGE & M-res256 & p05 of K mean (variant - Inception-v3 locked) & 100 & 0.043 & met \\
(1) & MIRAGE & all four & result & -- & -- & kept within this grid \\
(1) shown beside, not counted & MIRAGE & M-pad & p05 of K mean (variant - Inception-v3 locked) & 100 & 0.060 & known before this rule (G1, 18 Nov) \\
(2) & RETFound & R-squash & p95 of K mean (variant - Inception-v3 locked) < 0 & 100 & $-$0.072 & met \\
(2) & RETFound & R-squash & p05 of K mean (MIRAGE locked - variant) > 0 & 100 & 0.185 & met \\
(2) & RETFound & R-pad & p95 of K mean (variant - Inception-v3 locked) < 0 & 100 & $-$0.091 & met \\
(2) & RETFound & R-pad & p05 of K mean (MIRAGE locked - variant) > 0 & 100 & 0.199 & met \\
(2) & RETFound & R-bilinear & p95 of K mean (variant - Inception-v3 locked) < 0 & 100 & $-$0.059 & met \\
(2) & RETFound & R-bilinear & p05 of K mean (MIRAGE locked - variant) > 0 & 100 & 0.184 & met \\
(2) & RETFound & R-stretch & p95 of K mean (variant - Inception-v3 locked) < 0 & 100 & $-$0.075 & met \\
(2) & RETFound & R-stretch & p05 of K mean (MIRAGE locked - variant) > 0 & 100 & 0.199 & met \\
(2) & RETFound & R-norm05 & p95 of K mean (variant - Inception-v3 locked) < 0 & 100 & 0.110 & not met \\
(2) & RETFound & R-norm05 & p05 of K mean (MIRAGE locked - variant) > 0 & 100 & 0.018 & met \\
(2) & RETFound & R-res128 & p95 of K mean (variant - Inception-v3 locked) < 0 & 100 & $-$0.066 & met \\
(2) & RETFound & R-res128 & p05 of K mean (MIRAGE locked - variant) > 0 & 100 & 0.191 & met \\
(2) & RETFound & all six & result & -- & -- & position changed or not shown for R-norm05 \\
(3) & MIRAGE & M-crop & OCTID NORMAL false-rare count inside MIRAGE's locked 95\% interval & 206 & 17 & outside [0.005, 0.049] as a rate; NORMAL recall 0.917; R1 CSC 41/102, MH 7/102; ranks among four arms RETFound 1, Inception 2, DINOv2 3, MIRAGE 3 (description only) \\
(3) & MIRAGE & M-crop & OCTDL NORMAL false-rare count inside MIRAGE's locked 95\% interval & 332 & 329 & outside [0.055, 0.136] as a rate; NORMAL recall 0.006; R1 MH 35/39; ranks among four arms Inception 1, RETFound 2, DINOv2 3, MIRAGE 4 (description only) \\
(3) & MIRAGE & M-bicubic & OCTID NORMAL false-rare count inside MIRAGE's locked 95\% interval & 206 & 1 & outside [0.005, 0.049] as a rate; NORMAL recall 0.995; R1 CSC 41/102, MH 8/102; ranks among four arms RETFound 1, MIRAGE 2, Inception 3, DINOv2 4 (description only) \\
(3) & MIRAGE & M-bicubic & OCTDL NORMAL false-rare count inside MIRAGE's locked 95\% interval & 332 & 26 & inside [0.055, 0.136] as a rate; NORMAL recall 0.090; R1 MH 3/39; ranks among four arms Inception 1, RETFound 2, MIRAGE 3, DINOv2 4 (description only) \\
(3) & MIRAGE & M-nostretch & OCTID NORMAL false-rare count inside MIRAGE's locked 95\% interval & 206 & 0 & outside [0.005, 0.049] as a rate; NORMAL recall 1.000; R1 CSC 41/102, MH 10/102; ranks among four arms RETFound 1, MIRAGE 1, Inception 3, DINOv2 4 (description only) \\
(3) & MIRAGE & M-nostretch & OCTDL NORMAL false-rare count inside MIRAGE's locked 95\% interval & 332 & 27 & inside [0.055, 0.136] as a rate; NORMAL recall 0.087; R1 MH 3/39; ranks among four arms Inception 1, RETFound 2, MIRAGE 3, DINOv2 4 (description only) \\
(3) & MIRAGE & M-res256 & OCTID NORMAL false-rare count inside MIRAGE's locked 95\% interval & 206 & 16 & outside [0.005, 0.049] as a rate; NORMAL recall 0.922; R1 CSC 55/102, MH 23/102; ranks among four arms RETFound 1, Inception 2, MIRAGE 3, DINOv2 4 (description only) \\
(3) & MIRAGE & M-res256 & OCTDL NORMAL false-rare count inside MIRAGE's locked 95\% interval & 332 & 26 & inside [0.055, 0.136] as a rate; NORMAL recall 0.084; R1 MH 3/39; ranks among four arms Inception 1, RETFound 2, MIRAGE 3, DINOv2 4 (description only) \\
(3) & RETFound & R-squash & OCTID NORMAL false-rare count inside RETFound's locked 95\% interval & 206 & 0 & inside [0.000, 0.018] as a rate; NORMAL recall 0.990; R1 CSC 1/102, MH 0/102; ranks among four arms RETFound 1, MIRAGE 2, Inception 3, DINOv2 4 (description only) \\
(3) & RETFound & R-squash & OCTDL NORMAL false-rare count inside RETFound's locked 95\% interval & 332 & 11 & inside [0.029, 0.087] as a rate; NORMAL recall 0.696; R1 MH 3/39; ranks among four arms Inception 1, RETFound 1, MIRAGE 3, DINOv2 4 (description only) \\
(3) & RETFound & R-pad & OCTID NORMAL false-rare count inside RETFound's locked 95\% interval & 206 & 118 & outside [0.000, 0.018] as a rate; NORMAL recall 0.262; R1 CSC 95/102, MH 81/102; ranks among four arms MIRAGE 1, Inception 2, DINOv2 3, RETFound 4 (description only) \\
(3) & RETFound & R-pad & OCTDL NORMAL false-rare count inside RETFound's locked 95\% interval & 332 & 332 & outside [0.029, 0.087] as a rate; NORMAL recall 0.000; R1 MH 39/39; ranks among four arms Inception 1, MIRAGE 2, DINOv2 3, RETFound 4 (description only) \\
(3) & RETFound & R-bilinear & OCTID NORMAL false-rare count inside RETFound's locked 95\% interval & 206 & 0 & inside [0.000, 0.018] as a rate; NORMAL recall 0.961; R1 CSC 0/102, MH 0/102; ranks among four arms RETFound 1, MIRAGE 2, Inception 3, DINOv2 4 (description only) \\
(3) & RETFound & R-bilinear & OCTDL NORMAL false-rare count inside RETFound's locked 95\% interval & 332 & 5 & outside [0.029, 0.087] as a rate; NORMAL recall 0.645; R1 MH 1/39; ranks among four arms RETFound 1, Inception 2, MIRAGE 3, DINOv2 4 (description only) \\
(3) & RETFound & R-stretch & OCTID NORMAL false-rare count inside RETFound's locked 95\% interval & 206 & 0 & inside [0.000, 0.018] as a rate; NORMAL recall 0.922; R1 CSC 0/102, MH 0/102; ranks among four arms RETFound 1, MIRAGE 2, Inception 3, DINOv2 4 (description only) \\
(3) & RETFound & R-stretch & OCTDL NORMAL false-rare count inside RETFound's locked 95\% interval & 332 & 23 & inside [0.029, 0.087] as a rate; NORMAL recall 0.386; R1 MH 4/39; ranks among four arms Inception 1, RETFound 2, MIRAGE 3, DINOv2 4 (description only) \\
(3) & RETFound & R-norm05 & OCTID NORMAL false-rare count inside RETFound's locked 95\% interval & 206 & 0 & inside [0.000, 0.018] as a rate; NORMAL recall 1.000; R1 CSC 2/102, MH 0/102; ranks among four arms RETFound 1, MIRAGE 2, Inception 3, DINOv2 4 (description only) \\
(3) & RETFound & R-norm05 & OCTDL NORMAL false-rare count inside RETFound's locked 95\% interval & 332 & 43 & outside [0.029, 0.087] as a rate; NORMAL recall 0.636; R1 MH 3/39; ranks among four arms Inception 1, MIRAGE 2, RETFound 3, DINOv2 4 (description only) \\
(3) & RETFound & R-res128 & OCTID NORMAL false-rare count inside RETFound's locked 95\% interval & 206 & 0 & inside [0.000, 0.018] as a rate; NORMAL recall 0.908; R1 CSC 0/102, MH 0/102; ranks among four arms RETFound 1, MIRAGE 2, Inception 3, DINOv2 4 (description only) \\
(3) & RETFound & R-res128 & OCTDL NORMAL false-rare count inside RETFound's locked 95\% interval & 332 & 2 & outside [0.029, 0.087] as a rate; NORMAL recall 0.934; R1 MH 0/39; ranks among four arms RETFound 1, Inception 2, MIRAGE 3, DINOv2 4 (description only) \\
(3) & MIRAGE & M-pad & OCTID NORMAL false-rare count inside MIRAGE's locked 95\% interval & 206 & 206 & outside [0.005, 0.049] as a rate; NORMAL recall 0.000; R1 CSC 98/102, MH 81/102; ranks among four arms RETFound 1, Inception 2, DINOv2 3, MIRAGE 4 (description only); known before this rule (G1, 18 Nov) \\
(3) & MIRAGE & M-pad & OCTDL NORMAL false-rare count inside MIRAGE's locked 95\% interval & 332 & 306 & outside [0.055, 0.136] as a rate; NORMAL recall 0.003; R1 MH 33/39; ranks among four arms Inception 1, RETFound 2, DINOv2 3, MIRAGE 4 (description only); known before this rule (G1, 18 Nov) \\
(3) & all &  & result & 11 & -- & outside: MIRAGE M-crop OCTID NORMAL; MIRAGE M-crop OCTDL NORMAL; MIRAGE M-bicubic OCTID NORMAL; MIRAGE M-nostretch OCTID NORMAL; MIRAGE M-res256 OCTID NORMAL; RETFound R-pad OCTID NORMAL; RETFound R-pad OCTDL NORMAL; RETFound R-bilinear OCTDL NORMAL; RETFound R-norm05 OCTDL NORMAL; RETFound R-res128 OCTDL NORMAL; MIRAGE M-pad OCTID NORMAL (known before this rule (G1, 18 Nov)); MIRAGE M-pad OCTDL NORMAL (known before this rule (G1, 18 Nov)) \\
scope &  &  &  & -- & -- & input preprocessing only; pooling and token choice are not varied \\
\end{longtable}
}
\fi
\havetable{e3_development}\ifgen
\subsubsection*{E3. Development macro recall per variant}
Per-$K$ rows are in the CSV.

{\scriptsize\setlength{\tabcolsep}{3pt}% src: E3_development.csv in results/fm_v1/posthoc (generated by make_tables.py; formatting only);
% label: post hoc, Phase-1 rule E3 timestamped by commit 4b4e883 (20 Nov 2026), computed 2026-11-21; outside the locked protocol; never replaces a locked result; input preprocessing only; pooling and token choice are not varied
\begin{longtable}{@{}llcccp{3.6cm}@{}}
\toprule
Arm & Variant & \makecell{Macro recall, mean $K$ 1--8\\(SD; 5th--95th)} & \makecell{$-$ Inception-v3 locked\\(paired)} & \makecell{$-$ same arm locked\\(paired)} & Status \\
\midrule
\endfirsthead
\toprule
Arm & Variant & \makecell{Macro recall, mean $K$ 1--8\\(SD; 5th--95th)} & \makecell{$-$ Inception-v3 locked\\(paired)} & \makecell{$-$ same arm locked\\(paired)} & Status \\
\midrule
\endhead
\bottomrule
\endfoot
\bottomrule
\endlastfoot
MIRAGE & locked & \makecell{0.650\\{\scriptsize (0.028; 0.603--0.688)}} & -- & -- & interpreted \\
 & M-crop & \makecell{0.658\\{\scriptsize (0.029; 0.611--0.696)}} & \makecell{+0.124\\{\scriptsize (+0.071 to +0.190)}} & \makecell{+0.008\\{\scriptsize ($-$0.011 to +0.027)}} & interpreted \\
 & M-bicubic & \makecell{0.653\\{\scriptsize (0.029; 0.604--0.693)}} & \makecell{+0.119\\{\scriptsize (+0.075 to +0.188)}} & \makecell{+0.004\\{\scriptsize ($-$0.008 to +0.013)}} & interpreted \\
 & M-nostretch & \makecell{0.652\\{\scriptsize (0.029; 0.605--0.694)}} & \makecell{+0.118\\{\scriptsize (+0.067 to +0.192)}} & \makecell{+0.003\\{\scriptsize ($-$0.008 to +0.015)}} & interpreted \\
 & M-res256 & \makecell{0.630\\{\scriptsize (0.027; 0.588--0.669)}} & \makecell{+0.096\\{\scriptsize (+0.043 to +0.165)}} & \makecell{$-$0.019\\{\scriptsize ($-$0.038 to +0.001)}} & interpreted \\
 & M-pad & \makecell{0.646\\{\scriptsize (0.028; 0.607--0.691)}} & \makecell{+0.112\\{\scriptsize (+0.060 to +0.175)}} & \makecell{$-$0.004\\{\scriptsize ($-$0.029 to +0.024)}} & known before this rule (G1, 18 Nov) \\
 & range over variants & 0.630--0.658 (5 variants) &  &  &  \\
\addlinespace[2pt]
RETFound & locked & \makecell{0.406\\{\scriptsize (0.024; 0.367--0.441)}} & -- & -- & interpreted \\
 & R-squash & \makecell{0.405\\{\scriptsize (0.023; 0.368--0.440)}} & \makecell{$-$0.129\\{\scriptsize ($-$0.179 to $-$0.072)}} & \makecell{$-$0.001\\{\scriptsize ($-$0.014 to +0.013)}} & interpreted \\
 & R-pad & \makecell{0.397\\{\scriptsize (0.020; 0.366--0.426)}} & \makecell{$-$0.137\\{\scriptsize ($-$0.187 to $-$0.091)}} & \makecell{$-$0.009\\{\scriptsize ($-$0.036 to +0.013)}} & interpreted \\
 & R-bilinear & \makecell{0.414\\{\scriptsize (0.024; 0.376--0.450)}} & \makecell{$-$0.120\\{\scriptsize ($-$0.172 to $-$0.059)}} & \makecell{+0.008\\{\scriptsize ($-$0.009 to +0.024)}} & interpreted \\
 & R-stretch & \makecell{0.396\\{\scriptsize (0.023; 0.361--0.432)}} & \makecell{$-$0.138\\{\scriptsize ($-$0.197 to $-$0.075)}} & \makecell{$-$0.010\\{\scriptsize ($-$0.023 to +0.002)}} & interpreted \\
 & R-norm05 & \makecell{0.578\\{\scriptsize (0.025; 0.536--0.617)}} & \makecell{+0.044\\{\scriptsize ($-$0.013 to +0.110)}} & \makecell{+0.172\\{\scriptsize (+0.135 to +0.216)}} & interpreted \\
 & R-res128 & \makecell{0.408\\{\scriptsize (0.024; 0.369--0.444)}} & \makecell{$-$0.126\\{\scriptsize ($-$0.182 to $-$0.066)}} & \makecell{+0.002\\{\scriptsize ($-$0.039 to +0.048)}} & interpreted \\
 & range over variants & 0.396--0.578 (6 variants) &  &  &  \\
\end{longtable}
}
\fi
\havetable{e3_holdout}\ifgen
\subsubsection*{E3. Hold-out direction per variant}
Mean paired difference from locked Inception-v3 over 100 draws; read for direction only.

{\scriptsize\setlength{\tabcolsep}{3pt}% src: E3_holdout.csv in C:/Users/ThinkBook/.t3/worktrees/Few-shot/t3-03f3a2c4/results/fm_v1/posthoc (generated by make_tables.py; formatting only);
% label: post hoc, Phase-1 rule E3 timestamped by commit 4b4e883 (20 Nov 2026), computed 2026-11-21; outside the locked protocol; never replaces a locked result; input preprocessing only; pooling and token choice are not varied
\begin{longtable}{@{}llcccccclp{3cm}@{}}
\toprule
Arm & Variant & $K=1$ & $K=2$ & $K=4$ & $K=8$ & \makecell{$K=16$\\(not read)} & Mean $K$ 1--8 & Direction & Status \\
\midrule
\endfirsthead
\toprule
Arm & Variant & $K=1$ & $K=2$ & $K=4$ & $K=8$ & \makecell{$K=16$\\(not read)} & Mean $K$ 1--8 & Direction & Status \\
\midrule
\endhead
\bottomrule
\endfoot
\bottomrule
\endlastfoot
MIRAGE & M-crop & +0.108 & +0.197 & +0.250 & +0.256 & +0.280 & +0.202 & positive & interpreted \\
 & M-bicubic & +0.116 & +0.194 & +0.249 & +0.252 & +0.278 & +0.203 & positive & interpreted \\
 & M-nostretch & +0.113 & +0.197 & +0.253 & +0.253 & +0.275 & +0.204 & positive & interpreted \\
 & M-res256 & +0.099 & +0.164 & +0.224 & +0.229 & +0.254 & +0.179 & positive & interpreted \\
 & M-pad & +0.087 & +0.165 & +0.232 & +0.247 & +0.309 & +0.183 & positive & known before this rule (G1, 18 Nov) \\
\addlinespace[2pt]
RETFound & R-squash & $-$0.181 & $-$0.194 & $-$0.194 & $-$0.201 & $-$0.149 & $-$0.192 & negative & interpreted \\
 & R-pad & $-$0.174 & $-$0.193 & $-$0.208 & $-$0.229 & $-$0.153 & $-$0.201 & negative & interpreted \\
 & R-bilinear & $-$0.163 & $-$0.180 & $-$0.180 & $-$0.188 & $-$0.144 & $-$0.178 & negative & interpreted \\
 & R-stretch & $-$0.183 & $-$0.194 & $-$0.201 & $-$0.207 & $-$0.165 & $-$0.196 & negative & interpreted \\
 & R-norm05 & $-$0.020 & +0.031 & +0.089 & +0.127 & +0.179 & +0.057 & positive & interpreted \\
 & R-res128 & $-$0.157 & $-$0.220 & $-$0.256 & $-$0.285 & $-$0.265 & $-$0.229 & negative & interpreted \\
\end{longtable}
}
\fi
\havetable{e3_nine_way}\ifgen
\subsubsection*{E3. Nine-way audit per variant}
Cells as in Table~\ref{tab:frr}. $^{*}$known before this rule (G1, 18 Nov).

{\scriptsize\setlength{\tabcolsep}{2.5pt}% src: E3_nine_way.csv in C:/Users/ThinkBook/.t3/worktrees/Few-shot/t3-03f3a2c4/results/fm_v1/posthoc (generated by make_tables.py; formatting only);
% label: post hoc, Phase-1 rule E3 timestamped by commit 4b4e883 (20 Nov 2026), computed 2026-11-21; outside the locked protocol; never replaces a locked result; input preprocessing only; pooling and token choice are not varied
\begin{longtable}{@{}llcccccc@{}}
\toprule
Arm & Variant & \makecell{Kermany normals\\(patient-disjoint)} & OCTID normals & OCTDL normals & \makecell{Called NORMAL\\OCTID / OCTDL} & \makecell{R1: OCTID CSC / MH;\\OCTDL MH} & \makecell{Hold-out\\correct} \\
\midrule
\endfirsthead
\toprule
Arm & Variant & \makecell{Kermany normals\\(patient-disjoint)} & OCTID normals & OCTDL normals & \makecell{Called NORMAL\\OCTID / OCTDL} & \makecell{R1: OCTID CSC / MH;\\OCTDL MH} & \makecell{Hold-out\\correct} \\
\midrule
\endhead
\bottomrule
\endfoot
\bottomrule
\endlastfoot
MIRAGE & locked & \makecell{0/169\\{\scriptsize (0.000; 0.000--0.000)}} & \makecell{4/206\\{\scriptsize (0.019; 0.005--0.049)}} & \makecell{30/332\\{\scriptsize (0.090; 0.055--0.136)}} & 202 / 26 & 41 / 11; 3 & 20 \\
 & M-crop & \makecell{0/169\\{\scriptsize (0.000; 0.000--0.000)}} & \makecell{17/206\\{\scriptsize (0.083; 0.049--0.129)}} & \makecell{329/332\\{\scriptsize (0.991; 0.978--1.000)}} & 189 / 2 & 41 / 7; 35 & 20 \\
 & M-bicubic & \makecell{0/169\\{\scriptsize (0.000; 0.000--0.000)}} & \makecell{1/206\\{\scriptsize (0.005; 0.000--0.027)}} & \makecell{26/332\\{\scriptsize (0.078; 0.046--0.123)}} & 205 / 30 & 41 / 8; 3 & 20 \\
 & M-nostretch & \makecell{0/169\\{\scriptsize (0.000; 0.000--0.000)}} & \makecell{0/206\\{\scriptsize (0.000; 0.000--0.018)}} & \makecell{27/332\\{\scriptsize (0.081; 0.049--0.125)}} & 206 / 29 & 41 / 10; 3 & 20 \\
 & M-res256 & \makecell{0/169\\{\scriptsize (0.000; 0.000--0.000)}} & \makecell{16/206\\{\scriptsize (0.078; 0.045--0.123)}} & \makecell{26/332\\{\scriptsize (0.078; 0.045--0.123)}} & 190 / 28 & 55 / 23; 3 & 20 \\
 & M-pad$^{*}$ & \makecell{2/169\\{\scriptsize (0.012; 0.000--0.030)}} & \makecell{206/206\\{\scriptsize (1.000; 0.982--1.000)}} & \makecell{306/332\\{\scriptsize (0.922; 0.889--0.953)}} & 0 / 1 & 98 / 81; 33 & 21 \\
\addlinespace[2pt]
RETFound & locked & \makecell{1/169\\{\scriptsize (0.006; 0.000--0.019)}} & \makecell{0/206\\{\scriptsize (0.000; 0.000--0.018)}} & \makecell{18/332\\{\scriptsize (0.054; 0.029--0.087)}} & 189 / 137 & 0 / 0; 3 & 9 \\
 & R-squash & \makecell{2/169\\{\scriptsize (0.012; 0.000--0.030)}} & \makecell{0/206\\{\scriptsize (0.000; 0.000--0.018)}} & \makecell{11/332\\{\scriptsize (0.033; 0.012--0.063)}} & 204 / 231 & 1 / 0; 3 & 9 \\
 & R-pad & \makecell{4/169\\{\scriptsize (0.024; 0.006--0.048)}} & \makecell{118/206\\{\scriptsize (0.573; 0.502--0.641)}} & \makecell{332/332\\{\scriptsize (1.000; 1.000--1.000)}} & 54 / 0 & 95 / 81; 39 & 8 \\
 & R-bilinear & \makecell{1/169\\{\scriptsize (0.006; 0.000--0.019)}} & \makecell{0/206\\{\scriptsize (0.000; 0.000--0.018)}} & \makecell{5/332\\{\scriptsize (0.015; 0.003--0.034)}} & 198 / 214 & 0 / 0; 1 & 8 \\
 & R-stretch & \makecell{1/169\\{\scriptsize (0.006; 0.000--0.019)}} & \makecell{0/206\\{\scriptsize (0.000; 0.000--0.018)}} & \makecell{23/332\\{\scriptsize (0.069; 0.041--0.101)}} & 190 / 128 & 0 / 0; 4 & 9 \\
 & R-norm05 & \makecell{4/169\\{\scriptsize (0.024; 0.006--0.048)}} & \makecell{0/206\\{\scriptsize (0.000; 0.000--0.018)}} & \makecell{43/332\\{\scriptsize (0.130; 0.091--0.176)}} & 206 / 211 & 2 / 0; 3 & 17 \\
 & R-res128 & \makecell{2/169\\{\scriptsize (0.012; 0.000--0.030)}} & \makecell{0/206\\{\scriptsize (0.000; 0.000--0.018)}} & \makecell{2/332\\{\scriptsize (0.006; 0.000--0.016)}} & 187 / 310 & 0 / 0; 0 & 6 \\
\end{longtable}
}
\fi
\havetable{e3_checks}\ifgen
\subsubsection*{E3. Checks}
Hashes are shortened to 12 characters; the full values are in the CSV.

{\scriptsize% src: E3_checks.csv in C:/Users/ThinkBook/.t3/worktrees/Few-shot/t3-03f3a2c4/results/fm_v1/posthoc (generated by make_tables.py; formatting only);
% label: post hoc, Phase-1 rule E3 timestamped by commit 4b4e883 (20 Nov 2026), computed 2026-11-21; outside the locked protocol; never replaces a locked result; input preprocessing only; pooling and token choice are not varied
\begin{longtable}{@{}llp{5.8cm}p{3.7cm}p{3.5cm}@{}}
\toprule
Arm & Variant & Check & Value & Status \\
\midrule
\endfirsthead
\toprule
Arm & Variant & Check & Value & Status \\
\midrule
\endhead
\bottomrule
\endfoot
\bottomrule
\endlastfoot
MIRAGE & E3-ref & E3-ref against locked features (5927 images) & $3.0\times10^{-7}$ & pass \\
MIRAGE & E3-ref & overrides & \{"fm\_extract.mirage\_input": "fm\_control\_extract.\_m\_reference"\} &  \\
MIRAGE & E3-ref & code\_matches\_lock & true &  \\
MIRAGE & E3-ref & environment\_matches\_lock & false &  \\
MIRAGE & E3-ref & repeat (recorded) & 0 & 5 images \\
MIRAGE & E3-ref & inputs changed (worker) & 0 &  \\
MIRAGE & E3-ref & constant min-max inputs (left unchanged) & 0 &  \\
MIRAGE & M-crop & overrides & \{"fm\_extract.mirage\_input": "fm\_control\_extract.\_m\_crop"\} &  \\
MIRAGE & M-crop & code\_matches\_lock & true &  \\
MIRAGE & M-crop & environment\_matches\_lock & false &  \\
MIRAGE & M-crop & repeat (recorded) & 0 & 5 images \\
MIRAGE & M-crop & inputs changed (worker) & 4895 &  \\
MIRAGE & M-crop & constant min-max inputs (left unchanged) & 0 &  \\
MIRAGE & M-crop & square main images unchanged (1032) & $2.0\times10^{-7}$ & pass \\
MIRAGE & M-crop & largest per-image difference on 4895 non-square images & 0.272 & change applied \\
MIRAGE & M-bicubic & overrides & \{"fm\_extract.mirage\_input": "fm\_control\_extract.\_m\_bicubic"\} &  \\
MIRAGE & M-bicubic & code\_matches\_lock & true &  \\
MIRAGE & M-bicubic & environment\_matches\_lock & false &  \\
MIRAGE & M-bicubic & repeat (recorded) & 0 & 5 images \\
MIRAGE & M-bicubic & inputs changed (worker) & 4909 &  \\
MIRAGE & M-bicubic & constant min-max inputs (left unchanged) & 0 &  \\
MIRAGE & M-bicubic & largest per-image difference on 4895 non-square images & 0.041 & change applied \\
MIRAGE & M-nostretch & overrides & \{"fm\_extract.mirage\_input": "fm\_control\_extract.\_m\_nostretch"\} &  \\
MIRAGE & M-nostretch & code\_matches\_lock & true &  \\
MIRAGE & M-nostretch & environment\_matches\_lock & false &  \\
MIRAGE & M-nostretch & repeat (recorded) & 0 & 5 images \\
MIRAGE & M-nostretch & inputs changed (worker) & 1614 &  \\
MIRAGE & M-nostretch & constant min-max inputs (left unchanged) & 0 &  \\
MIRAGE & M-nostretch & largest per-image difference on 4895 non-square images & 0.045 & change applied \\
MIRAGE & M-res256 & overrides & \{"fm\_extract.mirage\_input": "fm\_control\_extract.\_m\_res256"\} &  \\
MIRAGE & M-res256 & code\_matches\_lock & true &  \\
MIRAGE & M-res256 & environment\_matches\_lock & false &  \\
MIRAGE & M-res256 & repeat (recorded) & 0 & 5 images \\
MIRAGE & M-res256 & inputs changed (worker) & 5927 &  \\
MIRAGE & M-res256 & constant min-max inputs (left unchanged) & 0 &  \\
MIRAGE & M-res256 & largest per-image difference on 4895 non-square images & 0.133 & change applied \\
RETFound & E3-ref & E3-ref against locked features (5927 images) & $1.9\times10^{-7}$ & pass \\
RETFound & E3-ref & overrides & \{"fm\_extract.IMAGENET\_MEAN": [0.485, 0.456, 0.406], "fm\_extract.IMAGENET\_STD": [0.229, 0.224, 0.225], "prestep": "fm\_control\_extract.\_identity"\} &  \\
RETFound & E3-ref & code\_matches\_lock & true &  \\
RETFound & E3-ref & environment\_matches\_lock & false &  \\
RETFound & E3-ref & repeat (recorded) & 0 & 5 images \\
RETFound & E3-ref & inputs changed (worker) & 0 &  \\
RETFound & E3-ref & constant min-max inputs (left unchanged) & 0 &  \\
RETFound & R-squash & overrides & \{"fm\_extract.IMAGENET\_MEAN": [0.485, 0.456, 0.406], "fm\_extract.IMAGENET\_STD": [0.229, 0.224, 0.225]\} &  \\
RETFound & R-squash & code\_matches\_lock & true &  \\
RETFound & R-squash & environment\_matches\_lock & false &  \\
RETFound & R-squash & repeat (recorded) & 0 & 5 images \\
RETFound & R-squash & inputs changed (worker) & 5927 &  \\
RETFound & R-squash & constant min-max inputs (left unchanged) & 0 &  \\
RETFound & R-squash & square main images unchanged (1032) & $1.9\times10^{-7}$ & pass \\
RETFound & R-squash & largest per-image difference on 4895 non-square images & 0.108 & change applied \\
RETFound & R-pad & overrides & \{"fm\_extract.IMAGENET\_MEAN": [0.485, 0.456, 0.406], "fm\_extract.IMAGENET\_STD": [0.229, 0.224, 0.225]\} &  \\
RETFound & R-pad & code\_matches\_lock & true &  \\
RETFound & R-pad & environment\_matches\_lock & false &  \\
RETFound & R-pad & repeat (recorded) & 0 & 5 images \\
RETFound & R-pad & inputs changed (worker) & 4895 &  \\
RETFound & R-pad & constant min-max inputs (left unchanged) & 0 &  \\
RETFound & R-pad & square main images unchanged (1032) & $1.9\times10^{-7}$ & pass \\
RETFound & R-pad & largest per-image difference on 4895 non-square images & 0.222 & change applied \\
RETFound & R-bilinear & overrides & \{"fm\_extract.IMAGENET\_MEAN": [0.485, 0.456, 0.406], "fm\_extract.IMAGENET\_STD": [0.229, 0.224, 0.225]\} &  \\
RETFound & R-bilinear & code\_matches\_lock & true &  \\
RETFound & R-bilinear & environment\_matches\_lock & false &  \\
RETFound & R-bilinear & repeat (recorded) & 0 & 5 images \\
RETFound & R-bilinear & inputs changed (worker) & 5927 &  \\
RETFound & R-bilinear & constant min-max inputs (left unchanged) & 0 &  \\
RETFound & R-bilinear & largest per-image difference on 4895 non-square images & 0.063 & change applied \\
RETFound & R-stretch & overrides & \{"fm\_extract.IMAGENET\_MEAN": [0.485, 0.456, 0.406], "fm\_extract.IMAGENET\_STD": [0.229, 0.224, 0.225]\} &  \\
RETFound & R-stretch & code\_matches\_lock & true &  \\
RETFound & R-stretch & environment\_matches\_lock & false &  \\
RETFound & R-stretch & repeat (recorded) & 0 & 5 images \\
RETFound & R-stretch & inputs changed (worker) & 243 &  \\
RETFound & R-stretch & constant min-max inputs (left unchanged) & 0 &  \\
RETFound & R-stretch & largest per-image difference on 4895 non-square images & 0.048 & change applied \\
RETFound & R-norm05 & overrides & \{"fm\_extract.IMAGENET\_MEAN": [0.5, 0.5, 0.5], "fm\_extract.IMAGENET\_STD": [0.5, 0.5, 0.5]\} &  \\
RETFound & R-norm05 & code\_matches\_lock & true &  \\
RETFound & R-norm05 & environment\_matches\_lock & false &  \\
RETFound & R-norm05 & repeat (recorded) & 0 & 5 images \\
RETFound & R-norm05 & inputs changed (worker) & 0 &  \\
RETFound & R-norm05 & constant min-max inputs (left unchanged) & 0 &  \\
RETFound & R-norm05 & largest per-image difference on 4895 non-square images & 0.545 & change applied \\
RETFound & R-res128 & overrides & \{"fm\_extract.IMAGENET\_MEAN": [0.485, 0.456, 0.406], "fm\_extract.IMAGENET\_STD": [0.229, 0.224, 0.225]\} &  \\
RETFound & R-res128 & code\_matches\_lock & true &  \\
RETFound & R-res128 & environment\_matches\_lock & false &  \\
RETFound & R-res128 & repeat (recorded) & 0 & 5 images \\
RETFound & R-res128 & inputs changed (worker) & 5927 &  \\
RETFound & R-res128 & constant min-max inputs (left unchanged) & 0 &  \\
RETFound & R-res128 & largest per-image difference on 4895 non-square images & 0.271 & change applied \\
\end{longtable}
}
\fi

\havetable{e4_label}\ifgen
\subsection*{E4. Source separability}
Label of every E4 row: post hoc, Phase-1 rule E4 timestamped by commit 4b4e883 (20 Nov 2026), computed 2026-11-21; outside the locked protocol; never replaces a locked result; laptop run of record, scikit-learn 1.9.1
. E4 does not show whether the models use the source in the rare-disease task, and a low value does not mean that the features hold no source information.
\fi
\havetable{e4a_fewshot}\ifgen
\subsubsection*{E4a. Few-shot source test (main measure)}
Two-way macro recall of source against source (chance 0.50), 100 draws; paired rows: E1-M minus locked on the same draw. Sources: K, Kermany test (1,000; patient units); Y, the 146 rare images of Yoo et al.\ (image units); ID, OCTID (410; image units); DL, OCTDL (371; patient units); KN, IDN, DLN: their normal scans only. Pairs with Y: different sources, source and disease are mixed. 

{\tiny\let\scriptsize\tiny\setlength{\tabcolsep}{2pt}% src: E4a_fewshot_source_test.csv in results/fm_v1/posthoc (generated by make_tables.py; formatting only);
% label: post hoc, Phase-1 rule E4 timestamped by commit 4b4e883 (20 Nov 2026), computed 2026-11-21; outside the locked protocol; never replaces a locked result; laptop run of record, scikit-learn 1.9.1
\begin{longtable}{@{}llcccccp{3cm}@{}}
\toprule
Pair & Features & $K=1$ & $K=2$ & $K=4$ & $K=8$ & Mean $K$ 1--8 & Reading \\
\midrule
\endfirsthead
\toprule
Pair & Features & $K=1$ & $K=2$ & $K=4$ & $K=8$ & Mean $K$ 1--8 & Reading \\
\midrule
\endhead
\bottomrule
\endfoot
\bottomrule
\endlastfoot
KN--IDN & RETFound locked & \makecell{0.947\\{\scriptsize (0.079)}} & \makecell{0.981\\{\scriptsize (0.030)}} & \makecell{0.992\\{\scriptsize (0.016)}} & \makecell{0.997\\{\scriptsize (0.006)}} & \makecell{0.979\\{\scriptsize (0.026; 0.921--1.000)}} &  \\
 & RETFound E1-M & \makecell{0.951\\{\scriptsize (0.079)}} & \makecell{0.982\\{\scriptsize (0.030)}} & \makecell{0.993\\{\scriptsize (0.017)}} & \makecell{0.998\\{\scriptsize (0.005)}} & \makecell{0.981\\{\scriptsize (0.026; 0.924--1.000)}} &  \\
 & DINOv2 locked & \makecell{0.809\\{\scriptsize (0.129)}} & \makecell{0.917\\{\scriptsize (0.079)}} & \makecell{0.973\\{\scriptsize (0.030)}} & \makecell{0.989\\{\scriptsize (0.016)}} & \makecell{0.922\\{\scriptsize (0.048; 0.833--0.988)}} &  \\
 & DINOv2 E1-M & \makecell{0.820\\{\scriptsize (0.128)}} & \makecell{0.923\\{\scriptsize (0.071)}} & \makecell{0.972\\{\scriptsize (0.034)}} & \makecell{0.990\\{\scriptsize (0.015)}} & \makecell{0.926\\{\scriptsize (0.047; 0.833--0.988)}} &  \\
 & MIRAGE locked & \makecell{0.947\\{\scriptsize (0.089)}} & \makecell{0.996\\{\scriptsize (0.009)}} & \makecell{0.999\\{\scriptsize (0.003)}} & \makecell{1.000\\{\scriptsize (0.000)}} & \makecell{0.986\\{\scriptsize (0.024; 0.931--1.000)}} &  \\
 & MIRAGE E1-M & \makecell{0.989\\{\scriptsize (0.024)}} & \makecell{0.998\\{\scriptsize (0.003)}} & \makecell{0.999\\{\scriptsize (0.002)}} & \makecell{0.999\\{\scriptsize (0.001)}} & \makecell{0.996\\{\scriptsize (0.006; 0.985--1.000)}} &  \\
 & Inception-v3 locked & \makecell{0.787\\{\scriptsize (0.117)}} & \makecell{0.878\\{\scriptsize (0.087)}} & \makecell{0.939\\{\scriptsize (0.051)}} & \makecell{0.978\\{\scriptsize (0.017)}} & \makecell{0.895\\{\scriptsize (0.054; 0.797--0.960)}} &  \\
 & Inception-v3 E1-M & \makecell{0.761\\{\scriptsize (0.109)}} & \makecell{0.840\\{\scriptsize (0.081)}} & \makecell{0.900\\{\scriptsize (0.045)}} & \makecell{0.941\\{\scriptsize (0.031)}} & \makecell{0.860\\{\scriptsize (0.048; 0.779--0.933)}} &  \\
 & Thumbnail raw & \makecell{0.724\\{\scriptsize (0.116)}} & \makecell{0.794\\{\scriptsize (0.104)}} & \makecell{0.855\\{\scriptsize (0.087)}} & \makecell{0.906\\{\scriptsize (0.054)}} & \makecell{0.820\\{\scriptsize (0.063; 0.690--0.910)}} &  \\
 & Thumbnail E1-M & \makecell{0.745\\{\scriptsize (0.127)}} & \makecell{0.826\\{\scriptsize (0.098)}} & \makecell{0.889\\{\scriptsize (0.057)}} & \makecell{0.919\\{\scriptsize (0.042)}} & \makecell{0.845\\{\scriptsize (0.053; 0.747--0.918)}} &  \\
 & RETFound E1-M $-$ RETFound locked & \makecell{+0.004\\{\scriptsize ($-$0.008 to +0.021)}} & \makecell{+0.002\\{\scriptsize ($-$0.004 to +0.008)}} & \makecell{+0.001\\{\scriptsize ($-$0.002 to +0.004)}} & \makecell{+0.001\\{\scriptsize (0.000 to +0.004)}} & \makecell{+0.002\\{\scriptsize (0.004; $-$0.002 to +0.007)}} & no reduction shown \\
 & DINOv2 E1-M $-$ DINOv2 locked & \makecell{+0.011\\{\scriptsize ($-$0.042 to +0.132)}} & \makecell{+0.006\\{\scriptsize ($-$0.034 to +0.072)}} & \makecell{$-$0.001\\{\scriptsize ($-$0.022 to +0.021)}} & \makecell{+0.001\\{\scriptsize ($-$0.006 to +0.011)}} & \makecell{+0.004\\{\scriptsize (0.021; $-$0.014 to +0.035)}} & no reduction shown \\
 & MIRAGE E1-M $-$ MIRAGE locked & \makecell{+0.042\\{\scriptsize ($-$0.004 to +0.226)}} & \makecell{+0.002\\{\scriptsize ($-$0.005 to +0.015)}} & \makecell{$-$0.000\\{\scriptsize ($-$0.002 to +0.002)}} & \makecell{$-$0.001\\{\scriptsize ($-$0.003 to 0.000)}} & \makecell{+0.011\\{\scriptsize (0.022; $-$0.003 to +0.059)}} & no reduction shown \\
 & Inception-v3 E1-M $-$ Inception-v3 locked & \makecell{$-$0.026\\{\scriptsize ($-$0.210 to +0.128)}} & \makecell{$-$0.039\\{\scriptsize ($-$0.162 to +0.070)}} & \makecell{$-$0.038\\{\scriptsize ($-$0.132 to +0.050)}} & \makecell{$-$0.037\\{\scriptsize ($-$0.098 to $-$0.004)}} & \makecell{$-$0.035\\{\scriptsize (0.051; $-$0.116 to +0.040)}} & no reduction shown \\
 & Thumbnail E1-M $-$ Thumbnail raw & \makecell{+0.021\\{\scriptsize ($-$0.082 to +0.217)}} & \makecell{+0.032\\{\scriptsize ($-$0.042 to +0.227)}} & \makecell{+0.034\\{\scriptsize ($-$0.034 to +0.187)}} & \makecell{+0.013\\{\scriptsize ($-$0.026 to +0.087)}} & \makecell{+0.025\\{\scriptsize (0.051; $-$0.025 to +0.138)}} & no reduction shown \\
\addlinespace[2pt]
KN--DLN & RETFound locked & \makecell{0.885\\{\scriptsize (0.107)}} & \makecell{0.946\\{\scriptsize (0.056)}} & \makecell{0.975\\{\scriptsize (0.026)}} & \makecell{0.987\\{\scriptsize (0.013)}} & \makecell{0.948\\{\scriptsize (0.034; 0.884--0.985)}} &  \\
 & RETFound E1-M & \makecell{0.893\\{\scriptsize (0.105)}} & \makecell{0.950\\{\scriptsize (0.058)}} & \makecell{0.979\\{\scriptsize (0.027)}} & \makecell{0.992\\{\scriptsize (0.012)}} & \makecell{0.954\\{\scriptsize (0.034; 0.883--0.992)}} &  \\
 & DINOv2 locked & \makecell{0.674\\{\scriptsize (0.097)}} & \makecell{0.755\\{\scriptsize (0.087)}} & \makecell{0.810\\{\scriptsize (0.065)}} & \makecell{0.877\\{\scriptsize (0.051)}} & \makecell{0.779\\{\scriptsize (0.049; 0.708--0.852)}} &  \\
 & DINOv2 E1-M & \makecell{0.678\\{\scriptsize (0.095)}} & \makecell{0.767\\{\scriptsize (0.085)}} & \makecell{0.820\\{\scriptsize (0.067)}} & \makecell{0.881\\{\scriptsize (0.049)}} & \makecell{0.787\\{\scriptsize (0.046; 0.706--0.851)}} &  \\
 & MIRAGE locked & \makecell{0.967\\{\scriptsize (0.045)}} & \makecell{0.981\\{\scriptsize (0.011)}} & \makecell{0.988\\{\scriptsize (0.006)}} & \makecell{0.990\\{\scriptsize (0.004)}} & \makecell{0.981\\{\scriptsize (0.013; 0.957--0.992)}} &  \\
 & MIRAGE E1-M & \makecell{0.948\\{\scriptsize (0.084)}} & \makecell{0.975\\{\scriptsize (0.030)}} & \makecell{0.983\\{\scriptsize (0.013)}} & \makecell{0.987\\{\scriptsize (0.008)}} & \makecell{0.973\\{\scriptsize (0.022; 0.922--0.993)}} &  \\
 & Inception-v3 locked & \makecell{0.881\\{\scriptsize (0.085)}} & \makecell{0.931\\{\scriptsize (0.033)}} & \makecell{0.947\\{\scriptsize (0.019)}} & \makecell{0.952\\{\scriptsize (0.011)}} & \makecell{0.928\\{\scriptsize (0.027; 0.872--0.957)}} &  \\
 & Inception-v3 E1-M & \makecell{0.738\\{\scriptsize (0.093)}} & \makecell{0.790\\{\scriptsize (0.086)}} & \makecell{0.856\\{\scriptsize (0.060)}} & \makecell{0.904\\{\scriptsize (0.042)}} & \makecell{0.822\\{\scriptsize (0.048; 0.741--0.898)}} &  \\
 & Thumbnail raw & \makecell{0.803\\{\scriptsize (0.130)}} & \makecell{0.885\\{\scriptsize (0.090)}} & \makecell{0.932\\{\scriptsize (0.045)}} & \makecell{0.960\\{\scriptsize (0.019)}} & \makecell{0.895\\{\scriptsize (0.053; 0.780--0.954)}} &  \\
 & Thumbnail E1-M & \makecell{0.852\\{\scriptsize (0.126)}} & \makecell{0.914\\{\scriptsize (0.042)}} & \makecell{0.935\\{\scriptsize (0.023)}} & \makecell{0.940\\{\scriptsize (0.015)}} & \makecell{0.910\\{\scriptsize (0.035; 0.843--0.946)}} &  \\
 & RETFound E1-M $-$ RETFound locked & \makecell{+0.008\\{\scriptsize ($-$0.007 to +0.020)}} & \makecell{+0.005\\{\scriptsize ($-$0.010 to +0.018)}} & \makecell{+0.004\\{\scriptsize ($-$0.004 to +0.012)}} & \makecell{+0.005\\{\scriptsize ($-$0.002 to +0.011)}} & \makecell{+0.005\\{\scriptsize (0.006; $-$0.003 to +0.013)}} & no reduction shown \\
 & DINOv2 E1-M $-$ DINOv2 locked & \makecell{+0.004\\{\scriptsize ($-$0.035 to +0.047)}} & \makecell{+0.012\\{\scriptsize ($-$0.018 to +0.053)}} & \makecell{+0.011\\{\scriptsize ($-$0.021 to +0.043)}} & \makecell{+0.004\\{\scriptsize ($-$0.019 to +0.036)}} & \makecell{+0.008\\{\scriptsize (0.017; $-$0.015 to +0.034)}} & no reduction shown \\
 & MIRAGE E1-M $-$ MIRAGE locked & \makecell{$-$0.019\\{\scriptsize ($-$0.100 to +0.026)}} & \makecell{$-$0.006\\{\scriptsize ($-$0.027 to +0.024)}} & \makecell{$-$0.005\\{\scriptsize ($-$0.022 to +0.014)}} & \makecell{$-$0.003\\{\scriptsize ($-$0.020 to +0.013)}} & \makecell{$-$0.008\\{\scriptsize (0.018; $-$0.025 to +0.010)}} & no reduction shown \\
 & Inception-v3 E1-M $-$ Inception-v3 locked & \makecell{$-$0.143\\{\scriptsize ($-$0.319 to $-$0.003)}} & \makecell{$-$0.140\\{\scriptsize ($-$0.281 to $-$0.025)}} & \makecell{$-$0.091\\{\scriptsize ($-$0.192 to $-$0.007)}} & \makecell{$-$0.048\\{\scriptsize ($-$0.131 to $-$0.001)}} & \makecell{$-$0.106\\{\scriptsize (0.043; $-$0.186 to $-$0.035)}} & the control reduced separability \\
 & Thumbnail E1-M $-$ Thumbnail raw & \makecell{+0.049\\{\scriptsize ($-$0.074 to +0.275)}} & \makecell{+0.029\\{\scriptsize ($-$0.045 to +0.221)}} & \makecell{+0.004\\{\scriptsize ($-$0.038 to +0.091)}} & \makecell{$-$0.019\\{\scriptsize ($-$0.040 to +0.020)}} & \makecell{+0.016\\{\scriptsize (0.047; $-$0.030 to +0.107)}} & no reduction shown \\
\addlinespace[2pt]
IDN--DLN & RETFound locked & \makecell{0.922\\{\scriptsize (0.100)}} & \makecell{0.967\\{\scriptsize (0.025)}} & \makecell{0.982\\{\scriptsize (0.012)}} & \makecell{0.988\\{\scriptsize (0.008)}} & \makecell{0.965\\{\scriptsize (0.028; 0.911--0.988)}} &  \\
 & RETFound E1-M & \makecell{0.922\\{\scriptsize (0.099)}} & \makecell{0.967\\{\scriptsize (0.024)}} & \makecell{0.982\\{\scriptsize (0.012)}} & \makecell{0.988\\{\scriptsize (0.008)}} & \makecell{0.965\\{\scriptsize (0.028; 0.911--0.988)}} &  \\
 & DINOv2 locked & \makecell{0.845\\{\scriptsize (0.110)}} & \makecell{0.927\\{\scriptsize (0.077)}} & \makecell{0.973\\{\scriptsize (0.024)}} & \makecell{0.985\\{\scriptsize (0.008)}} & \makecell{0.932\\{\scriptsize (0.042; 0.855--0.983)}} &  \\
 & DINOv2 E1-M & \makecell{0.844\\{\scriptsize (0.107)}} & \makecell{0.927\\{\scriptsize (0.075)}} & \makecell{0.972\\{\scriptsize (0.022)}} & \makecell{0.985\\{\scriptsize (0.008)}} & \makecell{0.932\\{\scriptsize (0.041; 0.860--0.982)}} &  \\
 & MIRAGE locked & \makecell{0.982\\{\scriptsize (0.021)}} & \makecell{0.991\\{\scriptsize (0.011)}} & \makecell{0.995\\{\scriptsize (0.003)}} & \makecell{0.996\\{\scriptsize (0.003)}} & \makecell{0.991\\{\scriptsize (0.007; 0.982--0.997)}} &  \\
 & MIRAGE E1-M & \makecell{0.946\\{\scriptsize (0.075)}} & \makecell{0.975\\{\scriptsize (0.029)}} & \makecell{0.985\\{\scriptsize (0.009)}} & \makecell{0.988\\{\scriptsize (0.008)}} & \makecell{0.974\\{\scriptsize (0.022; 0.943--0.990)}} &  \\
 & Inception-v3 locked & \makecell{0.958\\{\scriptsize (0.035)}} & \makecell{0.977\\{\scriptsize (0.011)}} & \makecell{0.985\\{\scriptsize (0.008)}} & \makecell{0.988\\{\scriptsize (0.006)}} & \makecell{0.977\\{\scriptsize (0.010; 0.961--0.991)}} &  \\
 & Inception-v3 E1-M & \makecell{0.864\\{\scriptsize (0.091)}} & \makecell{0.937\\{\scriptsize (0.043)}} & \makecell{0.958\\{\scriptsize (0.027)}} & \makecell{0.972\\{\scriptsize (0.013)}} & \makecell{0.933\\{\scriptsize (0.031; 0.884--0.970)}} &  \\
 & Thumbnail raw & \makecell{0.836\\{\scriptsize (0.129)}} & \makecell{0.900\\{\scriptsize (0.052)}} & \makecell{0.922\\{\scriptsize (0.021)}} & \makecell{0.931\\{\scriptsize (0.014)}} & \makecell{0.897\\{\scriptsize (0.040; 0.834--0.930)}} &  \\
 & Thumbnail E1-M & \makecell{0.890\\{\scriptsize (0.113)}} & \makecell{0.924\\{\scriptsize (0.040)}} & \makecell{0.937\\{\scriptsize (0.016)}} & \makecell{0.942\\{\scriptsize (0.009)}} & \makecell{0.923\\{\scriptsize (0.035; 0.868--0.944)}} &  \\
 & RETFound E1-M $-$ RETFound locked & \makecell{$-$0.000\\{\scriptsize ($-$0.006 to +0.006)}} & \makecell{$-$0.000\\{\scriptsize ($-$0.003 to +0.003)}} & \makecell{$-$0.000\\{\scriptsize ($-$0.002 to +0.002)}} & \makecell{$-$0.000\\{\scriptsize ($-$0.002 to 0.000)}} & \makecell{$-$0.000\\{\scriptsize (0.001; $-$0.002 to +0.002)}} & no reduction shown \\
 & DINOv2 E1-M $-$ DINOv2 locked & \makecell{$-$0.001\\{\scriptsize ($-$0.046 to +0.029)}} & \makecell{0.000\\{\scriptsize ($-$0.015 to +0.018)}} & \makecell{$-$0.000\\{\scriptsize ($-$0.008 to +0.006)}} & \makecell{$-$0.000\\{\scriptsize ($-$0.003 to +0.003)}} & \makecell{$-$0.000\\{\scriptsize (0.008; $-$0.015 to +0.013)}} & no reduction shown \\
 & MIRAGE E1-M $-$ MIRAGE locked & \makecell{$-$0.036\\{\scriptsize ($-$0.167 to $-$0.001)}} & \makecell{$-$0.016\\{\scriptsize ($-$0.032 to $-$0.001)}} & \makecell{$-$0.010\\{\scriptsize ($-$0.022 to +0.002)}} & \makecell{$-$0.007\\{\scriptsize ($-$0.020 to +0.002)}} & \makecell{$-$0.017\\{\scriptsize (0.017; $-$0.044 to $-$0.005)}} & the control reduced separability \\
 & Inception-v3 E1-M $-$ Inception-v3 locked & \makecell{$-$0.094\\{\scriptsize ($-$0.275 to +0.004)}} & \makecell{$-$0.040\\{\scriptsize ($-$0.120 to +0.011)}} & \makecell{$-$0.027\\{\scriptsize ($-$0.086 to +0.003)}} & \makecell{$-$0.017\\{\scriptsize ($-$0.043 to 0.000)}} & \makecell{$-$0.044\\{\scriptsize (0.030; $-$0.094 to $-$0.007)}} & the control reduced separability \\
 & Thumbnail E1-M $-$ Thumbnail raw & \makecell{+0.054\\{\scriptsize ($-$0.006 to +0.262)}} & \makecell{+0.024\\{\scriptsize ($-$0.006 to +0.088)}} & \makecell{+0.015\\{\scriptsize ($-$0.005 to +0.036)}} & \makecell{+0.011\\{\scriptsize ($-$0.001 to +0.027)}} & \makecell{+0.026\\{\scriptsize (0.022; +0.006 to +0.074)}} & no reduction shown \\
\end{longtable}
}
\fi
\havetable{e4b_columns}\ifgen
\subsubsection*{E4b. Per-column separability, no fit}
Separability of one column with no fit, max(AUROC, 1 $-$ AUROC); reading: $\geq 0.95$ near-perfectly, 0.80--0.95 strongly, 0.60--0.80 partly, below 0.60 hardly separable. Sources: K, Kermany test (1,000; patient units); Y, the 146 rare images of Yoo et al.\ (image units); ID, OCTID (410; image units); DL, OCTDL (371; patient units); KN, IDN, DLN: their normal scans only. Pairs with Y: different sources, source and disease are mixed.

\begin{center}\scriptsize% src: E4b_column_auroc.csv in results/fm_v1/posthoc (generated by make_tables.py; formatting only);
% label: post hoc, Phase-1 rule E4 timestamped by commit 4b4e883 (20 Nov 2026), computed 2026-11-21; outside the locked protocol; never replaces a locked result; laptop run of record, scikit-learn 1.9.1
\begin{tabular}{@{}lccccccccc@{}}
\toprule
Column & K--Y & K--ID & K--DL & Y--ID & Y--DL & ID--DL & KN--IDN & KN--DLN & IDN--DLN \\
\midrule
G: width & 0.963 & 0.790 & 0.987 & 1.000 & 1.000 & 0.937 & 0.941 & 0.997 & 0.955 \\
G: height & 0.966 & 1.000 & 0.997 & 0.969 & 0.786 & 0.998 & 1.000 & 0.997 & 0.997 \\
G: aspect\_ratio & 0.710 & 0.769 & 0.999 & 0.993 & 1.000 & 1.000 & 0.940 & 1.000 & 1.000 \\
G: border\_share & 0.975 & 0.672 & 0.694 & 0.985 & 0.986 & 0.530 & 0.696 & 0.696 & 0.500 \\
S-raw: mean & 0.835 & 0.748 & 0.761 & 0.946 & 0.958 & 0.548 & 0.942 & 0.863 & 0.752 \\
S-raw: sd & 0.625 & 0.848 & 0.615 & 0.791 & 0.525 & 0.902 & 0.796 & 0.638 & 0.818 \\
S-raw: p01 & 0.506 & 0.994 & 0.977 & 0.969 & 0.952 & 0.842 & 1.000 & 0.982 & 0.946 \\
S-raw: p50 & 0.822 & 0.867 & 0.722 & 0.955 & 0.888 & 0.681 & 0.986 & 0.833 & 0.764 \\
S-raw: p99 & 0.814 & 0.797 & 0.753 & 0.611 & 0.733 & 0.519 & 0.851 & 0.885 & 0.884 \\
S-raw: share\_le8 & 0.954 & 1.000 & 0.963 & 0.973 & 0.972 & 0.744 & 1.000 & 0.971 & 0.896 \\
S-raw: share\_ge250 & 0.944 & 0.780 & 0.833 & 0.937 & 0.837 & 0.976 & 0.917 & 0.965 & 1.000 \\
S-raw: entropy\_bits & 0.989 & 0.836 & 0.789 & 0.985 & 0.993 & 0.967 & 0.704 & 0.845 & 0.955 \\
S-E1: mean & 0.646 & 0.832 & 0.875 & 0.622 & 0.668 & 0.530 & 0.984 & 0.948 & 0.661 \\
S-E1: sd & 0.580 & 0.838 & 0.628 & 0.851 & 0.507 & 0.962 & 0.782 & 0.662 & 0.940 \\
S-E1: p01 & 0.568 & 0.997 & 0.983 & 0.841 & 0.767 & 0.895 & 1.000 & 0.992 & 0.960 \\
S-E1: p50 & 0.505 & 0.893 & 0.783 & 0.745 & 0.681 & 0.606 & 0.992 & 0.857 & 0.683 \\
S-E1: p99 & 0.731 & 0.753 & 0.642 & 0.531 & 0.765 & 0.746 & 0.746 & 0.730 & 0.547 \\
S-E1: share\_le8 & 0.628 & 1.000 & 0.964 & 0.841 & 0.770 & 0.725 & 1.000 & 0.974 & 0.925 \\
S-E1: share\_ge250 & 0.891 & 0.872 & 0.778 & 0.560 & 0.816 & 0.771 & 0.911 & 0.901 & 0.596 \\
S-E1: entropy\_bits & 0.589 & 0.988 & 0.615 & 0.828 & 0.652 & 0.978 & 0.967 & 0.741 & 0.966 \\
\bottomrule
\end{tabular}
\end{center}
\fi
\havetable{e4c_probe}\ifgen
\subsubsection*{E4c. Logistic probe}
Logistic probe, out-of-fold AUROC pooled over five grouped folds (95\% interval from 2,000 resamples of units within each source, without refitting). $^{*}$Some fits did not converge (count in the CSV). $^{\dagger}$Status other than interpreted (in the CSV). Sources: K, Kermany test (1,000; patient units); Y, the 146 rare images of Yoo et al.\ (image units); ID, OCTID (410; image units); DL, OCTDL (371; patient units); KN, IDN, DLN: their normal scans only. Pairs with Y: different sources, source and disease are mixed.

\begin{center}\scriptsize\setlength{\tabcolsep}{3pt}\resizebox{\textwidth}{!}{% src: E4c_probe_auroc.csv in results/fm_v1/posthoc (generated by make_tables.py; formatting only);
% label: post hoc, Phase-1 rule E4 timestamped by commit 4b4e883 (20 Nov 2026), computed 2026-11-21; outside the locked protocol; never replaces a locked result; laptop run of record, scikit-learn 1.9.1
\begin{tabular}{@{}lccccccccc@{}}
\toprule
Features & K--Y & K--ID & K--DL & Y--ID & Y--DL & ID--DL & KN--IDN & KN--DLN & IDN--DLN \\
\midrule
G & \makecell{0.997\\{\scriptsize (0.992--1.000)}} & \makecell{0.999\\{\scriptsize (0.997--1.000)}} & \makecell{0.999\\{\scriptsize (0.997--1.000)}} & \makecell{1.000\\{\scriptsize (1.000--1.000)}} & \makecell{1.000\\{\scriptsize (1.000--1.000)}} & \makecell{1.000\\{\scriptsize (1.000--1.000)}} & \makecell{1.000\\{\scriptsize (1.000--1.000)}} & \makecell{1.000\\{\scriptsize (1.000--1.000)}} & \makecell{1.000\\{\scriptsize (1.000--1.000)}} \\
S-raw & \makecell{0.996\\{\scriptsize (0.987--1.000)}} & \makecell{1.000\\{\scriptsize (1.000--1.000)}} & \makecell{0.994\\{\scriptsize (0.987--1.000)}} & \makecell{0.996\\{\scriptsize (0.992--0.999)}} & \makecell{0.989\\{\scriptsize (0.975--0.999)}} & \makecell{1.000\\{\scriptsize (0.999--1.000)}} & \makecell{1.000\\{\scriptsize (1.000--1.000)}} & \makecell{0.999\\{\scriptsize (0.996--1.000)}} & \makecell{1.000\\{\scriptsize (1.000--1.000)}} \\
S-E1 & \makecell{0.991\\{\scriptsize (0.984--0.996)}} & \makecell{1.000\\{\scriptsize (1.000--1.000)}} & \makecell{1.000\\{\scriptsize (0.999--1.000)}} & \makecell{0.954\\{\scriptsize (0.923--0.980)}} & \makecell{0.957\\{\scriptsize (0.931--0.978)}} & \makecell{0.993\\{\scriptsize (0.988--0.998)}} & \makecell{1.000\\{\scriptsize (1.000--1.000)}} & \makecell{0.998\\{\scriptsize (0.994--1.000)}} & \makecell{0.994\\{\scriptsize (0.988--0.998)}} \\
RETFound locked & \makecell{0.991\\{\scriptsize (0.978--0.998)}} & \makecell{1.000\\{\scriptsize (1.000--1.000)}} & \makecell{0.999\\{\scriptsize (0.998--1.000)}} & \makecell{0.998\\{\scriptsize (0.994--1.000)}} & \makecell{0.993\\{\scriptsize (0.982--0.999)}} & \makecell{0.998\\{\scriptsize (0.997--1.000)}} & \makecell{1.000\\{\scriptsize (1.000--1.000)}} & \makecell{1.000\\{\scriptsize (0.999--1.000)}} & \makecell{1.000\\{\scriptsize (1.000--1.000)}} \\
RETFound E1-M & \makecell{0.996\\{\scriptsize (0.993--0.998)}} & \makecell{1.000\\{\scriptsize (1.000--1.000)}} & \makecell{1.000\\{\scriptsize (0.999--1.000)}} & \makecell{1.000\\{\scriptsize (1.000--1.000)}} & \makecell{0.991\\{\scriptsize (0.984--0.997)}} & \makecell{0.999\\{\scriptsize (0.997--1.000)}} & \makecell{1.000\\{\scriptsize (1.000--1.000)}} & \makecell{1.000\\{\scriptsize (1.000--1.000)}} & \makecell{1.000\\{\scriptsize (1.000--1.000)}} \\
DINOv2 locked & \makecell{0.987\\{\scriptsize (0.979--0.993)}} & \makecell{0.996\\{\scriptsize (0.993--0.998)}} & \makecell{0.994\\{\scriptsize (0.991--0.996)}} & \makecell{0.975\\{\scriptsize (0.962--0.985)}} & \makecell{0.960\\{\scriptsize (0.940--0.976)}} & \makecell{0.994\\{\scriptsize (0.991--0.997)}} & \makecell{1.000\\{\scriptsize (1.000--1.000)}} & \makecell{0.998\\{\scriptsize (0.995--0.999)}} & \makecell{1.000\\{\scriptsize (1.000--1.000)}} \\
DINOv2 E1-M & \makecell{0.993\\{\scriptsize (0.990--0.997)}} & \makecell{0.996\\{\scriptsize (0.993--0.998)}} & \makecell{0.994\\{\scriptsize (0.991--0.996)}} & \makecell{0.977\\{\scriptsize (0.965--0.986)}} & \makecell{0.960\\{\scriptsize (0.940--0.976)}} & \makecell{0.994\\{\scriptsize (0.990--0.997)}} & \makecell{1.000\\{\scriptsize (1.000--1.000)}} & \makecell{0.997\\{\scriptsize (0.994--0.999)}} & \makecell{1.000\\{\scriptsize (1.000--1.000)}} \\
MIRAGE locked & \makecell{1.000\\{\scriptsize (0.999--1.000)}} & \makecell{1.000\\{\scriptsize (1.000--1.000)}} & \makecell{0.998\\{\scriptsize (0.997--1.000)}} & \makecell{1.000\\{\scriptsize (0.999--1.000)}} & \makecell{0.999\\{\scriptsize (0.997--1.000)}} & \makecell{1.000\\{\scriptsize (0.999--1.000)}} & \makecell{1.000\\{\scriptsize (1.000--1.000)}} & \makecell{1.000\\{\scriptsize (0.999--1.000)}} & \makecell{1.000\\{\scriptsize (1.000--1.000)}} \\
MIRAGE E1-M & \makecell{0.997\\{\scriptsize (0.994--0.999)}} & \makecell{1.000\\{\scriptsize (0.999--1.000)}} & \makecell{0.999\\{\scriptsize (0.999--1.000)}} & \makecell{0.995\\{\scriptsize (0.991--0.998)}} & \makecell{0.987\\{\scriptsize (0.978--0.993)}} & \makecell{0.999\\{\scriptsize (0.997--1.000)}} & \makecell{1.000\\{\scriptsize (1.000--1.000)}} & \makecell{0.999\\{\scriptsize (0.999--1.000)}} & \makecell{1.000\\{\scriptsize (1.000--1.000)}} \\
Inception-v3 locked & \makecell{0.997\\{\scriptsize (0.995--0.999)}} & \makecell{0.996\\{\scriptsize (0.993--0.998)}} & \makecell{0.986\\{\scriptsize (0.978--0.993)}} & \makecell{0.998\\{\scriptsize (0.995--0.999)}} & \makecell{0.994\\{\scriptsize (0.988--0.999)}} & \makecell{0.997\\{\scriptsize (0.993--0.999)}} & \makecell{1.000\\{\scriptsize (1.000--1.000)}} & \makecell{0.996\\{\scriptsize (0.993--0.999)}} & \makecell{1.000\\{\scriptsize (1.000--1.000)}} \\
Inception-v3 E1-M & \makecell{0.971\\{\scriptsize (0.957--0.982)}} & \makecell{0.985\\{\scriptsize (0.979--0.990)}} & \makecell{0.979\\{\scriptsize (0.969--0.987)}} & \makecell{0.979\\{\scriptsize (0.969--0.988)}} & \makecell{0.959\\{\scriptsize (0.940--0.974)}} & \makecell{0.993\\{\scriptsize (0.987--0.997)}} & \makecell{1.000\\{\scriptsize (1.000--1.000)}} & \makecell{0.995\\{\scriptsize (0.991--0.998)}} & \makecell{0.999\\{\scriptsize (0.998--1.000)}} \\
\bottomrule
\end{tabular}
}\end{center}
\fi
\havetable{e4d_probe_diff}\ifgen
\subsubsection*{E4d. Logistic probe, paired differences}
Paired AUROC differences on the same resamples. Reading: $^{r}$the control reduced separability (interval wholly below 0); $^{n}$no reduction shown; $^{c}$at ceiling, the difference is not read; $^{d}$descriptive comparison. Sources: K, Kermany test (1,000; patient units); Y, the 146 rare images of Yoo et al.\ (image units); ID, OCTID (410; image units); DL, OCTDL (371; patient units); KN, IDN, DLN: their normal scans only. Pairs with Y: different sources, source and disease are mixed.

\begin{center}\scriptsize\setlength{\tabcolsep}{3pt}\resizebox{\textwidth}{!}{% src: E4d_probe_differences.csv in results/fm_v1/posthoc (generated by make_tables.py; formatting only);
% label: post hoc, Phase-1 rule E4 timestamped by commit 4b4e883 (20 Nov 2026), computed 2026-11-21; outside the locked protocol; never replaces a locked result; laptop run of record, scikit-learn 1.9.1
\begin{tabular}{@{}lccccccccc@{}}
\toprule
Comparison & K--Y & K--ID & K--DL & Y--ID & Y--DL & ID--DL & KN--IDN & KN--DLN & IDN--DLN \\
\midrule
S-E1 $-$ S-raw & \makecell{$-$0.005$^{c}$\\{\scriptsize ($-$0.014 to +0.006)}} & \makecell{0.000$^{c}$\\{\scriptsize (0.000 to 0.000)}} & \makecell{+0.006$^{c}$\\{\scriptsize (0.000 to +0.013)}} & \makecell{$-$0.042$^{r}$\\{\scriptsize ($-$0.073 to $-$0.016)}} & \makecell{$-$0.032$^{r}$\\{\scriptsize ($-$0.060 to $-$0.006)}} & \makecell{$-$0.006$^{c}$\\{\scriptsize ($-$0.012 to $-$0.002)}} & \makecell{0.000$^{c}$\\{\scriptsize (0.000 to 0.000)}} & \makecell{$-$0.000$^{c}$\\{\scriptsize ($-$0.002 to +0.001)}} & \makecell{$-$0.006$^{c}$\\{\scriptsize ($-$0.012 to $-$0.002)}} \\
RETFound E1-M $-$ RETFound locked & \makecell{+0.005$^{c}$\\{\scriptsize ($-$0.002 to +0.016)}} & \makecell{$-$0.000$^{c}$\\{\scriptsize ($-$0.000 to 0.000)}} & \makecell{0.000$^{c}$\\{\scriptsize (0.000 to +0.001)}} & \makecell{+0.002$^{c}$\\{\scriptsize ($-$0.000 to +0.006)}} & \makecell{$-$0.002$^{c}$\\{\scriptsize ($-$0.011 to +0.009)}} & \makecell{0.000$^{c}$\\{\scriptsize ($-$0.000 to 0.000)}} & \makecell{0.000$^{c}$\\{\scriptsize (0.000 to 0.000)}} & \makecell{0.000$^{c}$\\{\scriptsize (0.000 to +0.001)}} & \makecell{0.000$^{c}$\\{\scriptsize (0.000 to 0.000)}} \\
DINOv2 E1-M $-$ DINOv2 locked & \makecell{+0.007$^{n}$\\{\scriptsize (+0.002 to +0.012)}} & \makecell{$-$0.000$^{c}$\\{\scriptsize ($-$0.001 to 0.000)}} & \makecell{0.000$^{c}$\\{\scriptsize ($-$0.001 to +0.001)}} & \makecell{+0.002$^{n}$\\{\scriptsize ($-$0.004 to +0.009)}} & \makecell{$-$0.000$^{n}$\\{\scriptsize ($-$0.009 to +0.008)}} & \makecell{$-$0.000$^{c}$\\{\scriptsize ($-$0.001 to 0.000)}} & \makecell{0.000$^{c}$\\{\scriptsize (0.000 to 0.000)}} & \makecell{$-$0.001$^{c}$\\{\scriptsize ($-$0.002 to +0.001)}} & \makecell{0.000$^{c}$\\{\scriptsize ($-$0.000 to 0.000)}} \\
MIRAGE E1-M $-$ MIRAGE locked & \makecell{$-$0.003$^{c}$\\{\scriptsize ($-$0.005 to $-$0.001)}} & \makecell{$-$0.000$^{c}$\\{\scriptsize ($-$0.000 to $-$0.000)}} & \makecell{+0.001$^{c}$\\{\scriptsize ($-$0.000 to +0.003)}} & \makecell{$-$0.005$^{c}$\\{\scriptsize ($-$0.009 to $-$0.001)}} & \makecell{$-$0.012$^{r}$\\{\scriptsize ($-$0.021 to $-$0.006)}} & \makecell{$-$0.001$^{c}$\\{\scriptsize ($-$0.003 to $-$0.000)}} & \makecell{0.000$^{c}$\\{\scriptsize (0.000 to 0.000)}} & \makecell{$-$0.000$^{c}$\\{\scriptsize ($-$0.001 to 0.000)}} & \makecell{$-$0.000$^{c}$\\{\scriptsize ($-$0.000 to 0.000)}} \\
Inception-v3 E1-M $-$ Inception-v3 locked & \makecell{$-$0.027$^{r}$\\{\scriptsize ($-$0.040 to $-$0.016)}} & \makecell{$-$0.011$^{r}$\\{\scriptsize ($-$0.016 to $-$0.007)}} & \makecell{$-$0.007$^{n}$\\{\scriptsize ($-$0.017 to +0.002)}} & \makecell{$-$0.018$^{r}$\\{\scriptsize ($-$0.029 to $-$0.010)}} & \makecell{$-$0.035$^{r}$\\{\scriptsize ($-$0.054 to $-$0.020)}} & \makecell{$-$0.003$^{c}$\\{\scriptsize ($-$0.009 to 0.000)}} & \makecell{$-$0.000$^{c}$\\{\scriptsize ($-$0.000 to 0.000)}} & \makecell{$-$0.001$^{c}$\\{\scriptsize ($-$0.005 to +0.003)}} & \makecell{$-$0.001$^{c}$\\{\scriptsize ($-$0.002 to $-$0.000)}} \\
RETFound locked $-$ S-raw & \makecell{$-$0.005$^{c}$\\{\scriptsize ($-$0.019 to +0.006)}} & \makecell{$-$0.000$^{c}$\\{\scriptsize ($-$0.000 to 0.000)}} & \makecell{+0.005$^{c}$\\{\scriptsize ($-$0.000 to +0.012)}} & \makecell{+0.002$^{c}$\\{\scriptsize ($-$0.003 to +0.007)}} & \makecell{+0.005$^{d}$\\{\scriptsize ($-$0.010 to +0.020)}} & \makecell{$-$0.001$^{c}$\\{\scriptsize ($-$0.003 to 0.000)}} & \makecell{0.000$^{c}$\\{\scriptsize (0.000 to 0.000)}} & \makecell{+0.001$^{c}$\\{\scriptsize ($-$0.000 to +0.004)}} & \makecell{0.000$^{c}$\\{\scriptsize (0.000 to 0.000)}} \\
DINOv2 locked $-$ S-raw & \makecell{$-$0.009$^{d}$\\{\scriptsize ($-$0.018 to +0.002)}} & \makecell{$-$0.004$^{c}$\\{\scriptsize ($-$0.007 to $-$0.002)}} & \makecell{$-$0.000$^{c}$\\{\scriptsize ($-$0.006 to +0.007)}} & \makecell{$-$0.021$^{d}$\\{\scriptsize ($-$0.034 to $-$0.010)}} & \makecell{$-$0.029$^{d}$\\{\scriptsize ($-$0.051 to $-$0.008)}} & \makecell{$-$0.005$^{c}$\\{\scriptsize ($-$0.009 to $-$0.002)}} & \makecell{0.000$^{c}$\\{\scriptsize (0.000 to 0.000)}} & \makecell{$-$0.001$^{c}$\\{\scriptsize ($-$0.004 to +0.002)}} & \makecell{$-$0.000$^{c}$\\{\scriptsize ($-$0.000 to 0.000)}} \\
MIRAGE locked $-$ S-raw & \makecell{+0.004$^{c}$\\{\scriptsize ($-$0.000 to +0.013)}} & \makecell{$-$0.000$^{c}$\\{\scriptsize ($-$0.000 to 0.000)}} & \makecell{+0.004$^{c}$\\{\scriptsize ($-$0.001 to +0.012)}} & \makecell{+0.003$^{c}$\\{\scriptsize ($-$0.000 to +0.008)}} & \makecell{+0.010$^{d}$\\{\scriptsize ($-$0.000 to +0.025)}} & \makecell{$-$0.000$^{c}$\\{\scriptsize ($-$0.000 to 0.000)}} & \makecell{0.000$^{c}$\\{\scriptsize (0.000 to 0.000)}} & \makecell{+0.001$^{c}$\\{\scriptsize ($-$0.000 to +0.004)}} & \makecell{0.000$^{c}$\\{\scriptsize (0.000 to 0.000)}} \\
Inception-v3 locked $-$ S-raw & \makecell{+0.002$^{c}$\\{\scriptsize ($-$0.004 to +0.011)}} & \makecell{$-$0.004$^{c}$\\{\scriptsize ($-$0.007 to $-$0.002)}} & \makecell{$-$0.008$^{d}$\\{\scriptsize ($-$0.018 to +0.002)}} & \makecell{+0.002$^{c}$\\{\scriptsize ($-$0.003 to +0.006)}} & \makecell{+0.006$^{d}$\\{\scriptsize ($-$0.006 to +0.020)}} & \makecell{$-$0.003$^{c}$\\{\scriptsize ($-$0.006 to $-$0.001)}} & \makecell{0.000$^{c}$\\{\scriptsize (0.000 to 0.000)}} & \makecell{$-$0.002$^{c}$\\{\scriptsize ($-$0.006 to +0.001)}} & \makecell{0.000$^{c}$\\{\scriptsize (0.000 to 0.000)}} \\
\bottomrule
\end{tabular}
}\end{center}
\fi
% src: tables/e*.tex, generated by make_tables.py from the E CSVs of fm_posthoc.py (e1-e4, e2f-marks); each file's first
%      comment names its CSV and label. Rules: PR@fec827e l.587-590.
